% concatenated from main.tex, math_commands.tex
\documentclass[11pt,letterpaper]{article}
\usepackage{geometry}
\geometry{left=1.25in, top=1.25in, right=1.25in, bottom=1.25in}

% Recommended, but optional, packages for figures and better typesetting:
\usepackage[utf8]{inputenc}
\usepackage{comment}
\usepackage{microtype}
\usepackage{graphicx}
\usepackage{subfigure}
\usepackage{booktabs} % for professional tables
\usepackage{scalerel}
% \usepackage[margin=1in]{geometry}

% change style to "authoryear-comp" if would like to enable in-text citation citet
\usepackage[style=alphabetic, maxbibnames=15, giveninits=true, maxcitenames=15, natbib=true, maxalphanames=10, backend=biber, sorting=nty, backref=true, uniquename=false]{biblatex}
\addbibresource{refs.bib}
\DefineBibliographyStrings{english}{backrefpage = {page},backrefpages = {pages},}
\DeclareNameAlias{sortname}{given-family}
% Remove "in:" before journals
\renewbibmacro{in:}{%
\ifentrytype{article}{}{\printtext{\bibstring{in}\intitlepunct}}}
  

\usepackage{amssymb}
\usepackage{amsthm}
\usepackage{amsmath}
\usepackage{mathtools}

\usepackage{nicefrac}

\usepackage{algorithmic}
\usepackage{algorithm}
\usepackage[shortlabels]{enumitem}

\usepackage{url}
\usepackage{float}

\setlength{\parskip}{0.5em}
\setlength{\parindent}{0em}

\usepackage{pifont}
\usepackage{hhline}
\usepackage{multirow}

\newcommand{\blue}[1]{\textcolor{blue}{#1}}

\newcommand{\myalert}[1]{{\vspace{2mm}\noindent\textbf{#1}}}
\newcommand{\xmark}{\ding{55}}
\newcommand{\tickmark}{\ding{51}}
\newcommand{\reals}{\mathbb{R}}

\newcommand{\red}[1]{{\color{red}{#1}}}
\newcommand{\rj}[1]{{\color{blue}{#1}}}
\newcommand{\calT}{\mathcal{T}}
%\newcommand{\Tr}{\mathbf{Tr}}
\newcommand{\Det}{\mathbf{Det}}

% Adapted for use with thuthesis.
% Original code is at https://github.com/goodfeli/dlbook_notation/blob/master/math_commands.tex

%%%%% NEW MATH DEFINITIONS %%%%%

\newcommand\ceil[1]{\lceil #1 \rceil}
\newcommand\floor[1]{\lfloor #1 \rfloor}


% Vectors
% \newcommand\Vector[1]{\symbf{#1}}
\newcommand\Vector[1]{\mathbf{#1}}

\newcommand\0{{\Vector{0}}}
\newcommand\vzero{{\Vector{0}}}
\newcommand\1{{\Vector{1}}}
\newcommand\vone{{\Vector{1}}}

\newcommand\va{{\Vector{a}}}
\newcommand\vb{{\Vector{b}}}
\newcommand\vc{{\Vector{c}}}
\newcommand\vd{{\Vector{d}}}
\newcommand\ve{{\Vector{e}}}
\newcommand\vf{{\Vector{f}}}
\newcommand\vg{{\Vector{g}}}
\newcommand\vh{{\Vector{h}}}
\newcommand\vi{{\Vector{i}}}
\newcommand\vj{{\Vector{j}}}
\newcommand\vk{{\Vector{k}}}
\newcommand\vl{{\Vector{l}}}
\newcommand\vm{{\Vector{m}}}
\newcommand\vn{{\Vector{n}}}
\newcommand\vo{{\Vector{o}}}
\newcommand\vp{{\Vector{p}}}
\newcommand\vq{{\Vector{q}}}
\newcommand\vr{{\Vector{r}}}
\newcommand\vs{{\Vector{s}}}
\newcommand\vt{{\Vector{t}}}
\newcommand\vu{{\Vector{u}}}
\newcommand\vv{{\Vector{v}}}
\newcommand\vw{{\Vector{w}}}
\newcommand\vx{{\Vector{x}}}
\newcommand\vy{{\Vector{y}}}
\newcommand\vz{{\Vector{z}}}

\newcommand\valpha{{\Vector{\alpha}}}
\newcommand\vbeta{{\Vector{\beta}}}
\newcommand\vgamma{{\Vector{\gamma}}}
\newcommand\vdelta{{\Vector{\delta}}}
\newcommand\vepsilon{{\Vector{\epsilon}}}
\newcommand\vtheta{{\Vector{\theta}}}
\newcommand\viota{{\Vector{\iota}}}
\newcommand\vkappa{{\Vector{\kappa}}}
\newcommand\vlambda{{\Vector{\lambda}}}
\newcommand\vmu{{\Vector{\mu}}}
\newcommand\vnu{{\Vector{\nu}}}
\newcommand\vxi{{\Vector{\xi}}}
\newcommand\vpi{{\Vector{\pi}}}
\newcommand\vrho{{\Vector{\rho}}}
\newcommand\vsigma{{\Vector{\sigma}}}
\newcommand\vtau{{\Vector{\tau}}}
\newcommand\vupsilon{{\Vector{\upsilon}}}
\newcommand\vphi{{\Vector{\phi}}}
\newcommand\vchi{{\Vector{\chi}}}
\newcommand\vpsi{{\Vector{\psi}}}
\newcommand\vomega{{\Vector{\omega}}}
\newcommand\veta{{\Vector{\eta}}}

\newcommand\vha{{\hat{\Vector{a}}}}
\newcommand\vhb{{\hat{\Vector{b}}}}
\newcommand\vhc{{\hat{\Vector{c}}}}
\newcommand\vhd{{\hat{\Vector{d}}}}
\newcommand\vhe{{\hat{\Vector{e}}}}
\newcommand\vhf{{\hat{\Vector{f}}}}
\newcommand\vhg{{\hat{\Vector{g}}}}
\newcommand\vhh{{\hat{\Vector{h}}}}
\newcommand\vhi{{\hat{\Vector{i}}}}
\newcommand\vhj{{\hat{\Vector{j}}}}
\newcommand\vhk{{\hat{\Vector{k}}}}
\newcommand\vhl{{\hat{\Vector{l}}}}
\newcommand\vhm{{\hat{\Vector{m}}}}
\newcommand\vhn{{\hat{\Vector{n}}}}
\newcommand\vho{{\hat{\Vector{o}}}}
\newcommand\vhp{{\hat{\Vector{p}}}}
\newcommand\vhq{{\hat{\Vector{q}}}}
\newcommand\vhr{{\hat{\Vector{r}}}}
\newcommand\vhs{{\hat{\Vector{s}}}}
\newcommand\vht{{\hat{\Vector{t}}}}
\newcommand\vhu{{\hat{\Vector{u}}}}
\newcommand\vhv{{\hat{\Vector{v}}}}
\newcommand\vhw{{\hat{\Vector{w}}}}
\newcommand\vhx{{\hat{\Vector{x}}}}
\newcommand\vhy{{\hat{\Vector{y}}}}
\newcommand\vhz{{\hat{\Vector{z}}}}
% Matrix
% \newcommand\MATRIX[1]{\symbf{#1}}
\newcommand\MATRIX[1]{\mathbf{#1}}

\newcommand\mA{{\MATRIX{A}}}
\newcommand\mB{{\MATRIX{B}}}
\newcommand\mC{{\MATRIX{C}}}
\newcommand\mD{{\MATRIX{D}}}
\newcommand\mE{{\MATRIX{E}}}
\newcommand\mF{{\MATRIX{F}}}
\newcommand\mG{{\MATRIX{G}}}
\newcommand\mH{{\MATRIX{H}}}
\newcommand\mI{{\MATRIX{I}}}
\newcommand\mJ{{\MATRIX{J}}}
\newcommand\mK{{\MATRIX{K}}}
\newcommand\mL{{\MATRIX{L}}}
\newcommand\mM{{\MATRIX{M}}}
\newcommand\mN{{\MATRIX{N}}}
\newcommand\mO{{\MATRIX{O}}}
\newcommand\mP{{\MATRIX{P}}}
\newcommand\mQ{{\MATRIX{Q}}}
\newcommand\mR{{\MATRIX{R}}}
\newcommand\mS{{\MATRIX{S}}}
\newcommand\mT{{\MATRIX{T}}}
\newcommand\mU{{\MATRIX{U}}}
\newcommand\mV{{\MATRIX{V}}}
\newcommand\mW{{\MATRIX{W}}}
\newcommand\mX{{\MATRIX{X}}}
\newcommand\mY{{\MATRIX{Y}}}
\newcommand\mZ{{\MATRIX{Z}}}

\newcommand\mhA{{\hat{\MATRIX{A}}}}
\newcommand\mhB{{\hat{\MATRIX{B}}}}
\newcommand\mhC{{\hat{\MATRIX{C}}}}
\newcommand\mhD{{\hat{\MATRIX{D}}}}
\newcommand\mhE{{\hat{\MATRIX{E}}}}
\newcommand\mhF{{\hat{\MATRIX{F}}}}
\newcommand\mhG{{\hat{\MATRIX{G}}}}
\newcommand\mhH{{\hat{\MATRIX{H}}}}
\newcommand\mhI{{\hat{\MATRIX{I}}}}
\newcommand\mhJ{{\hat{\MATRIX{J}}}}
\newcommand\mhK{{\hat{\MATRIX{K}}}}
\newcommand\mhL{{\hat{\MATRIX{L}}}}
\newcommand\mhM{{\hat{\MATRIX{M}}}}
\newcommand\mhN{{\hat{\MATRIX{N}}}}
\newcommand\mhO{{\hat{\MATRIX{O}}}}
\newcommand\mhP{{\hat{\MATRIX{P}}}}
\newcommand\mhQ{{\hat{\MATRIX{Q}}}}
\newcommand\mhR{{\hat{\MATRIX{R}}}}
\newcommand\mhS{{\hat{\MATRIX{S}}}}
\newcommand\mhT{{\hat{\MATRIX{T}}}}
\newcommand\mhU{{\hat{\MATRIX{U}}}}
\newcommand\mhV{{\hat{\MATRIX{V}}}}
\newcommand\mhW{{\hat{\MATRIX{W}}}}
\newcommand\mhX{{\hat{\MATRIX{X}}}}
\newcommand\mhY{{\hat{\MATRIX{Y}}}}
\newcommand\mhZ{{\hat{\MATRIX{Z}}}}

\newcommand\mGamma{{\MATRIX{\Gamma}}}
\newcommand\mDelta{{\MATRIX{\Delta}}}
\newcommand\mTheta{{\MATRIX{\Theta}}}
\newcommand\mLambda{{\MATRIX{\Lambda}}}
\newcommand\mXi{{\MATRIX{\Xi}}}
\newcommand\mPi{{\MATRIX{\Pi}}}
\newcommand\mSigma{{\MATRIX{\Sigma}}}
\newcommand\mUpsilon{{\MATRIX{\Upsilon}}}
\newcommand\mPhi{{\MATRIX{\Phi}}}
\newcommand\mPsi{{\MATRIX{\Psi}}}
\newcommand\mOmega{{\MATRIX{\Omega}}}


% Tensor
% \newcommand\tens[1]{\symbfsf{#1}}
% \newcommand\tA{{\tens{A}}}
% \newcommand\tB{{\tens{B}}}
% \newcommand\tC{{\tens{C}}}
% \newcommand\tD{{\tens{D}}}
% \newcommand\tE{{\tens{E}}}
% \newcommand\tF{{\tens{F}}}
% \newcommand\tG{{\tens{G}}}
% \newcommand\tH{{\tens{H}}}
% \newcommand\tI{{\tens{I}}}
% \newcommand\tJ{{\tens{J}}}
% \newcommand\tK{{\tens{K}}}
% \newcommand\tL{{\tens{L}}}
% \newcommand\tM{{\tens{M}}}
% \newcommand\tN{{\tens{N}}}
% \newcommand\tO{{\tens{O}}}
% \newcommand\tP{{\tens{P}}}
% \newcommand\tQ{{\tens{Q}}}
% \newcommand\tR{{\tens{R}}}
% \newcommand\tS{{\tens{S}}}
% \newcommand\tT{{\tens{T}}}
% \newcommand\tU{{\tens{U}}}
% \newcommand\tV{{\tens{V}}}
% \newcommand\tW{{\tens{W}}}
% \newcommand\tX{{\tens{X}}}
% \newcommand\tY{{\tens{Y}}}
% \newcommand\tZ{{\tens{Z}}}


% Graph
\newcommand\gA{{\mathcal{A}}}
\newcommand\gB{{\mathcal{B}}}
\newcommand\gC{{\mathcal{C}}}
\newcommand\gD{{\mathcal{D}}}
\newcommand\gE{{\mathcal{E}}}
\newcommand\gF{{\mathcal{F}}}
\newcommand\gG{{\mathcal{G}}}
\newcommand\gH{{\mathcal{H}}}
\newcommand\gI{{\mathcal{I}}}
\newcommand\gJ{{\mathcal{J}}}
\newcommand\gK{{\mathcal{K}}}
\newcommand\gL{{\mathcal{L}}}
\newcommand\gM{{\mathcal{M}}}
\newcommand\gN{{\mathcal{N}}}
\newcommand\gO{{\mathcal{O}}}
\newcommand\gP{{\mathcal{P}}}
\newcommand\gQ{{\mathcal{Q}}}
\newcommand\gR{{\mathcal{R}}}
\newcommand\gS{{\mathcal{S}}}
\newcommand\gT{{\mathcal{T}}}
\newcommand\gU{{\mathcal{U}}}
\newcommand\gV{{\mathcal{V}}}
\newcommand\gW{{\mathcal{W}}}
\newcommand\gX{{\mathcal{X}}}
\newcommand\gY{{\mathcal{Y}}}
\newcommand\gZ{{\mathcal{Z}}}


% Sets
\newcommand\sA{{\mathbb{A}}}
\newcommand\sB{{\mathbb{B}}}
\newcommand\sC{{\mathbb{C}}}
\newcommand\sD{{\mathbb{D}}}
% Don't use a set called E, because this would be the same as our symbol
% for expectation.
\newcommand\sF{{\mathbb{F}}}
\newcommand\sG{{\mathbb{G}}}
\newcommand\sH{{\mathbb{H}}}
\newcommand\sI{{\mathbb{I}}}
\newcommand\sJ{{\mathbb{J}}}
\newcommand\sK{{\mathbb{K}}}
\newcommand\sL{{\mathbb{L}}}
\newcommand\sM{{\mathbb{M}}}
\newcommand\sN{{\mathbb{N}}}
\newcommand\sO{{\mathbb{O}}}
\newcommand\sP{{\mathbb{P}}}
\newcommand\sQ{{\mathbb{Q}}}
\newcommand\sR{{\mathbb{R}}}
\newcommand\sS{{\mathbb{S}}}
\newcommand\sT{{\mathbb{T}}}
\newcommand\sU{{\mathbb{U}}}
\newcommand\sV{{\mathbb{V}}}
\newcommand\sW{{\mathbb{W}}}
\newcommand\sX{{\mathbb{X}}}
\newcommand\sY{{\mathbb{Y}}}
\newcommand\sZ{{\mathbb{Z}}}

\renewcommand\sA{{\mathcal{A}}}
\renewcommand\sS{{\mathcal{S}}}
\renewcommand\sQ{{\mathcal{Q}}}

% Random variables
\newcommand\RandomVariable[1]{\symit{#1}}

\newcommand\rA{{\RandomVariable{A}}}
\newcommand\rB{{\RandomVariable{B}}}
\newcommand\rC{{\RandomVariable{C}}}
\newcommand\rD{{\RandomVariable{D}}}
\newcommand\rE{{\RandomVariable{E}}}
\newcommand\rF{{\RandomVariable{F}}}
\newcommand\rG{{\RandomVariable{G}}}
\newcommand\rH{{\RandomVariable{H}}}
\newcommand\rI{{\RandomVariable{I}}}
\newcommand\rJ{{\RandomVariable{J}}}
\newcommand\rK{{\RandomVariable{K}}}
\newcommand\rL{{\RandomVariable{L}}}
\newcommand\rM{{\RandomVariable{M}}}
\newcommand\rN{{\RandomVariable{N}}}
\newcommand\rO{{\RandomVariable{O}}}
\newcommand\rP{{\RandomVariable{P}}}
\newcommand\rQ{{\RandomVariable{Q}}}
\newcommand\rR{{\RandomVariable{R}}}
\newcommand\rS{{\RandomVariable{S}}}
\newcommand\rT{{\RandomVariable{T}}}
\newcommand\rU{{\RandomVariable{U}}}
\newcommand\rV{{\RandomVariable{V}}}
\newcommand\rW{{\RandomVariable{W}}}
\newcommand\rX{{\RandomVariable{X}}}
\newcommand\rY{{\RandomVariable{Y}}}
\newcommand\rZ{{\RandomVariable{Z}}}

% Random vectors
\newcommand\RandomVector[1]{\symbf{#1}}

\newcommand\rvA{{\RandomVector{A}}}
\newcommand\rvB{{\RandomVector{B}}}
\newcommand\rvC{{\RandomVector{C}}}
\newcommand\rvD{{\RandomVector{D}}}
\newcommand\rvE{{\RandomVector{E}}}
\newcommand\rvF{{\RandomVector{F}}}
\newcommand\rvG{{\RandomVector{G}}}
\newcommand\rvH{{\RandomVector{H}}}
\newcommand\rvI{{\RandomVector{I}}}
\newcommand\rvJ{{\RandomVector{J}}}
\newcommand\rvK{{\RandomVector{K}}}
\newcommand\rvL{{\RandomVector{L}}}
\newcommand\rvM{{\RandomVector{M}}}
\newcommand\rvN{{\RandomVector{N}}}
\newcommand\rvO{{\RandomVector{O}}}
\newcommand\rvP{{\RandomVector{P}}}
\newcommand\rvQ{{\RandomVector{Q}}}
\newcommand\rvR{{\RandomVector{R}}}
\newcommand\rvS{{\RandomVector{S}}}
\newcommand\rvT{{\RandomVector{T}}}
\newcommand\rvU{{\RandomVector{U}}}
\newcommand\rvV{{\RandomVector{V}}}
\newcommand\rvW{{\RandomVector{W}}}
\newcommand\rvX{{\RandomVector{X}}}
\newcommand\rvY{{\RandomVector{Y}}}
\newcommand\rvZ{{\RandomVector{Z}}}

% Vectors

\newcommand\ba{{\bar{a}}}
\newcommand\bb{{\bar{b}}}
\newcommand\bc{{\bar{c}}}
\newcommand\bd{{\bar{d}}}
\newcommand\be{{\bar{e}}}

\newcommand\bg{{\bar{g}}}
\newcommand\bh{{\bar{h}}}
\newcommand\bi{{\bar{i}}}
\newcommand\bj{{\bar{j}}}
\newcommand\bk{{\bar{k}}}
\newcommand\bl{{\bar{l}}}

\newcommand\bn{{\bar{n}}}
\newcommand\bo{{\bar{o}}}
\newcommand\bp{{\bar{p}}}
\newcommand\bq{{\bar{q}}}
% \newcommand\br{{\bar{r}}}
\newcommand\bs{{\bar{s}}}
\newcommand\bt{{\bar{t}}}
\newcommand\bu{{\bar{u}}}
\newcommand\bv{{\bar{v}}}
\newcommand\bw{{\bar{w}}}
\newcommand\bx{{\bar{x}}}
\newcommand\by{{\bar{y}}}
\newcommand\bz{{\bar{z}}}

\newcommand\balpha{{\bar{\alpha}}}
\newcommand\bbeta{{\bar{\beta}}}
\newcommand\bgamma{{\bar{\gamma}}}
\newcommand\bdelta{{\bar{\delta}}}
\newcommand\bepsilon{{\bar{\epsilon}}}
\newcommand\btheta{{\bar{\theta}}}
\newcommand\biota{{\bar{\iota}}}
\newcommand\bkappa{{\bar{\kappa}}}
\newcommand\blambda{{\bar{\lambda}}}
\newcommand\bmu{{\bar{\mu}}}
\newcommand\bnu{{\bar{\nu}}}
\newcommand\bxi{{\bar{\xi}}}
\newcommand\bpi{{\bar{\pi}}}
\newcommand\brho{{\bar{\rho}}}
\newcommand\bsigma{{\bar{\sigma}}}
\newcommand\btau{{\bar{\tau}}}
\newcommand\bupsilon{{\bar{\upsilon}}}
\newcommand\bphi{{\bar{\phi}}}
\newcommand\bchi{{\bar{\chi}}}
\newcommand\bpsi{{\bar{\psi}}}
\newcommand\bomega{{\bar{\omega}}}

\newcommand\bA{{\bar{A}}}
\newcommand\bB{{\bar{B}}}
\newcommand\bC{{\bar{C}}}
\newcommand\bD{{\bar{D}}}
\newcommand\bF{{\bar{F}}}
\newcommand\bG{{\bar{G}}}
\newcommand\bH{{\bar{H}}}
\newcommand\bI{{\bar{I}}}
\newcommand\bJ{{\bar{J}}}
\newcommand\bK{{\bar{K}}}
\newcommand\bL{{\bar{L}}}
\newcommand\bM{{\bar{M}}}
\newcommand\bN{{\bar{N}}}
\newcommand\bO{{\bar{O}}}
\newcommand\bP{{\bar{P}}}
\newcommand\bQ{{\bar{Q}}}
\newcommand\bR{{\bar{R}}}
\newcommand\bS{{\bar{S}}}
\newcommand\bT{{\bar{T}}}
\newcommand\bU{{\bar{U}}}
\newcommand\bV{{\bar{V}}}
\newcommand\bW{{\bar{W}}}
\newcommand\bX{{\bar{X}}}
\newcommand\bY{{\bar{Y}}}
\newcommand\bZ{{\bar{Z}}}
\newcommand\bLambda{{\bar{\Lambda}}}
\newcommand\bGamma{{\bar{\Gamma}}}

% bv
\newcommand\bva{{\bar{\va}}}
\newcommand\bvb{{\bar{\vb}}}
\newcommand\bvc{{\bar{\vc}}}
\newcommand\bvd{{\bar{\vd}}}
\newcommand\bve{{\bar{\ve}}}

\newcommand\bvg{{\bar{\vg}}}
\newcommand\bvh{{\bar{\vh}}}
\newcommand\bvi{{\bar{\vi}}}
\newcommand\bvj{{\bar{\vj}}}
\newcommand\bvk{{\bar{\vk}}}
\newcommand\bvl{{\bar{\vl}}}

\newcommand\bvn{{\bar{\vn}}}
\newcommand\bvo{{\bar{\vo}}}
\newcommand\bvp{{\bar{\vp}}}
\newcommand\bvq{{\bar{\vq}}}
\newcommand\bvr{{\bar{\vr}}}
\newcommand\bvs{{\bar{\vs}}}
\newcommand\bvt{{\bar{\vt}}}
\newcommand\bvu{{\bar{\vu}}}
\newcommand\bvv{{\bar{\vv}}}
\newcommand\bvw{{\bar{\vw}}}
\newcommand\bvx{{\bar{\vx}}}
\newcommand\bvy{{\bar{\vy}}}
\newcommand\bvz{{\bar{\vz}}}

\newcommand\bvalpha{{\bar{\valpha}}}
\newcommand\bvbeta{{\bar{\vbeta}}}
\newcommand\bvgamma{{\bar{\vgamma}}}
\newcommand\bvdelta{{\bar{\vdelta}}}
\newcommand\bvepsilon{{\bar{\vepsilon}}}
\newcommand\bvtheta{{\bar{\vtheta}}}
\newcommand\bviota{{\bar{\viota}}}
\newcommand\bvkappa{{\bar{\vkappa}}}
\newcommand\bvlambda{{\bar{\vlambda}}}
\newcommand\bvmu{{\bar{\vmu}}}
\newcommand\bvnu{{\bar{\vnu}}}
\newcommand\bvxi{{\bar{\vxi}}}
\newcommand\bvpi{{\bar{\vpi}}}
\newcommand\bvrho{{\bar{\vrho}}}
\newcommand\bvsigma{{\bar{\vsigma}}}
\newcommand\bvtau{{\bar{\vtau}}}
\newcommand\bvupsilon{{\bar{\vupsilon}}}
\newcommand\bvphi{{\bar{\vphi}}}
\newcommand\bvchi{{\bar{\vchi}}}
\newcommand\bvpsi{{\bar{\vpsi}}}
\newcommand\bvomega{{\bar{\vomega}}}

% \newcommand\bvA{{\bar{A}}}
% \newcommand\bvB{{\bar{B}}}
% \newcommand\bvC{{\bar{C}}}
% \newcommand\bvD{{\bar{D}}}
% \newcommand\bvF{{\bar{F}}}
% \newcommand\bvG{{\bar{G}}}
% \newcommand\bvH{{\bar{H}}}
% \newcommand\bvI{{\bar{I}}}
% \newcommand\bvJ{{\bar{J}}}
% \newcommand\bvK{{\bar{K}}}
% \newcommand\bvL{{\bar{L}}}
\newcommand\bmM{{\bar{\mM}}}
% \newcommand\bvN{{\bar{N}}}
% \newcommand\bvO{{\bar{O}}}
% \newcommand\bvP{{\bar{P}}}
% \newcommand\bvQ{{\bar{Q}}}
% \newcommand\bvR{{\bar{R}}}
% \newcommand\bvS{{\bar{S}}}
% \newcommand\bvT{{\bar{T}}}
% \newcommand\bvU{{\bar{U}}}
% \newcommand\bvV{{\bar{V}}}
% \newcommand\bvW{{\bar{W}}}
% \newcommand\bvX{{\bar{X}}}
% \newcommand\bvY{{\bar{Y}}}
% \newcommand\bvZ{{\bar{Z}}}
\newcommand\bmLambda{{\bar{\mLambda}}}
\newcommand\bmGamma{{\bar{\mGamma}}}



\newcommand\tlambda{{\tilde{\lambda}}}
\newcommand\tLambda{{\tilde{\Lambda}}}

\newcommand\tmu{{\tilde{\mu}}}

\newcommand\laplace{\mathrm{Laplace}} % Laplace distribution

\newcommand\bigO{\mathcal{O}}
\newcommand\semiS[1]{{\mathbb{S}_{+}^{#1}}}
\newcommand\Prob{\mathbf{Prob}}
% \newcommand\I{\mathbb{I}} % Indicator function
\newcommand\givenp[1][]{{\:#1\vert\:}}
\newcommand\Diag{{\mathrm{Diag}}}
\newcommand\Real{{\mathrm{Re}}}
\newcommand\Imag{{\mathrm{Im}}}
\DeclareMathOperator*{\E}{\mathbb{E}}
% \newcommand\E{\mathbb{E}}
\newcommand\Ls{\mathcal{L}}
% \newcommand\R{\mathbb{R}}
% \newcommand\C{\mathbb{C}}
\newcommand\emp{\tilde{p}}
\newcommand\lr{\alpha}
\newcommand\reg{\lambda}
\newcommand\rect{\mathrm{rectifier}}
\newcommand\softmax{\mathrm{softmax}}
\newcommand\sigmoid{\sigma}
\newcommand\softplus{\zeta}
\newcommand\KL{D_{\mathrm{KL}}}
\newcommand\Var{\mathrm{Var}}
\newcommand\standarderror{\mathrm{SE}}
\newcommand\Cov{\mathrm{Cov}}
% Wolfram Mathworld says $L^2$ is for function spaces and $\ell^2$ is for vectors
% But then they seem to use $L^2$ for vectors throughout the site, and so does
% wikipedia.
\newcommand\normlzero{L^0}
\newcommand\normlone{L^1}
\newcommand\normltwo{L^2}
\newcommand\normlp{L^p}
\newcommand\normmax{L^\infty}
\newcommand\imag{{\mathbf{i}}}
% \newcommand\T{{\mathsf{T}}}
\newcommand\op{{\mathrm{op}}}

\DeclareMathOperator*{\argmax}{arg\,max}
\DeclareMathOperator*{\argmin}{arg\,min}

\DeclareMathOperator{\sign}{sign}
\DeclareMathOperator{\Tr}{\mathbf{Tr}}
\DeclareMathOperator{\interior}{int\,}
\DeclareMathOperator{\dom}{dom\,}
\let\ab\allowbreak


\newtheorem{assumption}{Assumption}
\newtheorem{definition}{Definition}
\newtheorem{theorem}{Theorem}
\newtheorem{lemma}[theorem]{Lemma}
\newtheorem{proposition}[theorem]{Proposition}
\newtheorem{corollary}[theorem]{Corollary}
\newtheorem{remark}{Remark}
\numberwithin{assumption}{section}
\numberwithin{definition}{section}
\numberwithin{theorem}{section}
\numberwithin{remark}{section}

\usepackage{hyperref}
\hypersetup{colorlinks,citecolor=blue,linktocpage,breaklinks=true}

\begin{document}

\title{Affine-Invariant Global Non-Asymptotic Convergence Analysis of BFGS under Self-Concordance}

\author{Qiujiang Jin\thanks{UT Austin  \{qiujiangjin0@gmail.com\}} \qquad Aryan Mokhtari\thanks{UT Austin and Google Research  \{mokhtari@austin.utexas.edu\}}}

\maketitle

\begin{abstract}
In this paper, we establish global non-asymptotic convergence guarantees for the BFGS quasi-Newton method without requiring strong convexity or the Lipschitz continuity of the gradient or Hessian. Instead, we consider the setting where the objective function is strictly convex and strongly self-concordant. For an arbitrary initial point and any arbitrary positive-definite initial Hessian approximation, we prove global linear and superlinear convergence guarantees for BFGS when the step size is determined using a line search scheme satisfying the weak Wolfe conditions.  Moreover, all our global guarantees are affine-invariant, with the convergence rates depending solely on the initial error and the strongly self-concordant constant. Our results extend the global non-asymptotic convergence theory of BFGS beyond traditional assumptions and, for the first time, establish affine-invariant convergence guarantees—aligning with the inherent affine invariance of the BFGS method.
\end{abstract}

\section{Introduction}\label{sec:introduction}

In this paper, we consider the convex optimization problem
\begin{equation}
    \min_{x \in \mathbb{R}^d} f(x),
\end{equation}
where the function $f$ is twice differentiable and \textit{strictly} convex. We focus on quasi-Newton methods—iterative optimization algorithms that approximate the Hessian and its inverse using gradient information, making them efficient for large-scale problems where computing the Hessian is costly. Different variants update the Hessian approximation in distinct ways.
 The most famous quasi-Newton methods include the Davidon-Fletcher-Powell (DFP) method \cite{davidon1959variable,fletcher1963rapidly}, the Broyden-Fletcher-Goldfarb-Shanno (BFGS) method \cite{broyden1970convergence,fletcher1970new,goldfarb1970family,shanno1970conditioning}, the Symmetric Rank-One~(SR1) method \cite{conn1991convergence,khalfan1993theoretical},  and the Broyden method \cite{broyden1965class}. There are also variants of these methods, including limited memory BFGS~\cite{liu1989limited, nocedal1980updating}, randomized quasi-Newton methods \cite{gower2016stochastic,gower2017randomized,kovalev2020fast,lin2021greedy, lin2022explicit}, and greedy quasi-Newton methods \cite{lin2021greedy,lin2022explicit,rodomanov2020greedy,ji2023greedy}.%, and the online learning-guided variants~\cite{jiang2023online, jiang2023accelerated}. 

In this paper, we focus exclusively on the BFGS method, one of the most widely used and well-regarded quasi-Newton algorithms. Specifically, we analyze its convergence guarantees in the setting where the objective function is strictly convex and self-concordant and establish non-asymptotic guarantees for this case. Before highlighting our contributions, we first provide a summary of the existing convergence guarantees for BFGS as established in prior work.

\textbf{Classic asymptotic guarantees.} The local asymptotic superlinear convergence of quasi-Newton methods, including BFGS, has been established in several works~\cite{broyden1973local,dennis1974characterization,griewank1982local,dennis1989convergence,yuan1991modified,al1998global,li1999globally,yabe2007local,mokhtari2017iqn,gao2019quasi}. Similarly, their global convergence under globalization strategies like line search and trust-region methods has been analyzed~\cite{khalfan1993theoretical,QN_tool,powell1971convergence, Powell, byrd1987global, byrd1996analysis, SR12024}. However, these results are asymptotic and lack showing explicit rates.

\textbf{Non-asymptotic guarantees under stronger assumptions.} Recently, there were several breakthroughs regarding the non-asymptotic local superlinear convergence analysis of BFGS including~\cite{rodomanov2020rates,rodomanov2020ratesnew,ye2023towards,qiujiang2020quasinewton} for the case that the objective function is strongly convex. More precisely, these works established an explicit superlinear rate of $\mathcal{O}({1}/{\sqrt{t}})^t$ under the assumptions of strong convexity and Lipschitz continuity of the gradient and Hessian, given that the initial point is within a local neighborhood of the optimum and the initial Hessian approximation satisfies certain conditions. Later, these local analyses were extended and non-asymptotic global convergence rates of BFGS were established in~\cite{krutikov2023convergence,antonglobal,jin2024non,qiujiang2024quasinewton} under similar assumptions on the objective function. In particular, \cite{jin2024non} established global explicit superlinear convergence guarantees of the whole convex class of Broyden's family of quasi-Newton methods including both BFGS and DFP with step size satisfying the exact line search schemes. In a follow up work \cite{qiujiang2024quasinewton}, the explicit global convergence rates for BFGS was established when deployed with an inexact line search satisfying the Armijo-Wolfe conditions. Specifically, these works show that when the objective is $\mu$-strongly convex, its gradient is $L$-Lipschitz smooth, and its Hessian is $K$-Lipschitz continuous, a global linear convergence rate of $(1 - {1}/{\kappa})^t$ can be achieved—matching that of gradient descent, where $\kappa = L/\mu$ is the condition number. Moreover, global superlinear convergence rates of $({(d\kappa + C_0\kappa)}/{t})^t$ and $({(C_0 d\log{\kappa} + C_0\kappa)}/{t})^t$ were established under specific choices of the initial Hessian approximation, where $d$ is the problem dimension,  and $C_0$ is the initial function value gap between the initial iterate $x_0$ and the unique optimal solution $x_*$.

%In these works, it was shown when the objective function is $\mu-$strongly convex, the gradient is $L-$Lipschitz smooth and the Hessian is $K-$Lipschitz continuous, a global linear convergence rate of $(1 - {1}/{\kappa})^{t}$ was proven, which matches the global linear rate of the gradient descent method, and the global superlinear convergence rates of $({(d\kappa + C_0\kappa)}/{t})^{t}$ and $({(C_0 d\log{\kappa} + C_0\kappa)}/{t})^{t}$ were established for particular choice of the initial Hessian approximation matrices, where $d$ is the dimension, $\kappa = L/\mu$ is the condition number of the objective function and $C_0$ is the initial function value difference between $x_0$ and $x_*$. 

While these results represent significant progress in studying quasi-Newton methods, the established non-asymptotic guarantees for BFGS, and most quasi-Newton methods in general, have two major limitations. First, these results rely on relatively strong assumptions that may not hold in many practical settings. For instance, in the case of logistic regression, the loss is strictly convex but not necessarily strongly convex. Similarly, a log-barrier function does not satisfy the global Lipschitz condition for gradient. Second, all previously established non-asymptotic convergence rates for BFGS are not affine invariant, as they depend on parameters such as the strong convexity constant \(\mu\), gradient Lipschitz constant \(L\), and Hessian Lipschitz constant \(K\), all of which vary under a change of basis or coordinate system in \(\mathbb{R}^d\). In contrast, BFGS is affine invariant with respect to linear transformations of the variables. This means that the convergence behavior of BFGS remains unaffected by the choice of coordinate system and instead depends solely on the topological structure of  \(f\). 

\noindent\textbf{Contributions:} We aim to address the discussed issues, and our main contributions are as follows:
\begin{itemize}

\item 
    We establish global non-asymptotic linear and superlinear convergence rates for BFGS without requiring strong convexity or Lipschitz continuity of the gradient or Hessian. Instead, we consider functions that are strictly convex and strongly self-concordant. Our analysis provides explicit global convergence guarantees for BFGS when the step size is selected via a line search satisfying the weak Wolfe conditions. These guarantees hold for any initial point \(x_0\) and any positive-definite initial Hessian approximation \(B_0\). %Notably, our results show that strong convexity is not required for BFGS to achieve superlinear convergence.

\item 
   We derive explicit convergence rates for the BFGS  method that are affine invariant. Specifically, our results show that both global linear and superlinear convergence rates depend solely on the strongly self-concordant constant, which remains invariant under linear transformations of the variables. To the best of our knowledge, these are the first theoretical convergence rates consistent with the affine invariance property of the BFGS method, reflecting its independence from the choice of coordinate system. 

\end{itemize}

%\textcolor{red}{
%The primary contribution and novelty of our submission lie in extending the global non-asymptotic convergence analysis of quasi-Newton methods to the setting of self-concordant functions without requiring strong convexity and gradient smoothness. Our results are also consistent with the affine-invariance property of the BFGS method. It is well known that global convergence rates for Newton's method under strong convexity and Lipschitz continuity of the gradient and Hessian have been established for many years. However, in the context of interior-point methods involving log-barrier functions, the strong convexity assumption does not hold. To address this, Nesterov and Nemirovski developed a global convergence theory for Newton’s method under the self-concordant framework—a major breakthrough that enabled complexity bounds for interior-point methods. This development was especially important as it preserved the affine-invariance property of Newton’s method. Extending the theory of BFGS to this setting is equally fundamental. Unlike the classical analysis of Newton's method, which was insufficient for self-concordant functions, prior convergence results for BFGS, such as those in \cite{jin2024non, qiujiang2024quasinewton}, do not apply in this regime due to the absence of strong convexity. To make BFGS applicable to interior-point methods and consistent with its affine-invariance property, it is crucial to establish a superlinear convergence analysis under self-concordance. In fact, the analysis in this setting is fundamentally different from the standard case involving strong convexity and Lipschitz continuity, just as it was for Newton's method. Furthermore, the proofs of Lemmas A.5, A.7, A.8, and A.9 are entirely original and independent of previous literature. These results provide new theoretical insights that are not present in earlier works, including \cite{jin2024non, qiujiang2024quasinewton}.
%}

%\textcolor{red}{
%Regarding the relation to the connections to previous papers \cite{jin2024non, qiujiang2024quasinewton}, we emphasize that our submission is not derivative of them. Rather, our paper extends the analysis in \cite{jin2024non, qiujiang2024quasinewton} to a significantly more general setting: we remove the assumptions of strong convexity and Lipschitz continuity of the gradient and Hessian. Once these assumptions are dropped, the theoretical tools and techniques developed in \cite{jin2024non, qiujiang2024quasinewton} are no longer applicable. Therefore, our analysis requires entirely new proof strategies. Extending a method’s analysis to a more general setting—especially when the existing techniques break down—is not plagiarism. To the best of our knowledge, there is no prior result establishing the global iteration complexity of BFGS under the assumptions we consider (Assumption 2.2), and thus we strongly believe the reviewer’s claim of plagiarism is unfounded.
%}

%We summarize the structure of this paper as well as some important notations as below:

%\noindent\textbf{Outline.} In Section~\ref{sec:background}, we provide an overview of the BFGS method with inexact line search, outline our assumptions, and introduce some preliminary lemmas for the convergence analysis. In Section~\ref{sec:linear}, we establish the global linear convergence rates of BFGS. Based on these linear convergence rates, Section~\ref{sec:superlinear} details our global superlinear convergence results. In Section~\ref{sec:complexity}, we show the iteration complexity and the line search complexity of BFGS under the considered assumptions. Section~\ref{sec:experiments} displays our numerical experiments. Finally, we state our concluding remarks in Section~\ref{sec:conclusion}.

\noindent\textbf{Notation.} We denote the $l_2$-norm by $\|\cdot\|$ and the set of $d \times d$ symmetric positive definite matrices by $\mathbb{S}^d_{++}$. We write $A \preceq B$ if $B - A$ is positive semi-definite, and $A \prec B$ if it is positive definite. The trace and determinant of matrix $A$ are represented as $\Tr(A)$ and $\Det(A)$, respectively. For function $f$ that is strictly convex, we define the weighted norm $\|.\|_x$ as $\|u\|_x := \sqrt{u^\top \nabla^2{f(x)} u}$

\section{Background and Preliminaries}\label{sec:background}

%\subsection{BFGS Quasi-Newton Method}

In this section, we provide a brief overview of the BFGS quasi-Newton method. At iteration \( t \), \( x_t \) denotes the current iterate, \( g_t = \nabla f(x_t) \) the gradient of the objective function, and \( B_t \) the Hessian approximation matrix. The general template of quasi-Newton methods update is given by
\begin{equation}\label{QN_update}
    x_{t + 1} = x_t + \eta_t d_t, \qquad d_t = -B_t^{-1} g_t,
\end{equation}
where $\eta_t > 0$ is the step size. By defining the variable difference and the gradient difference as
\begin{equation}\label{QN_difference}
    s_t := x_{t + 1} - x_t, \qquad y_t := \nabla f(x_{t + 1}) - \nabla f(x_t),
\end{equation}
we can present the Hessian approximation matrix update for BFGS as follows:
\begin{equation}\label{eq:BFGS_update}
    B_{t + 1} := B_t - \frac{B_t s_t s_t^\top B_t}{s_t^\top B_t s_t} + \frac{y_t y_t^\top}{s_t^\top y_t}.
\end{equation}
Further, if we define the inverse of Hessian approximation  as $H_t := B_t^{-1}$, using the Sherman-Morrison-Woodbury formula, we have $H_{t+1} = (I-\frac{s_t y_t^\top}{y_t^\top s_t}) H_t (I-\frac{ y_t s_t^\top}{s_t^\top y_t}) +\frac{s_t s_t^\top}{y_t^\top s_t}$. Note that if the function $f$ is strictly convex -- as considered in this paper -- and the initial Hessian approximation matrix is positive definite, then $B_{t} \in \mathbb{S}^d_{++}$ for any iterations $t > 0$ (Chapter 6 \cite{nocedal2006numerical}). In this paper, we focus on the analysis of BFGS when $\eta_t$ is selected based on the Armijo-Wolfe conditions, given by
\begin{align}
    f(x_t + \eta_t d_t) & \leq f(x_t) + \alpha \eta_t\nabla{f(x_t)}^\top d_t, \label{sufficient_decrease}\\
    \nabla{f(x_t + \eta_t d_t)}^\top d_t & \geq \beta \nabla{f(x_t)}^\top d_t, \label{curvature_condition}
\end{align}
where $\alpha$ and $\beta$ are the line search parameters, satisfying $0 < \alpha < \beta < 1$ and $0 < \alpha < {1}/{2}$. %Note that there always exists $\eta_t > 0$ satisfying the Armijo-Wolfe conditions as long as $f$ is bounded below; see, e.g.,  Lemma 3.1 in \cite{nocedal2006numerical}.

\noindent \textbf{Affine Invariance property of BFGS.} From \cite{affine1, affine2}, it is known that the iterates of BFGS are \textit{affine invariant}. This property underscores the necessity of an analysis framework aligned with affine invariance, which is the main focus of our paper. We state the following proposition for completeness. %(check proof in Appendix~\ref{proof_affine})

\begin{proposition}\label{BFGS_affine}
    Let the iterations $\{x_t\}_{t=0}^{+\infty}$ be generated by the BFGS algorithm applied to the objective function $f(x)$, as defined in \eqref{QN_update}-\eqref{eq:BFGS_update}. Consider the iterates $\{\dot{x}_t\}_{t=0}^{+\infty}$ produced by applying BFGS to the transformed function $\phi(x) = f(Ax)$, where $A \in \mathbb{R}^{d \times d}$ is a non-singular matrix. Assume that the initializations satisfy $\dot{x}_0 = A^{-1}x_0$ and $\dot{B}_0 = A^\top B_0 A$. Then, for any $t \geq 0$, the following relationships hold: $\dot{x}_t = A^{-1}x_t$, $\dot{B}_t = A^\top B_t A$ and $\phi(\dot{x}_t) = f(x_t)$.
\end{proposition}

%Although this paper focuses specifically on BFGS, it is worth noting that other quasi-Newton methods, such as the SR1 method, the DFP method, and the broader convex class of Broyden's family, are also affine invariant. Their proofs follow a similar argument to that presented in Proposition~\ref{BFGS_affine}.

\subsection{Assumptions}
Next, we state our assumptions and compare them with those used in prior work.

\begin{assumption}\label{ass_self_concodant}
The function $f$ satisfies the following conditions: (i) it is twice differentiable and strictly convex, and (ii) it is strongly self-concordant with parameter $M > 0$, i.e., for any $x, y, z \in \mathbb{R}^d$
\begin{equation}
\nabla^2{f(x)} - \nabla^2{f(y)} \preceq M\|x - y\|_{z}\nabla^2{f(y)}.
\end{equation}
\end{assumption}
Our first assumption requires the objective function to be strictly convex, i.e., $\nabla^2 f(x) \succ 0$. This is indeed a weaker condition than the strong convexity assumptions used in prior works that establish non-asymptotic guarantees for BFGS, such as \cite{rodomanov2020rates,rodomanov2020ratesnew,ye2023towards,qiujiang2020quasinewton,krutikov2023convergence,antonglobal,jin2024non,qiujiang2024quasinewton}. The second condition concerns strong self-concordance, which defines a subclass of self-concordant functions. Specifically, if $f$ is $M$-strongly self-concordant, then it is also $M/2$-self-concordant. To see this, fix $x \in \mathbb{R}^d$ and $u \in \mathbb{R}^d$. The inequality $u^\top (\nabla^2{f(x + tu)} - \nabla^2{f(x)}) u \leq tM\|u\|^{3}_{x}$ holds, and dividing by $t$ and taking the limit as $t \to 0$ yields $D^{3}f(x)[u, u, u] \leq M\|u\|^{3}_{x}$. A symmetric argument shows $|D^{3}f(x)[u, u, u]| \leq M\|u\|^{3}_{x}$, implying that $f$ is self-concordant with parameter $M/2$. Moreover, Theorem 5.1.2 of \cite{convexlecture} shows that the strong self-concordance parameter $M$ is affine invariant: for any non-singular $A \in \mathbb{R}^{d\times d}$, the function $\phi(x) = f(Ax)$ remains $M$-strongly self-concordant.

Next, we explain why our assumptions are strictly weaker than the more common conditions of strong convexity, Lipschitz gradient, and Lipschitz Hessian. Prior work (e.g., Example 4.1 in~\cite{rodomanov2020greedy}) shows that if a function is strongly convex and its Hessian is Lipschitz with respect to a matrix $B \succeq 0$, then it is also strongly self-concordant. However, the converse does not hold: strong self-concordance does not imply strong convexity, gradient smoothness, or Lipschitz Hessian continuity.

%Regarding the second assumption, we should clarify that strongly self-concordant functions are a specific subclass of self-concordant functions. More precisely, if \(f\) is \(M\)-strongly self-concordant, it can be verified that \(f\) is also \(M/2\)-self-concordant. To illustrate this, consider a point \(x \in \mathbb{R}^d\) and a direction \(u \in \mathbb{R}^d\). We have that $u^\top (\nabla^2{f(x + tu)} - \nabla^2{f(x)}) u \leq tM\|u\|^{3}_{x}$. Dividing this inequality by $t$ and computing the limit as $t$ approach $0$, we obtain that $D^{3}f(x)[u, u, u] \leq M\|u\|^{3}_{x}$. Similarly, we can prove that $|D^{3}f(x)[u, u, u]| \leq M\|u\|^{3}_{x}$. Thus, function $f$ is self-concordant with parameter ${M}/{2}$. Also, we should add that, as shown in Theorem 5.1.2 of \cite{convexlecture}, strongly self-concordant parameter $M$ is affine invariant, i.e., if we define function $\phi(x) = f(Ax)$ for any non-singular matrix $A \in \mathbb{R}^{d\times d}$, then it can be verified that $\phi(x)$ is also strongly self-concordant with  parameter $M$.

%{\color{blue}{Instead of the following remarks you need a subsection in which in detail you explain why our assumptions are weaker than assuming strong convexity, grad lip and Hessian lip. You provide detailed examples for which our assumptions are satisfied but these conditions do not hold. you can use part of the text in the following remarks for this purpose, if needed. another function that you can mention is log barrier $f(x)=- log(1-x^2)$, for which $M=4$ when $\|x-y\|\leq 1/2$, while both grad and hessian are not Lipschitz}}

%The first condition requires the objective function to be strictly convex, i.e., \$\nabla^2 f(x) \succ 0\$, which is weaker than the strong convexity assumptions typically used in prior works that establish non-asymptotic guarantees for BFGS\~\cite{rodomanov2020rates,rodomanov2020ratesnew,ye2023towards,qiujiang2020quasinewton,krutikov2023convergence,antonglobal,jin2024non,qiujiang2024quasinewton}.

%The second condition concerns strong self-concordance, which defines a subclass of self-concordant functions. Specifically, if \$f\$ is \$M\$-strongly self-concordant, then it is also \$M/2\$-self-concordant. To see this, fix \$x \in \mathbb{R}^d\$ and \$u \in \mathbb{R}^d\$. The inequality \$u^\top (\nabla^2{f(x + tu)} - \nabla^2{f(x)}) u \leq tM|u|^{3}*{x}\$ holds, and dividing by \$t\$ and taking the limit as \$t \to 0\$ yields \$D^{3}f(x)\[u, u, u] \leq M|u|^{3}*{x}\$. A symmetric argument shows \$|D^{3}f(x)\[u, u, u]| \leq M|u|^{3}\_{x}\$, implying that \$f\$ is self-concordant with parameter \$M/2\$. Moreover, Theorem 5.1.2 of \cite{convexlecture} shows that the strong self-concordance parameter \$M\$ is affine invariant: for any non-singular matrix \$A \in \mathbb{R}^{d\times d}\$, the function \$\phi(x) = f(Ax)\$ remains \$M\$-strongly self-concordant.

%We now explain why our assumptions are strictly weaker than the more common conditions of strong convexity, Lipschitz gradient, and Lipschitz Hessian. Prior work (e.g., Example 4.1 in\~\cite{rodomanov2020greedy}) shows that if a function is strongly convex and has a Hessian that is smooth with respect to a matrix \$B \succeq 0\$, then it is also strongly self-concordant. However, the converse does not hold: strong self-concordance does not imply strong convexity, gradient smoothness, or Lipschitz Hessian continuity.

As a concrete example, we can consider the log-sum-exp function formally defined as $f(x) = \log{(\sum_{i = 1}^{n}\exp({c_i^\top x - b_i}))} + \sum_{i = 1}^{n}(c_i^\top x)^2$, where $\{c_i\}_{i = 1}^{n} \in \mathbb{R}^d$ and $\{b_i\}_{i = 1}^{n} \in \mathbb{R}$. This function is not strongly convex with respect to the identity matrix $I$, due to the absence of explicit $\ell_2$ regularization. However, it can be shown to be strongly convex and have Lipschitz Hessian with respect to the matrix $B = \sum_{i = 1}^{n}c_ic_i^\top$ (Note that this matrix could be possibly singular). As a result, it is strongly self-concordant but not strongly convex in the standard sense; check Appendix~\ref{proof_of_concordance}. Other examples include the hard cubic function and the logistic regression objective discussed in Section~\ref{sec:experiments}. Another illustrative case is the log-barrier function $f(x) = -\log(1 - x^2)$, which is strongly self-concordant with $M = 4$ for $|x - y| \leq 1/2$, yet its gradient and Hessian are not Lipschitz continuous. Full detailed discussion for these examples is provided in Appendix~\ref{proof_of_concordance}.

%These examples underscore that our assumptions—strict convexity and strong self-concordance—are  weaker than those made in previous works~\cite{qiujiang2020quasinewton, jin2024non, qiujiang2024quasinewton}, which require strong convexity, Lipschitz gradient, and Lipschitz Hessian.

%One key observation is that our assumptions of strict convexity and strong self-concordance are weaker than the assumptions of strong convexity, gradient smoothness and Hessian Lipschitz continuity from previous works \cite{rodomanov2020rates,rodomanov2020ratesnew,ye2023towards,qiujiang2020quasinewton,krutikov2023convergence,antonglobal,jin2024non,qiujiang2024quasinewton}. This is because from Example 4.1 in \cite{rodomanov2020greedy}, it is established that if a function is strongly convex and its Hessian is smooth with respect to some matrix $B \succeq 0$, then the function is strongly self-concordant. However, the reverse does not hold: if a function is strongly self-concordant, there is no guarantee that it is strongly convex, its gradient is smooth or its Hessian is Lipschitz continuous. One example is the log-sum-exp function, which is formally defined as: $f(x) = \log{(\sum_{i = 1}^{n}e^{c_i^\top x - b_i})} + \sum_{i = 1}^{n}(c_i^\top x)^2$, where $\{c_i\}_{i = 1}^{n} \in \mathbf{R}^d$ and $\{b_i\}_{i = 1}^{n} \in \mathbf{R}$. This function is not strongly convex with respect to the standard linear system defined by the identity matrix $I$ since there are no $l_2$ regularization. However, it can be shown that this function is strongly convex and its Hessian is Lipschitz smooth with respect to the operator matrix $B = \sum_{i = 1}^{n}c_ic_i^\top$. Hence, from Example 4.1 in \cite{rodomanov2020greedy}, the log-sum-exp function is strongly self-concordant but not strongly convex. Other such examples include the hard cubic function and the logistic regression function defined in section~\ref{sec:experiments}. Another function is the log barrier function $f(x)=- log(1-x^2)$, which is strongly self-concordant with $M=4$ when $\|x-y\|\leq 1/2$, while both of its gradient and its Hessian are not Lipschitz smooth. Please check all detailed proof in Appendix~\ref{proof_of_concordance}. Consequently, our assumption is significantly weaker than those made in previous works \cite{qiujiang2020quasinewton, jin2024non, qiujiang2024quasinewton}, which require the objective function to be strongly convex, its gradient to be Lipschitz smooth, and its Hessian to be Lipschitz smooth.

%\begin{remark}

%\textcolor{red}{From Example 4.1 in \cite{rodomanov2020greedy}, it is established that if a function is strongly convex and its Hessian is smooth with respect to some matrix $B \succeq 0$, then the function is strongly self-concordant. However, the reverse does not hold: if a function is strongly self-concordant, there is no guarantee that it is strongly convex for some constant \(\mu > 0\). One such example is the log-sum-exp function, which is formally defined as: $f(x) = \log{(\sum_{i = 1}^{n}e^{c_i^\top x - b_i})} + \sum_{i = 1}^{n}(c_i^\top x)^2$, where $\{c_i\}_{i = 1}^{n} \in \mathbf{R}^d$ and $\{b_i\}_{i = 1}^{n} \in \mathbf{R}$. This function is not strongly convex with respect to the standard linear system defined by the identity matrix $I$ since there are no $l_2$ regularization. However, it can be shown that this function is strongly convex and its Hessian is Lipschitz smooth with respect to the operator matrix $B = \sum_{i = 1}^{n}c_ic_i^\top$. (Please check proof details in the numerical experiment section of \cite{rodomanov2020greedy}) Hence, by Example 4.1 in \cite{rodomanov2020greedy}, the log-sum-exp function is strongly self-concordant but not strongly convex. Consequently, our assumption is significantly weaker than those made in previous works \cite{qiujiang2020quasinewton, jin2024non, qiujiang2024quasinewton}, which require the objective function to be strongly convex, its gradient to be Lipschitz smooth, and its Hessian to be Lipschitz smooth.}

%\end{remark}

%\begin{remark}
%\textcolor{red}{This assumption of strong self-concordance in Assumption~\ref{ass_self_concodant} cannot be relaxed to the semi-strong self-concordance condition from \cite{dampedNewton} and standard self-concordance from \cite{convexlecture}. The strong self-concordance condition implies the inequality of $C_1 \nabla^2{f(y)} \preceq \nabla^2{f(x)} \preceq C_2 \nabla^2{f(y)}$ with some $C_1, C_2 > 0$ for any two points $x, y \in \mathbb{R}^d$, which plays a critical role in the theoretical analysis for the BFGS method. However, we can not obtain similar inequality by semi-strong self-concordance condition without the strong convexity because semi-strong self-concordance condition only assumes the upper bound of the operator norm of the difference of two Hessian matrices on two different points. We cannot relax this strong self-concordance assumption to the standard self-concordance neither, which is applied for the convergence analysis of Newton's method. This is because the standard  self-concordance assumption only holds locally, while the strong self-concordance assumption holds globally. The convergence analysis of BFGS method relies heavily on this global property of the strong self-concordance condition.}   
%\end{remark}

\subsection{Definitions}\label{sec:intermediate}

Next, we state our definitions and notations. For any $A \in \mathbb{S}^d_{++}$, we define  $\Psi(A)$ as
\begin{equation}\label{potential_function}
    \Psi(A) := \Tr(A) - d - \log{\Det(A)}.
\end{equation}
This function characterizes the distance between matrix $A$ and the identity matrix $I$. Note that $\Psi(A) \geq 0$ for any $A \in \mathbb{S}^d_{++}$ and $\Psi(A) = 0$ if and only if $A = I$.

A common technique in the analysis of quasi-Newton methods involves the use of a reweighting matrix; see, e.g.,  \cite{QN_tool}. We also use this approach in our analysis. Specifically, given any weight matrix $P \in \mathbb{S}^d_{++}$, we define the weighted versions of the vectors $g_t$, $s_t$, $y_t$, $d_t$ and the matrix $B_t$ as 
\begin{equation}\label{weighted_vector}
\hat{g}_t\! := \!P^{-\frac{1}{2}}g_t,\quad  
\hat{s}_t \!:= \! P^{\frac{1}{2}} s_t, \quad 
\hat{y}_t \! := \! P^{-\frac{1}{2}}y_t,\quad  
\hat{d}_t \! := \! P^{\frac{1}{2}} d_t,\quad 
\hat{B}_t := P^{-\frac{1}{2}}B_t P^{-\frac{1}{2}}.
\end{equation}
The weight matrix $P$ plays fundamental role in our proof and the global linear and superlinear convergence rates are based on \textit{different choices of $P$}. Note that the update rule for the weighted version of Hessian approximation matrices $\hat{B}_t$ is similar to the update rule of the unweighted $B_t$, i.e., $\hat{B}_{t+ 1} = \hat{B}_t - \frac{\hat{B}_t \hat{s}_t \hat{s}_t^\top \hat{B}_t}{\hat{s}_t^\top \hat{B}_t \hat{s}_t} + \frac{\hat{y}_t \hat{y}_t^\top}{\hat{s}_t^\top \hat{y}_t}$. We next introduce a common function in self-concordant analysis:
\begin{equation}\label{definition_omega}
    \omega(x) := x - \log{(x + 1)}.
\end{equation}
As shown in Lemma~\ref{lemma_omega}, $\omega(x)$ is strictly increasing for $x > 0$. Hence, we can define its inverse function $\omega^{-1}(.)$  such that $\omega^{-1}(\omega(x)) = x$  for $x > 0$. It can be verified that \(\omega^{-1}(x)\) is also strictly increasing for \(x > 0\). Further, since \(\omega(x)\) is a convex function, \(\omega^{-1}(x)\) is concave. We use $\omega^{-1}$ to measure suboptimality of the iterates $\{x_t\}_{t = 0}^{+\infty}$ and define the sequences $\{C_t\}_{t = 0}^{+\infty}$ and $\{D_t\}_{t = 0}^{+\infty}$ as
\begin{equation}\label{suboptimality}
    C_t := f(x_t) - f(x_*),\qquad 
    D_t := 2\omega^{-1}\left({M^2}C_t/4\right),
\end{equation} 
Indeed, both of the above sequences are always non-negative. 

\begin{remark}
  The expression $\omega^{-1}(.)$ frequently appears in our complexity bounds. To better understand this function and its approximation, as shown in Lemma \ref{lemma_omega}, we can use the approximation $\omega^{-1}(a) \approx (a + \sqrt{2a})$. Consequently, if $a < 1$, $\omega^{-1}(a) = \mathcal{O}(\sqrt{a})$, and if $a > 1$, $\omega^{-1}(a) = \mathcal{O}(a)$.
\end{remark}

With these preliminaries, the next two sections prove global linear and superlinear convergence rates of BFGS for strictly convex, strongly self-concordant functions—rates that remain invariant under linear transformations, consistent with BFGS's affine invariance.

\section{Global Linear Convergence Rates}\label{sec:linear}

In this section, we present the global linear convergence results of BFGS when the step size is selected based on the weak Wolfe conditions introduced in \eqref{sufficient_decrease} and \eqref{curvature_condition}. Before we begin, we need to define the following weighted versions of the initial Hessian approximation matrix \( B_0 \):
%%
\begin{equation}\label{def_B_bar}
    \bar{B_0} = \frac{\nabla^2 f(x_*)^{-\frac{1}{2}} B_0 \nabla^2 f(x_*)^{-\frac{1}{2}}}{1 + D_0}, \qquad 
    \tilde{B}_0 = \nabla^2 f(x_*)^{-\frac{1}{2}} B_0 \nabla^2 f(x_*)^{-\frac{1}{2}}.
\end{equation} 
%%
These two weighted versions of \( B_0 \) correspond to the weight matrices \( P = (1 + D_0)\nabla^2 f(x_*) \) and \( P = \nabla^2 f(x_*) \), respectively. They play a key role in the non-asymptotic analysis of BFGS for self-concordant functions. Next, we present our first global explicit linear convergence rate of BFGS for any initial point $x_0$ and any initial Hessian approximation matrix $B_0 \in \mathbb{S}^d_{++}$. %The proof is presented in Appendix~\ref{proof_thm_first_linear}.

\begin{theorem}\label{thm:first_linear_rate}
Suppose Assumption~\ref{ass_self_concodant} holds. Let $\{x_t\}_{t\geq 0}$ be the iterates generated by BFGS, where the step size satisfies the Armijo-Wolfe conditions in \eqref{sufficient_decrease} and \eqref{curvature_condition}. Recall $\Psi(\cdot)$ in \eqref{potential_function}, $D_0$ in \eqref{suboptimality} and $\bar{B}_0$ in \eqref{def_B_bar}. For any initial point $x_0 \in \mathbb{R}^{d}$ and any initial Hessian approximation $B_0 \in \mathbb{S}^d_{++}$, we have
\begin{equation}\label{eq:first_linear_rate}
    \frac{f(x_t) - f(x_*)}{f(x_0) - f(x_*)} \leq \left(1 - \frac{\alpha(1 - \beta)e^{-\frac{\Psi(\bar{B}_0)}{t}}}{(1 + D_0)^2}\right)^{t}.
\end{equation}
Moreover, when $t \geq \Psi(\bar{B}_0)$, we obtain that
\begin{equation}\label{eq:first_linear_rate_2}
    \frac{f(x_t) - f(x_*)}{f(x_0) - f(x_*)} \leq \left(1 - \frac{\alpha(1 - \beta)}{3(1 + D_0)^2}\right)^{t}.
\end{equation}
\end{theorem}

Theorem~\ref{thm:first_linear_rate} states that BFGS converges globally at a linear rate, influenced by the line search parameters (as expected), the term $\Psi(\bar{B}_0)$, which quantifies the discrepancy between the initial Hessian approximation and the optimal one, and $D_0$, which depends on the suboptimality of the initial function value and the strongly self-concordance parameter.
To further simplify the expression, as shown in the second result, when $t \geq \Psi(\bar{B}_0)$, the linear convergence rate can be further simplified as  $\mathcal{O}(1 - 1/(1 + D_0)^2)$. Hence, $D_0=2\omega^{-1}({M^2}(f(x_0)-f(x_*)/4)$ indicates the rate. %This bound implies that to achieve an $\epsilon$ accurate solution, BFGS requires at most $\mathcal{O}((1+D_0)^2\log(1/\epsilon) \ +\ \Psi(\bar{B}_0) )$ which can be simplified as $\mathcal{O}((1+2\omega^{-1}(\frac{M^2}{4}(f(x_0)-f(x_*)))^2\log(1/\epsilon) \ +\ \Psi(\bar{B}_0))$. 

Two remarks follow the above result. First, our global linear convergence rate does not require assuming strong convexity or gradient Lipschitz-ness. Second, the linear convergence rate is affine invariant across different linear systems, consistent with the affine invariance property of BFGS. 

We emphasize that the proof of Theorem~\ref{thm:first_linear_rate} for showing global linear convergence rate is fundamentally different from the analyses in prior work. Specifically, the results in \cite{jin2024non, qiujiang2020quasinewton, qiujiang2024quasinewton} heavily depend on the strong convexity and gradient Lipschitz-ness assumptions to showcase a linear convergence rate: they use the Lipschitz continuity of the gradient to upper bound  $\|y_t\|^2/s_t^\top y_t$ by $L$, and use $\mu$-strong convexity to establish the following lower bound $\|g_t\|^2/(f(x_t) - f(x_*)) \geq 2\mu$. These bounds are key to establishing the global linear rate of BFGS in prior work. In our setting such bounds do not hold and we do not have a universal upper bound on $\|y_t\|^2/s_t^\top y_t$ and a lower bound on $\|g_t\|^2/(f(x_t) - f(x_*))$. Instead, for the first bound, we transfer the inequality to the norm induced by the  weight matrix $P = (1 + D_0)\nabla^2{f(x_*)}$ and show under this norm and strong self-concordance assumption we have $\|\hat{y}_t\|^2/\hat{s}_t^\top \hat{y}_t\leq 1$. For the lower bound on $\|g_t\|^2/(f(x_t) - f(x_*))$, instead of a uniform lower bound, we show that it can be bounded below by $1/(1 + D_t)$, which is dependent on $x_t$, but we show that even this time-dependent lower bound is sufficient to establish a linear convergence rate for BFGS. For more details check the proofs of  Lemma~\ref{lemma_g} and Section~\ref{proof_thm_first_linear} in the Appendix. 

The linear convergence result depends on $\Psi(\bar{B}_0)$, and hence the choice of $B_0$ affects the convergence rate. In practice, it is often a scaled identity and a common choice is $B_0 = cI$, where $c = ({s^\top y})/{\|s\|^2}$, with $s = x_2 - x_1$, $y = \nabla f(x_2) - \nabla f(x_1)$, and $x_1, x_2 $ as two randomly selected vectors. In the next corollary, we present our global linear rate when $B_0 = a I$ where $a > 0$ is an arbitrary constant.

\begin{corollary}\label{coro:first_linear_rate}
Suppose Assumptions~\ref{ass_self_concodant} holds, $\{x_t\}_{t\geq 0}$  are generated by BFGS with step size satisfying the Armijo-Wolfe conditions in \eqref{sufficient_decrease} and \eqref{curvature_condition}, and $x_0 \in \mathbb{R}^{d}$ is an arbitrary initial point. If the initial Hessian approximation matrix is set as $B_0 = a I$ for any $a > 0$, then we have that
\begin{equation}\label{eq:first_linear_rate_a}
    \frac{f(x_t) - f(x_*)}{f(x_0) - f(x_*)} \leq \left(1 - \frac{\alpha(1 - \beta)e^{-\frac{\Delta_1}{t}}}{(1 + D_0)^2}\right)^{t},
\end{equation}
where $\Delta_1 := \Psi(\frac{a\nabla^{2}{f(x_*)^{-1}}}{1 + D_0})$ can be written as 
\begin{equation}\label{Delta_1}
\Delta_1 = \Tr\left[\frac{a\nabla^{2}{f(x_*)^{-1}}}{1 + D_0}\right] - d - \log{\Det\!\left[\frac{a\nabla^{2}{f(x_*)^{-1}}}{1 + D_0}\right]}.
\end{equation}
Moreover, when $t \geq \Delta_1$, we obtain that
\begin{equation}\label{eq:first_linear_rate_a_2}
    \frac{f(x_t) - f(x_*)}{f(x_0) - f(x_*)} \leq \left(1 - \frac{\alpha(1 - \beta)}{3(1 + D_0)^2}\right)^{t}.
\end{equation}
\end{corollary}

Note that the proof of this corollary simply follows by setting $B_0 = a I$ in Theorem~\ref{thm:first_linear_rate}. The above result shows that by selecting $B_0 = a I$, the linear convergence rates of the BFGS method is totally determined by the initial suboptimality $D_0$ and the trace and determinant of the inverse matrix of the Hessian at $x_*$, which are also consistent with the affine invariance property of BFGS.

Next, we proceed to present an improved version of the result in Theorem \ref{thm:first_linear_rate}, showing that after a sufficient number of iterations, the linear rate of BFGS becomes independent of  $D_0$ and $B_0$.

\begin{theorem}\label{thm:second_linear_rate}
Suppose Assumptions~\ref{ass_self_concodant} holds, and let $\{x_t\}_{t\geq 0}$ be the iterates generated by the BFGS method with the Armijo-Wolfe line search in \eqref{sufficient_decrease} and \eqref{curvature_condition}. Recall the definition of $\Psi(\cdot)$ in \eqref{potential_function}, $D_0$ in \eqref{suboptimality} and $\bar{B}_0$, $\tilde{B}_0$ in \eqref{def_B_bar}. Then, for any initial point $x_0 \in \mathbb{R}^{d}$ and any initial Hessian approximation matrix $B_0 \in \mathbb{S}^d_{++}$, when $
    t \geq \Psi(\Tilde{B}_{0}) + 3D_0(\Psi(\bar{B}_{0}) + \frac{3(1 + D_0)^2}{\alpha(1 - \beta)})$,
we have 
\begin{equation}\label{eq:second_linear_rate}
    \frac{f(x_t) - f(x_*)}{f(x_0) - f(x_*)} \leq \left(1 - \frac{2\alpha(1 - \beta)}{3} \right)^{t}.
\end{equation}
\end{theorem}

This theorem demonstrates that when the number of iterations is larger than $\Psi(\Tilde{B}_{0}) + 3D_0(\Psi(\bar{B}_{0}) + \frac{3(1 + D_0)^2}{\alpha(1 - \beta)})$, BFGS with stepsize satisfying the Armijo-Wolfe conditions achieves an explicit linear convergence rate that is independent of the initial suboptimality $D_0$ and only determined by the line search parameters $\alpha$ and $\beta$ defined in \eqref{sufficient_decrease} and \eqref{curvature_condition}. That said, the point that transition to this fast rate happens still depends on the choice of $x_0$ and $B_0$, as stated in Theorem~\ref{thm:second_linear_rate}.
Similar to Corollary~\ref{coro:first_linear_rate}, next we present the special case of Theorem~\ref{thm:second_linear_rate} of $B_0 = a I$ for any $a > 0$.

\begin{corollary}\label{coro:second_linear_rate}
Suppose Assumptions~\ref{ass_self_concodant} holds, $\{x_t\}_{t\geq 0}$  are generated by BFGS with step size satisfying the Armijo-Wolfe conditions in \eqref{sufficient_decrease} and \eqref{curvature_condition}, and $x_0 \in \mathbb{R}^{d}$ is an arbitrary initial point. If the initial Hessian approximation matrix is set as $B_0 = a I$ for any $a > 0$, then the following rate holds
\begin{equation}\label{eq:second_linear_rate_a}
    \frac{f(x_t) - f(x_*)}{f(x_0) - f(x_*)} \leq \left(1 - \frac{2\alpha(1 - \beta)}{3} \right)^{t},
\end{equation}
for all iterates satisfying $
    t \geq \Delta_2 + 3D_0\left(\Delta_1 + \frac{3(1 + D_0)^2}{\alpha(1 - \beta)}\right)$,
where $\Delta_1$ is defined in \eqref{Delta_1} and 
\begin{equation}\label{Delta_2}
\Delta_2= \Tr(a\nabla^{2}{f(x_*)^{-1}}) - d - \log{\Det(a\nabla^{2}{f(x_*)^{-1}})}.
\end{equation}
\end{corollary}

Note that both $\Delta_1$ and $\Delta_2$ are determined by the Hessian at the optimal solution $x_*$, while $\Delta_1$ also depends on the initial suboptimality error through $D_0$. In general, we do expect the convergence rates of BFGS to depend on the distance between $x_0$ and $x_*$, which is characterized by $D_0$ defined in \eqref{suboptimality} as well as the distance between the initial Hessian approximation matrix $B_0$ and the exact Hessian at optimal solution $x_*$, which is characterized by $\Delta_1$ and $\Delta_2$ when $B_0 = \alpha I$.

%\textcolor{blue}{here you need a remark to explain how the results for linear convergence in this section were different from those in prior analysis of BFGS and how you had to establish new bounds without those assumptions. you basically need to mention that global linear rate of prior work heavily depends on grad lip, and without that something will be missing. then explain here you did what to overcome this issue and still managed to get a global linear rate without grad lip and strong convexity}



\section{Global Superlinear Convergence Rates}\label{sec:superlinear}

Building on the established linear convergence results, we next establish our global superlinear convergence rate of BFGS. A key point in our analysis is that to reach the superlinear convergence stage, the unit step size must be chosen after some iterations. This is a necessary condition, as noted in several prior works \cite{powell1971convergence,Powell,byrd1987global,QN_tool}. The fundamental methodology is to first establish the sufficient conditions of when the unit step size can be selected, i.e., when $\eta_t = 1$ satisfies the conditions in \eqref{sufficient_decrease} and \eqref{curvature_condition}. Then, based on these conditions, we can prove that after some specific iterations $t_0$, the unit step size $\eta_t = 1$ is admissible for the inexact line search scheme except for a finite number of iterations, which leads to the final proof of the global non-asymptotic superlinear convergence rate.

Next, we proceed to establish under what conditions $\eta=1$ is admissible. 
First, define $\rho_t$ as
\begin{equation}\label{definition_rho}
    \rho_t := \frac{-g_t^\top d_t}{\|\tilde{d}_t\|^2}, \quad \tilde{d}_t := \nabla^2{f(x_*)}^{\frac{1}{2}}d_t, \quad \forall t \geq 0.
\end{equation}
In the following lemma, we demonstrate that when $C_t = f(x_t) - f(x_*)$ is small enough and  $\rho_t$ is close enough to $1$, the unit step size $\eta_t=1$ is admissible and meets the Armijo-Wolfe conditions. %The proof is presented in Appendix~\ref{proof_lem:unit_step_size}.

\begin{lemma}\label{lem:unit_step_size}
    Suppose Assumption~\ref{ass_self_concodant} holds and define
    \begin{equation}\label{definition_delta_1_2_3}
    \begin{split}
        &\delta_1 := \min \bigg\{\frac{1}{16},  \frac{4}{M^2}\omega\left(\frac{1}{32}\right),  \frac{4}{M^2}\omega\left(\frac{\sqrt{2(1 - \alpha)} - 1}{2}\right)\!,  \frac{4}{M^2}\omega\left(\frac{1}{2}\left(\frac{1}{\sqrt{1 - \beta}} -\! 1\right)\right) \bigg\}, \\
        & \delta_2 := \max\left\{\frac{15}{16}, \ \frac{1}{\sqrt{2(1 - \alpha)}}\right\}, \ \delta_3 := \frac{1}{\sqrt{1 - \beta}},
    \end{split}
    \end{equation}
    which satisfy $0 < \delta_1 < \delta_2 < 1< \delta_3$.  If $C_t \leq \delta_1$ and $\delta_2 \leq \rho_t \leq \delta_3$, then $\eta_t = 1$ satisfies \eqref{sufficient_decrease} and \eqref{curvature_condition}.
\end{lemma}


First, we highlight the key difference between Lemma~\ref{lem:unit_step_size} and prior results in \cite{qiujiang2020quasinewton, qiujiang2024quasinewton, jin2024non}. The proof of Lemma~\ref{lem:unit_step_size} hinges on ensuring $f(x_t + d_t) \leq f(x_t)$, i.e., that a unit step yields a decrease in function value. Under Lipschitz continuity of the Hessian with constant $K$, the error of approximating $f(y)$ by its second-order Taylor expansion at $x$ is bounded by $\frac{K}{6}\|y - x\|^3$. Without this assumption, and under $M$-strongly self-concordant assumption, we instead use the bound
$f(y) \leq f(x) + g(x)^\top(y - x) + \frac{4}{M^2}\omega_{*}\left(\frac{M}{2}\|y - x\|_x\right)$
for $\|y - x\|_x < \frac{2}{M}$, where $\omega_{*}(x) = - x - \log(1 - x)$ is defined for $x < 1$. As a result, the error is no longer cubic in $\|y - x\|$, making it more challenging to ensure a function decrease. Nevertheless, we can still guarantee this property, with the main difference being that the error bound $\delta_1$ now depends on $\omega(x)$ defined above. See Lemma~\ref{lemma_decrease} and Section~\ref{proof_lem:unit_step_size} for details.


The result in Lemma~\ref{lem:unit_step_size} shows that when $C_t\leq\delta_1$ and $\rho_t \in [\delta_2, \delta_3]$, we can choose the step size $\eta_t = 1$ at iteration $t$ of BFGS, as it  satisfies the weak Wolfe conditions. 
Moreover, from the global non-asymptotic linear convergence rates of the last section, we can specify the $t_0$ such that for any $t \geq t_0$, the first condition $C_t \leq \delta_1$ always holds. Moreover, we can demonstrate that the second condition on $\rho_t$ is violated only for a finite number of iterations, i.e., the set of the indices that $\rho_t \notin [\delta_2, \delta_3]$ can be upper bounded by some constants. We formally present these results in the following lemma and the proofs are available in Appendix~\ref{proof_lem:bad_set}. 

\begin{lemma}\label{lem:bad_set}
    Suppose Assumptions~\ref{ass_self_concodant} holds and  $\{x_t\}_{t \geq 0}$ are generated by BFGS with step size satisfying the Armijo-Wolfe conditions in \eqref{sufficient_decrease}-\eqref{curvature_condition}. Recall the definition of $C_t$ in \eqref{suboptimality}, $D_t$ in \eqref{suboptimality}, $\Psi(\cdot)$ in \eqref{potential_function}, $\{\delta_i\}_{i = 1}^{3}$ in \eqref{definition_delta_1_2_3}, and $\bar{B}_0$, $\tilde{B}_0$ in \eqref{def_B_bar}. We have $C_t \leq \delta_1$ when $t \geq t_0$, where $t_0$ is defined as
    \begin{equation}\label{definition_t_0}
        t_0 := \max\left\{\Psi(\bar{B}_0), \ \frac{3(1 + D_0)^2}{\alpha(1 - \beta)}\log{\frac{C_0}{\delta_1}}\right\}.
    \end{equation}
    Moreover, the size of the set $I = \{t_0 \leq i \leq t - 1:\; \rho_t \notin [\delta_2,\delta_3]\}$ is at most 
    \begin{equation}\label{definition_delta_4}
        |I| \leq \delta_4\left(\Psi(\tilde{B}_{0}) + 2D_0\left(\Psi(\bar{B}_{0}) + \frac{3(1 + D_0)^2}{\alpha(1 - \beta)}\right)\right),\quad \!\!\text{where}\ \ \delta_4 := \frac{1}{\min\{\omega(\delta_2 - 1), \omega(\delta_3 - 1)\}}.
    \end{equation}
\end{lemma}

The above lemma specifies the time instance $t_0$ for which $C_t \leq \delta_1$ is satisfied for any $t \geq t_0$ and for only a finite number of indices, the condition $\rho_t \in [\delta_2, \delta_3]$ does not hold. In practice, we always start with the unit step size when we implement the inexact line search scheme at iteration $t$ to check if $\eta_t = 1$ satisfies the Armijo-Wolfe conditions in \eqref{sufficient_decrease} and \eqref{curvature_condition}. Hence, when $t \geq t_0$, only for a finite number of iterations that $\rho_t \notin [\delta_2, \delta_3]$, the unit step size is not selected. With all these points, we present the global superlinear convergence rate of BFGS for self-concordant functions. %The proof is in Appendix~\ref{proof_thm:superlinear_rate_2}.

\begin{theorem}\label{thm:superlinear_rate_2}
    Suppose Assumptions~\ref{ass_self_concodant} holds and the iterates $\{x_t\}_{t \geq 0}$ are generated by BFGS with step size satisfying the Armijo-Wolfe conditions in \eqref{sufficient_decrease} and \eqref{curvature_condition}. Recall the definition of $D_t$ in \eqref{suboptimality}, $\Psi(\cdot)$ in \eqref{potential_function}, $\bar{B}_0$, $\tilde{B}_0$ in \eqref{def_B_bar}, and $\{\delta_i\}_{i = 1}^{4}$ in \eqref{definition_delta_1_2_3}, \eqref{definition_delta_4}. Then, for any initial point $x_0 \in \mathbb{R}^{d}$ and any initial Hessian approximation matrix $B_0 \in \mathbb{S}^d_{++}$, the following global superlinear result holds:
    \begin{equation*}
    \begin{split}
        \frac{f(x_t) - f(x_*)}{f(x_0) - f(x_*)} \leq  \left(\frac{\delta_6 t_0 + \delta_7 \Psi(\tilde{B}_0) + \delta_8 D_0 (\Psi(\bar{B}_{0}) + \frac{3(1 + D_0)^2}{\alpha(1 - \beta)})}{t}\right)^{t},
    \end{split}
    \end{equation*}
    where $t_0$ is defined in \eqref{definition_t_0}, $\{\delta_i\}_{i = 5}^{8}$ defined below only depend on line search parameters $\alpha$ and $\beta$,
    \begin{equation}\label{def_delta_5678}
    \begin{split}
        & \delta_5 := \max\left\{\frac{2 + ({2}/{\delta_2})}{2\delta_2 - 17/16}, \ \frac{4\delta_3}{2\delta_2 - 17/16}\right\}, \qquad \delta_6 := \log{\frac{1}{2\alpha(1 - \beta)}},\\
        &  \qquad \delta_7 := 1 + \delta_4\delta_6 + \delta_5, \qquad \delta_8 := 2 + 2\delta_4\delta_6 + 2\delta_5 + \frac{2\delta_2 - 1/16 - \log{\delta_2}}{2\delta_2 - 17/16}.
    \end{split}
    \end{equation}
\end{theorem}

Theorem~\ref{thm:superlinear_rate_2} shows that the superlinear convergence rate of BFGS for a self-concordant function is of the form $({C}/{t})^{t}$ for some constant $C > 0$. Notice that from the definition of $t_0$ in \eqref{definition_t_0}, we know that $t_0 = \mathcal{O}(\Psi(\bar{B}_{0}) + (1 + D_0)^2\log{D_0})$. Hence, the superlinear convergence rate is of the order
$
    \mathcal{O}((\frac{\Psi(\tilde{B}_0) + D_0 (\Psi(\bar{B}_{0}) + (1 + D_0)^2)}{t})^t)$,
and we reach the superlinear convergence stage when $
    t \geq \Omega(\Psi(\tilde{B}_0) + D_0 (\Psi(\bar{B}_{0}) + (1 + D_0)^2))$,
which depends on the initial suboptimality $D_0$ and the initial Hessian approximation matrix $B_0$. To our knowledge, this is the first non-asymptotic global superlinear convergence rate of a quasi-Newton method without the assumption of strong convexity. Moreover, the superlinear rate in Theorem~\ref{thm:superlinear_rate_2} is independent of the linear system chosen for the variables, and, hence, it is consistent with the affine invariance property of BFGS. Next, we present the superlinear convergence rate of BFGS for the special case of $B_0 = a I$, where $a > 0$.

\begin{corollary}\label{coro:superlinear_rate}
Suppose Assumptions~\ref{ass_self_concodant} holds, $\{x_t\}_{t\geq 0}$  are generated by BFGS with step size satisfying the Armijo-Wolfe conditions in \eqref{sufficient_decrease} and \eqref{curvature_condition}, and $x_0 \in \mathbb{R}^{d}$ is an arbitrary initial point. If the initial Hessian approximation matrix is $B_0 = a I$ where $a > 0$, the following result holds:
\begin{equation*}
 \frac{f(x_t) \!-\! f(x_*)}{f(x_0)\! -\! f(x_*)} 
     \!\leq \!\left(\!\frac{\delta_6 t_0 \!+\! \delta_7 \Delta_2 \!+\! \delta_8 D_0 (\Delta_1 \!+\! \frac{3(1 + D_0)^2}{\alpha(1 - \beta)})}{t}\!\right)^{\!\!t}\!,
\end{equation*}
where $t_0$ is defined in \eqref{definition_t_0}, $\{\delta_i\}_{i = 5}^{8}$ are defined in \eqref{def_delta_5678} and $\Delta_1$, $\Delta_2$ are defined in \eqref{Delta_1}, \eqref{Delta_2}.
\end{corollary}

%\textcolor{blue}{again here you need a remark to explain how the results for superlinear convergence in this section were different from those in prior analysis of BFGS and how you had to establish new bounds without those assumptions. you basically need to mention that global superlinear rate of prior work heavily depends on grad and hessian lip and strong convexity, and without those the arguments would break. then explain here you did what to overcome this issue and still managed to get a global superlinear rate without those assumptions}

%\textcolor{red}{Similar to the linear rate analysis, the global superlinear convergence analysis in this section is totally different from the previous works \cite{jin2024non, qiujiang2020quasinewton, qiujiang2024quasinewton} requiring strong convexity, gradient smoothness and Hessian Lipschitz continuity. These arguments imply that $C_1 \nabla^2{f(y)} \preceq \nabla^2{f(x)} \preceq C_2 \nabla^2{f(y)}$ with some $C_1, C_2 > 0$ for any two points $x, y \in \mathbb{R}^d$, which plays a fundamental role in the final non-asymptotic superlinear convergence analysis. However, without the strong convexity, we can not establish such relationships between Hessian matrices at different points because Hessian smoothness condition only assumes the upper bound of the operator norm of the difference between two Hessian matrices on two different points. We applied the strong self-concordance condition to overcome these difficulties by proving the similar inequalities of $C_1 \nabla^2{f(y)} \preceq \nabla^2{f(x)} \preceq C_2 \nabla^2{f(y)}$ in Lemma~\ref{lemma_Hessian} from Appendix. We managed to demonstrate a global superlinear rate without those assumptions from \cite{jin2024non, qiujiang2020quasinewton, qiujiang2024quasinewton}.}

%Similar to our linear rate analysis, the  superlinear convergence analysis presented here differs fundamentally from prior work \cite{jin2024non, qiujiang2020quasinewton, qiujiang2024quasinewton}, which rely on strong convexity as well as gradient and Hessian Lipschitzness. In the superlinear convergence analysis, the key quantity to control is
%$\frac{\|\tilde{g}_t\|^2}{2(f(x_t) - f(x_*))}$ which must converge to $1$.  In standard analyses, this follows from the smoothness of the Hessian and strong convexity, as we can lower bound the denominator by  $f(x_t) - f(x_*) \geq \frac{1}{2}\mu\|x_t - x_*\|^2$ and upper bound the numerator by ?????. 

%Without strong convexity and Hessian continuity, we need to apply the strong self-concordance from Assumption~\ref{ass_self_concodant} with inequality $f(y) \geq f(x) + g(x)^\top(y - x) + \frac{4}{M^2}\omega(\frac{M}{2}\|y - x\|_x)$ where $\omega$ defined in \eqref{definition_omega}. Then, we can bound the term as $\frac{1}{(1 + D_t)^2} \leq \frac{\|\tilde{g}_t\|^2}{2(f(x_t) - f(x_*))} \leq (1 + D_t)^2$, which implies that $\lim_{t \to \infty}\frac{\|\tilde{g}_t\|^2}{2(f(x_t) - f(x_*))} = 1$ since $D_t$ converges to $0$. These bounds play critical roles in the final proof of the non-asymptotic superlinear convergence rate. Please check details in Lemma~\ref{lemma_Hessian} and Lemma~\ref{lemma_eigenvalue} from Appendix. 

\section{Complexity Analysis}\label{sec:complexity}

\noindent\textbf{Iteration Complexity.} Using Theorems~\ref{thm:first_linear_rate}, \ref{thm:second_linear_rate}, and \ref{thm:superlinear_rate_2}, we characterize the global iteration complexity of BFGS with inexact line search on self-concordant functions.
 These three results provide upper bounds, and the smallest of these bounds determines the complexity of BFGS. The smallest bound depends on the required accuracy relative to the problem and algorithm parameters. Specifically, for any initial point $x_0 \in \mathbb{R}^d$ and initial Hessian approximation matrix $B_0 \in \mathbb{S}^d_{++}$, to achieve a function value accuracy of $\epsilon > 0$, i.e., $f(x_T) - f(x_*) \leq \epsilon$, the number of iterations required, as per Theorem~\ref{thm:first_linear_rate}, is at most $
    T_1=\mathcal{O}\left(\Psi(\bar{B_0}) + (1 + D_0)^2\log{\frac{1}{\epsilon}}\right) $. 
The result in Theorem \ref{thm:superlinear_rate_2} eliminates the multiplicative factor in the $\log(1/\epsilon)$ term but requires a possibly larger additive constant, resulting in a complexity of $
T_2=\mathcal{O}(\Psi(\tilde{B_0}) + (\Psi(\bar{B_0}) + (1 + D_0)^2)D_0 + \log{\frac{1}{\epsilon}})$
Indeed, $T_2$ is smaller than $T_1$ when $\epsilon$ is small and $\log{\frac{1}{\epsilon}}$ becomes the dominant term. When $\epsilon$ is very small, the superlinear bound from Theorem~\ref{thm:superlinear_rate_2} provides the best complexity, which is
$
T_3=\mathcal{O}\!\!\left(\nicefrac{(\log{\frac{1}{\epsilon}})}{\log\!{\left(\frac{1}{2} \!+\! \sqrt{\frac{1}{4} \!+\! \frac{1}{\Psi(\tilde{B_0}) + (\Psi(\bar{B_0}) + (1 + D_0)^2)D_0}\!\log{\frac{1}{\epsilon}}}\right)}}\! \!\right).
$
Given these three bounds the overall iteration complexity of BFGS for the considered setting is 
%\begin{equation*}
$T= \min\{T_1,T_2,T_3\}$.
%\end{equation*}
Note that, for the special case of $B_0 = a I$ where $a > 0$ is an arbitrary positive constant, the complexity bounds denoted by $T_1,T_2,T_3$ can be further simplified as 
\begin{equation*}
\begin{split}
    T_1= \mathcal{O}\!\left( \Delta_1 \!+\! (1 \!+\! D_0)^2\log{\frac{1}{\epsilon}}\right),\ T_2=\mathcal{O}\!\left( C_1 \!+\! \log{\frac{1}{\epsilon}}\right),\  T_3 = \mathcal{O}\!\left( \frac{\log{\frac{1}{\epsilon}}}{\log{\left(\frac{1}{2} \!+\! \sqrt{\frac{1}{4} \!+\! \frac{1}{C_1}\log{\frac{1}{\epsilon}}}\right)}}\right)\!,
\end{split}
\end{equation*}
where $\Delta_1,\Delta_2$ are defined in \eqref{Delta_1}, \eqref{Delta_2}, and $C_1:= \Delta_2 + (\Delta_1 + (1 + D_0)^2)D_0$. For full iteration complexity details, see Appendix~\ref{proof_of_complexity}.


\noindent\textbf{Line Search Complexity.}
While the previous section characterized the  complexity of BFGS under Assumption~\ref{ass_self_concodant}, analyzing its gradient complexity requires determining the number of gradient queries needed per iteration to obtain an admissible step size. In \cite{qiujiang2024quasinewton}, the authors proposed an efficient log-bisection approach for step size selection in BFGS, satisfying the line search conditions in \eqref{sufficient_decrease} and \eqref{curvature_condition}, and provided a complexity analysis. However, their results apply only to strongly convex functions with Lipschitz-continuous gradients and Hessians. In this section, we examine the line-search complexity of the log-bisection approach from \cite{qiujiang2024quasinewton} when the objective function is strictly convex and strongly self-concordant. Let $\Lambda_t$ denote the average number of iterations in Algorithm~\ref{algo_wolfe} required to terminate after $t$ iterations. The following proposition provides an upper bound for $\Lambda_t$.%, and its proof is in Appendix~\ref{proof_of_linesearch}. 

\begin{proposition}\label{prop:line_search}
Suppose Assumptions~\ref{ass_self_concodant} holds. Let $\{x_t\}_{t\geq 0}$ be generated by BFGS  with step size satisfying the Armijo-Wolfe conditions in \eqref{sufficient_decrease} and \eqref{curvature_condition} and is chosen by Algorithm~\ref{algo_wolfe}. Let $\Lambda_t$ be the average number of the function value and gradient evaluations per iteration in Algorithm~\ref{algo_wolfe} after $t$ iterations. For any initial point $x_0 \!\in\! \mathbb{R}^{d}$ and initial Hessian approximation $B_0\! \in\! \mathbb{S}^d_{++}$,  we have that
\begin{equation*}
    \Lambda_t = \mathcal{O}\left( 1 + \log\left(1 + \frac{\Gamma}{t}\right) + \log\left(1 + \log( 1 + \frac{\Psi(\tilde{B}_{0}) + \Gamma}{t})\right)\right),
\end{equation*}
where $\Gamma = \mathcal{O}(D_0(\Psi(\bar{B}_{0}) + (1 + D_0)^2))$. As a corollary, for the special case of $B_0 = a I$ where $a > 0$, we have $\Lambda_t = \mathcal{O}( 1 + \log(1 + \frac{\tilde{\Gamma}}{t}) + \log(1 + \log( 1 + \frac{\Delta_2 + \tilde{\Gamma}}{t})))$, where $\tilde{\Gamma} = \mathcal{O}\left(D_0(\Delta_1 + (1 + D_0)^2)\right)$. 
\end{proposition}

This proposition implies the average number of iterations in Algorithm~\ref{algo_wolfe} is at most $\mathcal{O}(\log{(1 + \Gamma)})$, which is a constant depending on the initial suboptimality $D_0$ and the initial matrix $B_0$. Moreover, when the number of iterations $T$ exceeds $\Omega(\Psi(\tilde{B}_{0}) + \Gamma)$, the average number of function  and gradient evaluations per iteration for Algorithm~\ref{algo_wolfe} is an absolute constant of $\mathcal{O}(1)$. Thus, even in the worst case, the gradient and iteration complexities remain of the same order, up to logarithmic factors.

\section{Numerical Experiments}\label{sec:experiments}

Next, we present numerical experiments applying BFGS to two functions satisfying Assumptions~\ref{ass_self_concodant}. We report our results using two different choices of initial Hessian approximation  $B_0$: (i) $B_0 = I$, and (ii) $B_0 = cI$, where $c = \frac{s^\top y}{\|s\|^2}$, with $s = x_2 - x_1$, $y = \nabla f(x_2) - \nabla f(x_1)$, where $x_1, x_2$ are randomly selected. The line search parameters are also set as $\alpha = 0.1$ and $\beta = 0.9$. In our experiments, we also report the convergence paths of gradient descent (GD) and accelerated gradient descent (AGD), with step sizes determined using backtracking line search.  

\begin{figure}[t!]
    \centering
    \subfigure[$d = 100$.]{\includegraphics[width=0.45\linewidth]{Figures/Hard_Cubic/d1e2.pdf}}
    \subfigure[$d = 1000$.]{\includegraphics[width=0.45\linewidth]{Figures/Hard_Cubic/d1e3.pdf}}
    \vspace{-3mm}
    \caption{Convergence rates of BFGS with different $B_0$, gradient descent and accelerated gradient descent for solving the hard cubic function with different dimensions.}
    \label{fig:1}  
    \vspace{-5mm}
\end{figure}

The first function that we study is the cubic function from \cite{hard_cubic}

\begin{equation*}
    f(x) = \frac{\omega_1}{12}\left[\sum_{i = 1}^{d - 1}g(v_i^\top x - v_{i + 1}^\top x) - \omega_2 v_1^\top x\right], \quad  \text{where} \ g(x) =
    \begin{cases}
        \frac{1}{3}|x|^3 & |x| \leq \Delta, \\
        \Delta x^2 - \Delta^2 |x| + \frac{1}{3}\Delta^3  & |x| > \Delta.
    \end{cases}  
\end{equation*}

Note that  $g: \mathbb{R} \to \mathbb{R}$. We set the hypermeters of the objective function as $\omega_1 = 4, \omega_2 = 3, \Delta = 1$ and the vectors $\{v_i\}_{i = 1}^{n}$ are set to be the orthogonal unit basis vectors of $\mathbb{R}^{d}$. We study this function as it serves as a benchmark for establishing lower bounds for second-order methods. The second loss is the logistic regression: $f(x) = \frac{1}{N}\sum_{i = 1}^{N}\ln{(1 + e^{-y_i z_i^\top x})}$, where $\{z_i\}_{i = 1}^{N}$ are the data points and $\{y_i\}_{i = 1}^{N}$ are their corresponding labels. We assume that $z_i \in \mathbb{R}^d$ generated with standard normal distribution and $y_i \in \{-1, 1\}$ generated with uniform distribution for all $1 \leq i \leq N$. We choose the number of data points as $N = d$. Note that both the hard cubic function and the logistic regression function are strictly convex and strongly self-concordant; see Appendix~\ref{proof_of_concordance}.%, making them well-suited for evaluating second-order optimization methods. 

\begin{figure}[t!]
    \centering
    \subfigure[$d = 100$.]{\includegraphics[width=0.45\linewidth]{Figures/Logistic/d1e2.pdf}}
    \subfigure[$d = 1000$.]{\includegraphics[width=0.45\linewidth]{Figures/Logistic/d1e3.pdf}}
    \vspace{-3mm}
    \caption{Convergence rates of BFGS with different $B_0$, gradient descent and accelerated gradient descent for solving the logistic regression function with different dimensions.}
    \label{fig:2}
    \vspace{-5mm}
\end{figure}

The convergence paths for the cubic problem are shown in Figure~\ref{fig:1} for various problem dimensions $d$. Initially, the performance of BFGS is worse than that of the first-order gradient descent and accelerated gradient descent methods. However, after approximately $d$ iterations, BFGS significantly outperforms the first-order methods. Notably, for this problem, the performance of BFGS with $B_0 = I$ and $B_0 = cI$ are nearly identical. Figure~\ref{fig:2} shows the convergence paths for the logistic loss across different problem dimensions $d$. Initially, BFGS performs similarly to first-order methods, but after several iterations, it outperforms them. Notably, in this experiment, BFGS with $B_0 = cI$  outperforms BFGS with $B_0 = I$. We also compared the performance of these different optimization methods with respect to the number of gradient evaluations and the time in seconds. Please check Figure~\ref{fig:7} and Figure~\ref{fig:8} in Appendix~\ref{add_experiments} and any other additional numerical experiments.

Moreover, we display the step sizes selected at each iteration by the inexact line search in the BFGS method in Figure~\ref{fig:3}. We observe that the step sizes are initially very small, then gradually increase, and after approximately $d$ iterations, they stabilize at 1 for nearly all subsequent iterations. This confirms our theoretical analysis: BEGS enters the superlinear convergence phase after about $d$ iterations, and there are only limited iterations where the unit step size didn't satisfy the weak Wolfe condition as proved in Lemma~\ref{lem:bad_set}.

Finally, in Figure~\ref{fig:4}, we compare the performance of BFGS, GD, and AGD under a transformation matrix $A$ chosen to be a non-singular ill-conditioned matrix. We observe that the convergence trajectory of BFGS with this transformation is identical to that of the vanilla BFGS method, consistent with the affine invariance of quasi-Newton methods proved in Proposition~\ref{BFGS_affine}. In contrast, the performance of GD and AGD degrades significantly under the transformation matrix, since first-order methods do not possess the affine-invariance property.

\begin{figure}[t!]
    \centering
    \subfigure[$d = 100$.]{\includegraphics[width=0.45\linewidth]{Figures/StepSize/d100.pdf}}
    \subfigure[$d = 1000$.]{\includegraphics[width=0.45\linewidth]{Figures/StepSize/d1000.pdf}}
    \vspace{-3mm}
    \caption{Step size of BFGS with different $B_0$ using inexact line search for solving the hard cubic function with different dimensions.}
    \label{fig:3}
    \vspace{-5mm}
\end{figure}

\begin{figure}[t!]
    \centering
    \subfigure[$d = 100$.]{\includegraphics[width=0.45\linewidth]{Figures/Transform/d100.pdf}}
    \subfigure[$d = 1000$.]{\includegraphics[width=0.45\linewidth]{Figures/Transform/d1000.pdf}}
    \vspace{-3mm}
    \caption{Convergence rates of BFGS with different $B_0$, gradient descent and accelerated gradient descent for solving the hard cubic function with transformation matrix $A$.}
    \label{fig:4}
    \vspace{-5mm}
\end{figure}

\section{Conclusions}\label{sec:conclusion}

We established non-asymptotic global linear and superlinear convergence rates for the BFGS method on strictly convex and strongly self-concordant functions, using Wolfe step sizes. Our guarantees hold for any initial point $x_0 \in \mathbb{R}^d$ and any positive-definite initial Hessian approximation $B_0 \in \mathbb{S}^{d}_{++}$. Our analysis also respects the affine invariance of BFGS. A limitation is the reliance on strong self-concordance; extending results to standard self-concordance is a potential future direction.

%%%%%%%%%%%%%%%%%%%%%%%%%%%%%%%%%%%%%%%%%%%%%%%%%%%%%%%%%%%%%%%%%%%%%%%%%%%%%%%
%%%%%%%%%%%%%%%%%%%%%%%%%%%%%%%%%%%%%%%%%%%%%%%%%%%%%%%%%%%%%%%%%%%%%%%%%%%%%%%
% APPENDIX
%%%%%%%%%%%%%%%%%%%%%%%%%%%%%%%%%%%%%%%%%%%%%%%%%%%%%%%%%%%%%%%%%%%%%%%%%%%%%%%
%%%%%%%%%%%%%%%%%%%%%%%%%%%%%%%%%%%%%%%%%%%%%%%%%%%%%%%%%%%%%%%%%%%%%%%%%%%%%%%

\appendix

\section*{Appendix}

\section{Notations}

Given weight matrix $P\in \mathbb{S}^d_{++}$, recall the weighted vectors defined in \eqref{weighted_vector}. We define the quantity $\hat{\theta}_t$, 
\begin{equation}\label{eq:theta}
    \cos(\hat{\theta}_t) = \frac{-\hat{g}_t^\top \hat{s}_t}{\|\hat{g}_t\|\|\hat{s}_t\|}.
\end{equation}
We also define the following terms which play important roles in all the convergence analysis,
\begin{equation}\label{eq:def_alpha_q_m}
    \hat{p}_t := \frac{f(x_t)-f(x_{t+1})}{-\hat{g}_t^\top \hat{s}_t}, \qquad \hat{q}_t := \frac{\|\hat{g}_t\|^2}{f(x_t)-f(x_*)}, \qquad \hat{m}_t := \frac{\hat{y}_t^\top \hat{s}_t}{\|\hat{s}_t\|^2}, \qquad \hat{n}_t= \frac{\hat{y}_t^\top \hat{s}_t}{-\hat{g}_t^\top \hat{s}_t}. 
\end{equation}
Moreover, we define that
\begin{equation}\label{definition_g}
    \tilde{g}_t := \nabla^2{f(x_*)}^{-\frac{1}{2}}g_t.
\end{equation}
For $\tau_1, \tau_2 \in [0, 1]$, we define the Hessian matrices $J_t$ and $G_t$ as 
\begin{align}\label{def_J}
    J_t := \nabla^2{f(x_t + \tau_1 (x_{t + 1} - x_t))},
\end{align}
\begin{align}\label{def_G}
    G_t := \nabla^2{f(x_t + \tau_2 (x_{*} - x_t))}.
\end{align}
We have that $y_t = \nabla f(x_{t + 1}) - \nabla f(x_t) = J_t(x_{k + 1} - x_k) = J_ts_t$ and $\nabla{f(x_t)} = \nabla{f(x_t)} - \nabla{f(x_*)} = G_t(x_t - x_*)$ for some $\tau_1, \tau_2 \in [0, 1]$ by mean value theorem.

\section{Lemmas and Propositions}

\begin{lemma}\label{lemma_line_search}
    Consider the BFGS method with Armijo-Wolfe inexact line search, where the step size satisfies the conditions in  \eqref{sufficient_decrease} and \eqref{curvature_condition}. If $B_t$ is symmetric positive definite, we have $f(x_{t + 1}) \leq f(x_t)$ and the following results hold:
    \begin{equation}
        \hat{p}_t = \frac{f(x_t)-f(x_{t + 1})}{-{g}_t^\top {s}_t} \geq \alpha, \qquad \hat{n}_t = \frac{{y}_t^\top {s}_t}{- {g}_t^\top {s}_t} \geq 1 - \beta.
    \end{equation}
\end{lemma}

\begin{proof}
Please check Lemma 2.1 in \cite{qiujiang2024quasinewton}.
\end{proof}

\begin{proposition}\label{lemma_BFGS}
Let $\{{B}_t\}_{t\geq 0}$ be the Hessian approximation matrices generated by the BFGS update in \eqref{eq:BFGS_update}. For a given weight matrix $P\in \mathbb{S}^d_{++}$, recall the weighted vectors and the weighted matrix in \eqref{weighted_vector}. Then, we have that for any $t \geq 0$,
\begin{equation}\label{eq:potential_decrease}
    \Psi(\hat{B}_{t+1}) \leq \Psi(\hat{B}_t) + \frac{\|\hat{y}_t\|^2}{\hat{y}_t^\top \hat{s}_t} -1 + \log \frac{\cos^2\hat{\theta}_t}{\hat{m}_t},
\end{equation}
where $\hat{m}_t$ is defined in \eqref{eq:def_alpha_q_m} and $\cos(\hat{\theta}_t)$ is defined in \eqref{eq:theta}. As a corollary, we have that for any $t \geq 1$,
\begin{equation}\label{eq:sum_of_logs}
    \sum_{i = 0}^{t - 1} \log{\frac{\cos^2(\hat{\theta}_i)}{\hat{m}_i}} \geq  - \Psi(\hat{B}_{0}) + \sum_{i = 0}^{t - 1}\left(1-\frac{\|\hat{y}_i\|^2}{\hat{y}_i^\top \hat{s}_i} \right).
\end{equation}
\end{proposition}

\begin{proof}
    Please check Proposition 2 in \cite{jin2024non}.
\end{proof}

\begin{lemma}\label{lemma_omega}
Recall the definition of function $\omega(x)$ in \eqref{definition_omega} and define the function $\omega_{*}(x)$ as
\begin{equation}\label{definition_omega_star}
    \omega_{*}(x) := - x - \log{(1 - x)}, \qquad \forall x < 1.
\end{equation}
We have that
\begin{enumerate}[(a)]
    \item $\omega(x)$ is increasing function for $x > 0$ and decreasing function for $-1 < x < 0$. Moreover, $\omega(x) \geq 0$ for all $x > -1$.
    \item When $x \geq 0$, we have that $\omega(x) \geq \frac{x^2}{2(1 + x)}$. 
    \item When $-1 < x \leq 0$, we have that $\omega(x) \geq \frac{x^2}{2 + x}$.
    \item When $0 < x < 1$, we have that $\omega_{*}(x) \leq \frac{x^2}{2(1 - x)}$.
    \item We have $\sqrt{2x} + \frac{2x}{3} \leq \omega^{-1}(x) \leq \sqrt{2x} + x$, where $\omega^{-1}$ is the inverse function of $\omega(x)$ when $x > 0.$
\end{enumerate}
\end{lemma}

\begin{proof}
    Please check Lemma G.1 in \cite{qiujiang2024quasinewton} for the proof of (a), (b), and (c). Please check Lemma 5.1.5 from \cite{convexlecture} for the proof of (d) and Lemma A.1 from \cite{antonglobal} for the proof of (e).
\end{proof}

\begin{lemma}\label{lemma_concordance}
    Recall the definition of function $\omega(x)$ and $\omega_*(x)$ in \eqref{definition_omega} and \eqref{definition_omega_star}. If Assumption~\ref{ass_self_concodant} holds, we have that
    \begin{equation}
        f(y) \geq f(x) + g(x)^\top(y - x) + \frac{4}{M^2}\omega(\frac{M}{2}\|y - x\|_x).
    \end{equation}
    Moreover, when $\|y - x\|_x < \frac{2}{M}$, we have that
    \begin{equation}
        f(y) \leq f(x) + g(x)^\top(y - x) + \frac{4}{M^2}\omega_{*}(\frac{M}{2}\|y - x\|_x).
    \end{equation}
\end{lemma}

\begin{proof}
    Check Theorem 5.1.8 and Theorem 5.1.9 of \cite{convexlecture}.
\end{proof}

\begin{lemma}\label{lemma_Hessian}
    Suppose Assumptions~\ref{ass_self_concodant} holds, and recall the definitions of the matrices $J_t$ and $G_t$ in \eqref{def_J} and \eqref{def_G}, and the quantity $D_t$ in \eqref{suboptimality}. The following statements hold: 
    \begin{enumerate}[(a)]
        \item For any $t \geq 0$, we have that
        \begin{equation}\label{eq:H_t_vs_H*}
            \frac{1}{1 + D_t}\nabla^2{f(x_*)} \preceq \nabla^2{f(x_t)} \preceq (1 + D_t)\nabla^2{f(x_*)} .
        \end{equation}
        \item For any $t \geq 0$, we have that
        \begin{equation}\label{eq:G_t_vs_H*}
            \frac{1}{1 + D_t}\nabla^2{f(x_*)} \preceq G_t \preceq (1 + D_t)\nabla^2{f(x_*)} .
        \end{equation}
        \item For any $t \geq 0$ and $\tau \in [0, 1]$, we have that
        \begin{equation}\label{eq:G_t_vs_H_t}
            \frac{1}{1 + D_t}G_t \preceq \nabla^2{f(x_t + \tau(x_* - x_t))} \preceq (1 + D_t)G_t .
        \end{equation}
        \item Suppose that $f(x_{t + 1}) \leq f(x_t)$ for $t \geq 0$, we have that
        \begin{equation}\label{eq:J_t_vs_H*}
            \frac{1}{1 + D_t}\nabla^2{f(x_*)} \preceq J_t \preceq (1 + D_t)\nabla^2{f(x_*)} .
        \end{equation}
    \end{enumerate}
\end{lemma}

\begin{proof}

    Proof for (a): In Lemma~\ref{lemma_concordance}, take $y = x_t$ and $x = x_*$, we have that $\omega(\frac{M}{2}\|x_t - x_*\|_{x_*}) \leq \frac{M^2}{4}(f(x_t) - f(x_*))$. Hence, we have that
    \begin{align*}
        \|x_t - x_*\|_{x_*} \leq \frac{2}{M}\omega^{-1}(\frac{M^2}{4}C_t).
    \end{align*}
    In Assumptions~\ref{ass_self_concodant}, take $x = x_t$, $y = x_*$ and $z = x_*$, we prove that 
    \begin{equation}
        \nabla^2{f(x_t)} \preceq (1 + M\|x_t - x_*\|_{x_*})\nabla^2{f(x_*)} \preceq (1 + 2\omega^{-1}(\frac{M^2}{4}C_t))\nabla^2{f(x_*)} = (1 + D_t)\nabla^2{f(x_*)}.
    \end{equation}
    Similarly, take $x = x_*$, $y = x_t$, $z = x_*$ and $w = x_t$, we prove that 
    \begin{equation}
        \nabla^2{f(x_t)} \succeq \frac{1}{1 + M\|x_t - x_*\|_{x_*}}\nabla^2{f(x_*)} \succeq \frac{1}{1 + 2\omega^{-1}(\frac{M^2}{4}C_t)}\nabla^2{f(x_*)} = \frac{1}{1 + D_t}\nabla^2{f(x_*)}.
    \end{equation}

    Proof for (b): Recall the definition of $G_t$ in \eqref{def_G}. Notice that $\|x_t + \tau_2 (x_{*} - x_t) - x_*\|_{x_*} = (1 - \tau_2)\|x_t - x_*\|_{x_*} \leq \|x_t - x_*\|_{x_*} \leq \frac{2}{M}\omega^{-1}(\frac{M^2}{4}C_t)$. Similar to the arguments in (a), in Assumptions~\ref{ass_self_concodant}, take $x = x_t + \tau_2 (x_{*} - x_t)$, $y = x_*$ and $z = x_*$, we prove that 
    \begin{equation}
    \begin{split}
        & G_t \preceq (1 + M\|x_t + \tau (x_{*} - x_t) - x_*\|_{x_*})\nabla^2{f(x_*)} \\
        & \preceq (1 + 2\omega^{-1}(\frac{M^2}{4}C_t))\nabla^2{f(x_*)} = (1 + D_t)\nabla^2{f(x_*)}.
    \end{split}
    \end{equation}
    Similarly, take $x = x_*$, $y = x_t + \tau_2 (x_{*} - x_t)$ and $z = x_*$, we prove that 
    \begin{equation}
    \begin{split}
        & G_t \succeq \frac{1}{1 + M\|x_t + \tau (x_{*} - x_t) - x_*\|_{x_*}}\nabla^2{f(x_*)} \\
        & \succeq \frac{1}{1 + 2\omega^{-1}(\frac{M^2}{4}C_t)}\nabla^2{f(x_*)} = \frac{1}{1 + D_t}\nabla^2{f(x_*)}.
    \end{split}
    \end{equation}
    Proof for (c): Notice that $\|x_t + \tau (x_{*} - x_t) - (x_t + \tau_2 (x_{*} - x_t))\|_{x_*} = |\tau - \tau_2|\|x_t - x_*\|_{x_*} \leq \|x_t - x_*\|_{x_*} \leq \frac{2}{M}\omega^{-1}(\frac{M^2}{4}C_t)$. Similar to the arguments in (a), in Assumptions~\ref{ass_self_concodant}, take $x = x_t + \tau (x_{*} - x_t)$, $y = x_t + \tau_2 (x_{*} - x_t)$ and $z = x_*$, we prove that 
    \begin{equation}
    \begin{split}
        & \nabla^2{f(x_t + \tau(x_* - x_t))}  \preceq (1 + M\|x_t + \tau (x_{*} - x_t) - x_*\|_{x_*})G_t \\
        & \preceq (1 + 2\omega^{-1}(\frac{M^2}{4}C_t))G_t = (1 + D_t)G_t.
    \end{split}
    \end{equation}
    Similarly, take $x = x_t + \tau_2 (x_{*} - x_t)$, $y = x_t + \tau (x_{*} - x_t)$ and $z = x_*$, we prove that 
    \begin{equation}
    \begin{split}
        & \nabla^2{f(x_t + \tau(x_* - x_t))}  \succeq \frac{1}{1 + M\|x_t + \tau (x_{*} - x_t) - x_*\|_{x_*}}G_t \\
        & \succeq \frac{1}{1 + 2\omega^{-1}(\frac{M^2}{4}C_t)}G_t = \frac{1}{1 + D_t}G_t.
    \end{split}
    \end{equation}
    Proof for (d): Recall the definition of $J_t$ in \eqref{def_J}. Notice that $\|x_t + \tau (x_{t + 1} - x_t) - x_*\|_{x_*} = (1 - \tau)\|x_t - x_*\|_{x_*} + \tau \|x_{t + 1} - x_*\|_{x_*} \leq (1 - \tau)\frac{2}{M}\omega^{-1}(\frac{M^2}{4}C_t) + \tau \frac{2}{M}\omega^{-1}(\frac{M^2}{4}C_t) \leq \frac{2}{M}\omega^{-1}(\frac{M^2}{4}C_t)$ where the last inequality holds since $f(x_{t + 1}) \leq f(x_t)$ and $\omega^{-1}(x)$ is increasing function. Hence, similar to the arguments in (a), in Assumptions~\ref{ass_self_concodant}, take $x = x_t + \tau (x_{t + 1} - x_t)$, $y = x_*$ and $z = x_*$, we prove that 
    \begin{equation}
    \begin{split}
        & J_t \preceq (1 + M\|x_t + \tau (x_{t + 1} - x_t) - x_*\|_{x_*})\nabla^2{f(x_*)} \\
        & \preceq (1 + 2\omega^{-1}(\frac{M^2}{4}C_t))\nabla^2{f(x_*)} = (1 + D_t)\nabla^2{f(x_*)}.
    \end{split}
    \end{equation}
    Similarly, take $x = x_*$, $y = x_t + \tau (x_{*} - x_t)$ and $z = x_*$, we prove that 
    \begin{equation}
    \begin{split}
        & J_t \succeq \frac{1}{1 + M\|x_t + \tau (x_{t + 1} - x_t) - x_*\|_{x_*}}\nabla^2{f(x_*)} \\
        & \succeq \frac{1}{1 + 2\omega^{-1}(\frac{M^2}{4}C_t)}\nabla^2{f(x_*)} = \frac{1}{1 + D_t}\nabla^2{f(x_*)}.
    \end{split}
    \end{equation}
\end{proof}

\begin{proposition}\label{lemma_bound}
    Let $\{x_t\}_{t\geq 0}$ be the iterates generated by BFGS. Recall the definitions of weighted vectors in \eqref{weighted_vector} and notations in \eqref{eq:def_alpha_q_m}. Then, for any weight matrix $P \in \mathbb{S}^{d}_{++}$ and any $t \geq 1$, we have
    \begin{equation}\label{eq:product}
        \frac{f(x_{t}) - f(x_*)}{f(x_0) - f(x_*)} \leq \bigg(1 - \bigg(\prod_{i = 0}^{t - 1} \hat{p}_i \hat{q}_{i} \hat{n}_i\frac{\cos^2(\hat{\theta}_{i})}{\hat{m}_{i}}\bigg)^{\frac{1}{t}}\bigg)^{t}.
    \end{equation}  
\end{proposition}

\begin{proof}
Please check Proposition 1 in \cite{jin2024non}
\end{proof}

\begin{lemma}\label{lemma_g}
    Suppose Assumption~\ref{ass_self_concodant} holds for the convex objective function $f(x)$ and recall the definition $\tilde{g}_t$ in \eqref{definition_g} and $D_t$ in \eqref{suboptimality}. We have the following condition,
    \begin{equation}
        \frac{\|\tilde{g}_t\|^2}{f(x_t) - f(x_*)} \geq \frac{1}{1 + D_t}, \qquad \forall t \geq 0.
    \end{equation}
\end{lemma}

\begin{proof}
    Since $f$ is convex, we know that for any $x, y \in \mathbb{R}^d$, we have $f(y) \geq f(x) + g(x)^\top(y - x)$. Take  $x = x_t$ and $y = x_*$, we obtain that
    \begin{equation*}
        f(x_t) - f(x_*) \leq g_t^\top(x_t - x_*).
    \end{equation*}
    Using mean value theorem and the fact that $\nabla{f(x_*)} = 0$, we have that
    \begin{equation*}
        g_t = \nabla{f(x_t)} = \nabla{f(x_t)} - \nabla{f(x_*)} = G_t(x_t - x_*),
    \end{equation*}
    where $G_t$ is defined in \eqref{def_G} for some $\tau_2 \in [0, 1]$. Hence, we prove that
    \begin{equation*}
    \begin{split}
        & f(x_t) - f(x_*) \leq g_t^\top(x_t - x_*) = g_t^\top G_t^{-1}g_t \\
        & \leq \left(1 + 2\omega^{-1}(\frac{M^2}{4}C_t)\right)g_t^\top \nabla^2{f(x_*)}^{-1}g_t = \left(1 + D_t\right)\|\tilde{g}_t\|^2,
    \end{split}
    \end{equation*}
    where we use the result in \eqref{eq:G_t_vs_H*} from Lemma~\ref{lemma_Hessian}.
\end{proof}

\begin{lemma}\label{lemma_eigenvalue}
    Suppose Assumption~\ref{ass_self_concodant} holds for the convex objective function $f(x)$ and recall the definition $\tilde{g}_t$ in \eqref{definition_g} and $D_t$ in \eqref{suboptimality}. We have the following condition,
    \begin{equation}
        \frac{2}{(1 + D_t)^2} \leq \frac{\|\tilde{g}_t\|^2}{f(x_t) - f(x_*)} \leq 2(1 + D_t)^2, \qquad \forall t \geq 0.
    \end{equation}
\end{lemma}

\begin{proof}    
    By applying Taylor's theorem with Lagrange remainder, there exists $\tilde{\tau}_t \in [0,1]$ such that 
    \begin{equation}\label{lemma_eigenvalue_proof_1}
        \begin{split}
            f(x_t) & = f(x_*) + \nabla{f(x_*)}^\top (x_t - x_*) + \frac{1}{2}(x_t - x_*)^\top \nabla^2{f(x_t + \tilde{\tau}_t(x_* - x_t))}(x_t - x_*) \\
            & = f(x_*) + \frac{1}{2}(x_t - x_*)^\top \nabla^2{f(x_t + \tilde{\tau}_t(x_* - x_t))}(x_t - x_*),
    \end{split}
    \end{equation}
    where we used the fact that $\nabla{f(x_*)} = 0$ in the last equality. Moreover, by the fundamental theorem of calculus, we have
    \begin{equation*}
        \nabla{f(x_t)} - \nabla{f(x_*)} = \nabla^2{f(x_t + \tau (x_{*} - x_t))} (x_t - x_*) = G_t(x_t - x^*),
    \end{equation*}
    where we use the definition of $G_t$ in \eqref{def_G}. Since $\nabla f(x_*) = 0$ and we denote $g_t = \nabla{f(x_t)}$, this further implies that
    \begin{equation}\label{lemma_eigenvalue_proof_2}
        x_t - x_* = G_t^{-1}(\nabla{f(x_t)} - \nabla{f(x_*)}) = G_t^{-1} g_t.
    \end{equation}
    Combining \eqref{lemma_eigenvalue_proof_1} and \eqref{lemma_eigenvalue_proof_2} leads to 
    \begin{equation}\label{eq:f_x_k-f_*}
        f(x_t) - f(x_*) = \frac{1}{2}g_t^\top G_t^{-1}\nabla^2{f(x_t + \tilde{\tau}_t(x_* - x_t))}G_t^{-1}g_t. 
    \end{equation}
    Based on \eqref{eq:G_t_vs_H_t} in Lemma~\ref{lemma_Hessian}, we have $\nabla^2{f(x_t + \tilde{\tau}_t (x_{*} - x_{t}))} \preceq (1 + D_t)G_t$, which implies
    \begin{equation}\label{lemma_eigenvalue_proof_3}
        G_t^{-1}\nabla^2{f(x_t + \tilde{\tau}_t (x_{*} - x_{t}))}G_t^{-1} \preceq (1 + D_t)^2G_t^{-1}.
    \end{equation}
    Moreover, it follows from \eqref{eq:G_t_vs_H*} in Lemma~\ref{lemma_Hessian} that $\frac{1}{1 + D_t} \nabla^2{f(x_*)} \preceq G_t$, which implies that
    \begin{equation}\label{lemma_eigenvalue_proof_4}
        G_t^{-1} \preceq (1 + D_t)(\nabla^2{f(x_*)})^{-1}.
    \end{equation}
    Combining \eqref{lemma_eigenvalue_proof_3} and \eqref{lemma_eigenvalue_proof_4}, we obtain that
    \begin{equation*}
        G_t^{-1}\nabla^2{f(x_t + \tilde{\tau}_t (x_{*} - x_{t}))}G_t^{-1} \preceq (1 + D_t)^2(\nabla^2{f(x_*)})^{-1},
    \end{equation*}
    and hence 
    \begin{equation*}
        g_t^\top G_t^{-1}\nabla^2{f(x_t + \tilde{\tau}_t (x_{*} - x_{t}))}G_t^{-1}g_t \leq (1 + D_t)^2 g_t^\top(\nabla^2{f(x_*)})^{-1} g_t.
    \end{equation*}
    By using \eqref{eq:f_x_k-f_*} and the fact that $g_t^\top(\nabla^2{f(x_*)})^{-1} g_t = \|\tilde{g}_t\|^2$, we obtain that
    \begin{equation*}
        \frac{\|\tilde{g}_t\|^2}{f(x_t)-f(x_*)} \geq \frac{2}{(1 + D_t)^2},
    \end{equation*}
    and the left claim follows. Using the similar method, we can prove the right claim. 
\end{proof}

\begin{lemma}\label{lemma_decrease}
    Suppose Assumption~\ref{ass_self_concodant} holds and $C_t \leq \min\{\frac{1}{16}, \frac{4}{M^2}\omega(\frac{1}{32})\}$ and $\rho_t \geq \frac{15}{16}$ at iteration $t$, then we have that
    \begin{equation}
        f(x_t + d_t) \leq f(x_t).
    \end{equation}
\end{lemma}

\begin{proof}
    Notice that using \eqref{eq:H_t_vs_H*} from Lemma~\ref{lemma_Hessian}, we have that
    \begin{equation}\label{lemma_decrease_1}
        d_t^\top\nabla^2{f}(x_t)d_t \leq (1 + D_t)d_t^\top\nabla^2{f}(x_*)d_t = (1 + D_t)\|\tilde{d}_t\|^2.
    \end{equation}
    Since $-\tilde{g}_t^\top \tilde{d}_t \leq \|\tilde{g}_t\|\|\tilde{d}_t\|$ by Cauchy–Schwarz inequality where $\tilde{g}_t = \nabla^2{f(x_*)}^{-\frac{1}{2}}g_t$, we obtain
    \begin{equation}\label{lemma_decrease_2}
        \|\tilde{d}_t\| = \|\tilde{g}_t\|\frac{\|\tilde{d}_t\|}{\|\tilde{g}_t\|} \leq \|\tilde{g}_t\|\frac{\|\tilde{d}_t\|^2}{-\tilde{g}_t^\top \tilde{d}_t} = \frac{1}{\rho_t}\|\tilde{g}_t\|.
    \end{equation}
    Using the right inequality in Lemma~\ref{lemma_eigenvalue}, we have that
    \begin{equation}\label{lemma_decrease_3}
        \|\tilde{g}_t\|^2 \leq 2(1 + D_t)^{2}(f(x_t) - f(x_*)) = 2(1 + D_t)^{2}C_t.
    \end{equation}
    Leveraging \eqref{lemma_decrease_1}, \eqref{lemma_decrease_2} and \eqref{lemma_decrease_3}, we obtain that
    \begin{align*}
        d_t^\top\nabla^2{f}(x_t)d_t \leq (1 + D_t)\|\tilde{d}_t\|^2 \leq \frac{1 + D_t}{\rho_t^2}\|\tilde{g}_t\|^2 \leq \frac{2(1 + D_t)^3}{\rho_t^2}C_t.
    \end{align*}
    Since $C_t \leq \frac{1}{16}$, $D_t = 2\omega^{-1}(\frac{M^2}{4}C_t) \leq \frac{1}{16}$ and $\rho_t \geq \frac{15}{16}$, we have that
    \begin{align*}
        \sqrt{d_t^\top\nabla^2{f}(x_t)d_t} \leq \sqrt{\frac{2(1 + D_t)^3}{\rho_t^2}C_t} \leq \frac{13}{30} < 1.
    \end{align*}
    Applying the second inequality from Lemma~\ref{lemma_concordance} with $x = x_t$ and $y = x_t + d_t$, we have that
    \begin{equation*}
        f(x_t +d_t) \leq f(x_t) + g_t^\top d_t + \omega_{*}(\|d_t\|_{x_t}),
    \end{equation*}
    since $\|d_t\|_{x_t} = \sqrt{d_t^\top\nabla^2{f}(x_t)d_t} < 1$. Using (d) from Lemma~\ref{lemma_omega}, we have that
    \begin{equation*}
        f(x_t +d_t) - f(x_t) \leq g_t^\top d_t + \omega_{*}(\|d_t\|_{x_t}) \leq g_t^\top d_t + \frac{\|d_t\|_{x_t}^2}{2(1 - \|d_t\|_{x_t})}.
    \end{equation*}
    Applying the fact that $\|d_t\|_{x_t} = \sqrt{d_t^\top\nabla^2{f}(x_t)d_t} \leq \frac{13}{30}$ and \eqref{lemma_decrease_1}, we have that
    \begin{equation}\label{lemma_decrease_4}
    \begin{split}
        & f(x_t +d_t) - f(x_t) \leq g_t^\top d_t + \frac{\|d_t\|_{x_t}^2}{2(1 - \|d_t\|_{x_t})} \leq g_t^\top d_t + \frac{15}{17}\|d_t\|_{x_t}^2 \leq g_t^\top d_t + \frac{15}{17}(1 + D_t)\|\tilde{d}_t\|^2 \\
        & = g_t^\top d_t + \frac{15}{17}(1 + D_t)\frac{\|\tilde{d}_t\|^2}{-\tilde{g}_t^\top \tilde{d}_t}(-g_t^\top d_t) = -g_t^\top d_t(\frac{15}{17}(1 + D_t)\frac{\|\tilde{d}_t\|^2}{-\tilde{g}_t^\top \tilde{d}_t} - 1) \\
        & = -g_t^\top d_t(\frac{15}{17}\frac{1 + D_t}{\rho_t} - 1)
    \end{split}
    \end{equation}
    Notice that $-g_t^\top d_t = -g_t^\top B_t^{-1} g_t > 0$ and when $D_t \leq \frac{1}{16}$ and $\rho_t \geq \frac{15}{16}$, we can verify that
    \begin{equation*}
        \frac{15}{17}\frac{1 + D_t}{\rho_t} \leq 1.
    \end{equation*}
    Therefore, \eqref{lemma_decrease_4} implies the conclusion that
    \begin{equation*}
        f(x_t + d_t) - f(x_t) \leq 0.
    \end{equation*}
\end{proof}

\begin{lemma}\label{lem:lower_bound}
    Recall $\hat{p}_t = \frac{f(x_t) - f(x_{t+1})}{-\hat{g}_t^\top \hat{s}_t}$ and $\hat{n}_t= \frac{\hat{y}_t^\top \hat{s}_t}{-\hat{g}_t^\top \hat{s}_t}$ defined in \eqref{eq:def_alpha_q_m}. If the unit step size $\eta_t = 1$ satisfies the Armijo-Wolfe conditions \eqref{sufficient_decrease} and \eqref{curvature_condition}, then we have that
    \begin{equation}\label{eq:lower_bound}
        \hat{p}_t \geq 1 - \frac{1 + D_t}{2\rho_t}, \qquad \hat{n}_t\geq \frac{1}{(1 + D_t)\rho_t}.
    \end{equation}
\end{lemma}

\begin{proof}
    Please check Lemma 6.1 in \cite{qiujiang2024quasinewton}. The only difference is that $C_t$ is replaced by $D_t$ defined in \eqref{suboptimality}.
\end{proof}

\begin{proposition}\label{lemma_BFGS_superlinear}
Let $\{{B}_t\}_{t\geq 0}$ be the Hessian approximation matrices generated by the BFGS update in \eqref{eq:BFGS_update}. Suppose Assumptions~\ref{ass_self_concodant} holds and $f(x_{t + 1}) \leq f(x_t)$ for any $t \geq 0$. Recall the definition of $\Psi(.)$ in \eqref{potential_function} and $D_t$ in \eqref{suboptimality}, we have
\begin{equation}\label{eq:omega}
    \sum_{i = 0}^{t - 1} \omega(\rho_i - 1) \leq  \Psi(\tilde{B}_{0}) + 2\sum_{i = 0}^{t - 1}D_i.
\end{equation}
\end{proposition}

\begin{proof}
    Please check Proposition G.2 in \cite{qiujiang2024quasinewton}. The only difference is $C_t$ is replaced by $D_t$ defined in \eqref{suboptimality}.
\end{proof}

\section{Proof of Lemmas, Propositions and Theorems}

\subsection{Proof of Proposition~\ref{BFGS_affine}}\label{proof_affine}

    We use induction to prove that $x_t = A\dot{x}_t$ and $B_t = (A^{-1})^\top \dot{B}_t A^{-1}$ for any $t \geq 0$. Notice that when $t = 0$, we already have that $x_0 = A\dot{x}_0$ and $B_0 = (A^{-1})^\top \dot{B}_0 A^{-1}$ since $\dot{x}_0 = A^{-1}x_0$ and $\dot{B}_0 = A^\top B_0 A$ where $A$ is non-singular. Suppose that the conditions hold for $t = k$ with $k \geq 0$, i.e., $x_k = A\dot{x}_k$ and $B_k = (A^{-1})^\top \dot{B}_k A^{-1}$. We consider the case $t = k + 1$. We have
    \begin{equation*}
    \begin{split}
        x_{k + 1} & = x_k - \eta_k B_k^{-1}\nabla{f(x_k)} = A\dot{x}_k - \eta_k A \dot{B}_k^{-1} A^\top \nabla{f(A\dot{x}_k)} \\
        & = A\dot{x}_k - \eta_k A \dot{B}_k^{-1} A^\top (A^\top)^{-1} \nabla{\phi(\dot{x}_k)} = A(\dot{x}_k - \eta_k \dot{B}_k^{-1} \nabla{\phi(\dot{x}_k)}) = A \dot{x}_{k + 1}.
    \end{split}
    \end{equation*}
    Suppose that $\dot{s}_k = \dot{x}_{k + 1} - \dot{x}_k$ and $\dot{y}_k = \nabla{\phi(\dot{x}_{k + 1})} - \nabla{\phi(\dot{x}_k)}$, we have that $s_k = x_{k + 1} - x_k = A\dot{s}_k$ and $y_k = \nabla{f(x_{k + 1})} - \nabla{f(x_k)} = \nabla{f(A\dot{x}_{k + 1})} - \nabla{f(A\dot{x}_k)} = (A^\top)^{-1}(\nabla{\phi(\dot{x}_{k + 1})} - \nabla{\phi(\dot{x}_k)}) = (A^\top)^{-1}\dot{y}_k$. Hence, we have that
    \begin{equation}
    \begin{split}
        B_{k + 1} & = B_k - \frac{B_k s_k s_k^\top B_k}{s_k^\top B_k s_k} + \frac{y_k y_k^\top}{s_k^\top y_k} \\
        & = (A^{-1})^\top \dot{B}_k A^{-1} - \frac{(A^{-1})^\top \dot{B}_k A^{-1} A \dot{s}_k \dot{s}_k^\top A^\top (A^{-1})^\top \dot{B}_k A^{-1}}{\dot{s}_k A^\top (A^{-1})^\top \dot{B}_k A^{-1} A \dot{s}_k} + \frac{(A^\top)^{-1}\dot{y}_k \dot{y}_k^\top A^{-1}}{\dot{s}_k A^\top (A^\top)^{-1}\dot{y}_k} \\
        & = (A^{-1})^\top \left(\dot{B}_k - \frac{\dot{B}_k \dot{s}_k \dot{s}_k^\top \dot{B}_k}{\dot{s}_k^\top \dot{B}_k \dot{s}_k} + \frac{\dot{y}_k \dot{y}_k^\top}{\dot{s}_k^\top \dot{y}_k}\right)A^{-1} \\
        & = (A^{-1})^\top \dot{B}_{k + 1} A^{-1}.
    \end{split}
    \end{equation}
    We prove that $x_{k + 1} = A\dot{x}_{k + 1}$ and $B_{k + 1} = (A^{-1})^\top \dot{B}_{k + 1} A^{-1}$. Therefore, we prove that $\dot{x}_t = A^{-1}x_t$ and $\dot{B}_t = A^\top B_t A$ for any $t \geq 0$ using induction. It is obvious that $\phi(\dot{x}_t) = f(A\dot{x}_t) = f(AA^{-1}x_t) = f(x_t)$ for any $t \geq 0$. Therefore, The BFGS quasi-Newton method is affine invariant.

\subsection{Proof of Theorem~\ref{thm:first_linear_rate}}\label{proof_thm_first_linear}

    We choose $P = \left(1 + D_0)\right)\nabla^2{f(x_*)}$ throughout the proof. Note that given this weight matrix $P$, it can be easily verified that for any $t \geq 0$,
    \begin{equation}
        \frac{\|\hat{y}_t\|^2}{\hat{s}_t^\top \hat{y}_t} = \frac{\hat{s}_t^\top \hat{J}_t^2 \hat{s}_t}{\hat{s}_t^\top \hat{J}_t \hat{s}_t} \leq \|\hat{J}_t\| = \frac{\|\nabla^2{f(x_*)}^{-\frac{1}{2}}J_t \nabla^2{f(x_*)}^{-\frac{1}{2}}\|}{1 + D_0} \leq \frac{1 + D_t}{1 + D_0} \leq 1,
    \end{equation}
    where $J_t$ is defined in \eqref{def_J} and we use \eqref{eq:J_t_vs_H*} in Lemma~\ref{lemma_Hessian} as well as the fact that $f(x_{t + 1}) \leq f(x_t)$, $D_t \leq D_0$ and $\omega^{-1}$ is increasing function. Hence, we use \eqref{eq:sum_of_logs} in Proposition~\ref{lemma_BFGS} with $\hat{B}_0 = \bar{B}_0$ defined in \eqref{def_B_bar} to obtain
    \begin{equation*}
        \sum_{i = 0}^{t - 1} \log{\frac{\cos^2(\hat{\theta}_i)}{\hat{m}_i}} \geq  - \Psi(\bar{B}_{0}) + \sum_{i = 0}^{t - 1}\left(1-\frac{\|\hat{y}_i\|^2}{\hat{s}_i^\top \hat{y}_i} \right) \geq - \Psi(\bar{B}_{0}),
    \end{equation*}
    which further implies that 
    \begin{equation*}%\label{lemma_6_proof_3}
    \qquad \prod_{i = 0}^{t-1} \frac{\cos^2(\hat{\theta}_i)}{\hat{m}_i} \geq e^{-\Psi(\bar{B}_{0})}.
    \end{equation*}
    Moreover, for the choice $P = \left(1 + D_0\right)\nabla^2{f(x_*)}$, it can be shown that
    \begin{equation}
        \hat{q}_t = \frac{\|\tilde{g}_t\|^2}{(1 + D_0)(f(x_t) - f(x_*))} \geq \frac{1}{(1 + D_0)(1 + D_t)} \geq \frac{1}{(1 + D_0)^2}.
    \end{equation}
    by using result in Lemma~\ref{lemma_g}. From Lemma~\ref{lemma_line_search}, we know $\hat{p}_t \geq \alpha$ and $\hat{n}_t\geq 1 - \beta$, which lead to
    \begin{equation*}
        \prod_{i=0}^{t-1} \frac{\hat{p}_i \hat{n}_i \hat{q}_i}{\hat{m}_i} \cos^2(\hat{\theta}_i) \geq \prod_{i=0}^{t-1} \hat{p}_i \prod_{i=0}^{t-1} \hat{q}_i \prod_{i=0}^{t-1} \hat{n}_i \prod_{i=0}^{t-1} \frac{\cos^2(\hat{\theta}_i)}{\hat{m}_i} \geq \left(\frac{\alpha(1 - \beta)}{(1 + D_0)^2}\right)^t  e^{-\Psi(\bar{B}_{0})}.
    \end{equation*}
    Thus, it follows from Proposition~\ref{lemma_bound} that 
    \begin{align*}
        \frac{f(x_{t}) - f(x_*)}{f(x_0) - f(x_*)} \leq \left[1 - \left(\prod_{i = 0}^{t-1} \frac{\hat{p}_i\hat{q}_i\hat{n}_i}{\hat{m}_i} \cos^2(\hat{\theta}_i)\right)^{\frac{1}{t}}\right]^{t} \leq \left(1 - \frac{\alpha(1 - \beta)e^{-\frac{\Psi(\bar{B}_0)}{t}}}{(1 + D_0)^2}\right)^{t}.
    \end{align*}
    This completes the proof. \eqref{eq:first_linear_rate_2} can be easily verified since when $t \geq \Psi(\bar{B}_0)$, we have $e^{-\frac{\Psi(\bar{B}_0)}{t}} \geq \frac{1}{3}$.

\subsection{Proof of Theorem~\ref{thm:second_linear_rate}}\label{proof_thm_second_linear}

    First, we prove the following result holds:
    \begin{equation}\label{proposition_3_1}
        \frac{f(x_t) - f(x_*)}{f(x_0) - f(x_*)} \leq \left(1 - 2\alpha(1 - \beta)e^{- \frac{\Psi(\tilde{B}_{0}) + 3\sum_{i = 0}^{t - 1}D_i}{t}} \right)^{t}.
    \end{equation}

    We choose the weight matrix as $P = \nabla^2 f(x_*)$ throughout the proof. Similar to the proof of Theorem~\ref{thm:first_linear_rate}, we start from the key inequality in \eqref{eq:product}, but we apply different bounds on the $\hat{q}_t$ and $\frac{\cos^2(\hat{\theta}_t)}{\hat{m}_t}$. Specifically, we have that
    \begin{equation}
        \frac{\|\hat{y}_t\|^2}{\hat{s}_t^\top \hat{y}_t} = \frac{\hat{s}_t^\top \hat{J}_t^2 \hat{s}_t}{\hat{s}_t^\top \hat{J}_t \hat{s}_t} \leq \|\hat{J}_t\| = \|\nabla^2{f(x_*)}^{-\frac{1}{2}}J_t \nabla^2{f(x_*)}^{-\frac{1}{2}}\| \leq 1 + D_t.
    \end{equation}
    where $J_t$ is defined in \eqref{def_J} and we use \eqref{eq:J_t_vs_H*} in Lemma~\ref{lemma_Hessian} as well as the fact that $f(x_{t + 1}) \leq f(x_t)$. Hence, we use \eqref{eq:sum_of_logs} in Proposition~\ref{lemma_BFGS} with $\hat{B}_0 = \tilde{B}_0$ defined in \eqref{def_B_bar} to obtain
    \begin{equation*}
        \sum_{i = 0}^{t - 1} \log{\frac{\cos^2(\hat{\theta}_i)}{\hat{m}_i}} \geq  - \Psi(\Tilde{B}_{0}) + \sum_{i = 0}^{t - 1}\left(1-\frac{\|\hat{y}_i\|^2}{\hat{s}_i^\top \hat{y}_i} \right) \geq - \Psi(\Tilde{B}_{0}) - \sum_{i = 0}^{t - 1}D_i,
    \end{equation*}
    which further implies that
    \begin{equation}\label{eq:prodcut_cos}
        \prod_{i = 0}^{t - 1} \frac{\cos^2(\hat{\theta}_i)}{\hat{m}_i} \geq e^{- \Psi(\Tilde{B}_{0}) - \sum_{i = 0}^{t - 1}D_i}.
    \end{equation}
    Moreover, since $\hat{q}_t = \frac{\|\tilde{g}_t\|^2}{f(x_t) - f(x_*)} \geq \frac{2}{(1 + 2\omega^{-1}(\frac{M^2}{4}C_t))^2}$ for any $t\geq 0$ by using Lemma~\ref{lemma_eigenvalue}, we obtain that
    \begin{equation}\label{eq:product_alpha_q}
        \prod_{i = 0}^{t - 1} \hat{q}_i \geq \prod_{i = 0}^{t - 1} \frac{2}{(1 + D_i)^2} \geq 2^t\prod_{i = 0}^{t - 1} e^{-2D_i} = 2^te^{-2\sum_{i = 0}^{t - 1}D_i},
    \end{equation}
    where we use the inequality $1+x \leq e^x$ for any $x \in \mathbb{R}$. From Lemma~\ref{lemma_line_search}, we know $\hat{p}_t \geq \alpha$ and $\hat{n}_t\geq 1 - \beta$, which lead to
    \begin{equation}\label{eq:product_kappa}
        \prod_{i = 0}^{t - 1} \hat{p}_i\hat{n}_i \geq \alpha^{t} (1 - \beta)^{t}.
    \end{equation}
    Combining \eqref{eq:prodcut_cos}, \eqref{eq:product_alpha_q}, \eqref{eq:product_kappa} and \eqref{eq:product} from Proposition~\ref{lemma_bound}, we prove that
    \begin{align*}
        \frac{f(x_{t}) - f(x_*)}{f(x_0) - f(x_*)} \leq \left[1 - \left(\prod_{i = 0}^{t-1} \frac{\hat{p}_i\hat{q}_i\hat{n}_i}{\hat{m}_i} \cos^2(\hat{\theta}_i)\right)^{\frac{1}{t}}\right]^{t} \leq \left(1 - 2\alpha(1 - \beta)e^{- \frac{\Psi(\Tilde{B}_{0}) + 3\sum_{i = 0}^{t - 1}D_i}{t}} \right)^{t}.
    \end{align*}
    This completes the proof. Notice that when
    \begin{equation}
        t \geq \Psi(\Tilde{B}_{0}) + 3\sum_{i = 0}^{t - 1}D_i,
    \end{equation}
    \eqref{proposition_3_1} implies the condition that 
    \begin{equation}
        \frac{f(x_t) - f(x_*)}{f(x_0) - f(x_*)} \leq \left(1 -\frac{2\alpha(1 - \beta)}{e} \right)^{t}\leq \left(1 - \frac{2\alpha(1 - \beta)}{3} \right)^{t},
    \end{equation}
    which leads to the linear rate in \eqref{eq:second_linear_rate}. 
    
    Hence, it is sufficient to establish an upper bound on $\sum_{i = 0}^{t - 1}D_i$. We decompose the sum into two parts: $\sum_{i=0}^{\ceil{\Psi(\bar{B}_0)} - 1} D_i$ and $\sum_{i = \ceil{\Psi(\bar{B}_0)}}^t D_i$. For the first part, note that since $f(x_{i + 1}) \leq f(x_i)$ by Lemma~\ref{lemma_line_search}, we also have $D_{i + 1} \leq D_i$ for $i \geq 0$ using the fact that $\omega^{-1}$ is increasing. Hence, we have $\sum_{i=0}^{\ceil{\Psi(\bar{B}_0)} - 1} D_i \leq D_0 \ceil{\Psi(\bar{B}_0)} \leq D_0 (\Psi(\bar{B}_0) + 1)$. Moreover, by Theorem~\ref{thm:first_linear_rate}, when $t \geq \Psi(\bar{{B}}_0)$ we have 
    \begin{equation*}
        \frac{f(x_t) - f(x_*)}{f(x_0) - f(x_*)} \leq \left(1 - e^{-\frac{\Psi(\bar{B}_0)}{t}}\frac{\alpha(1 - \beta)}{(1 + D_0)^2}\right)^{t} \leq \left(1 - \frac{\alpha(1 - \beta)}{3(1 + D_0)^2}\right)^{t}.
    \end{equation*}
    Hence, using the fact that $\omega^{-1}(x) \leq \sqrt{2x} + x$ and the definition of $D_t$ in \eqref{suboptimality}, we obtain that 
    \begin{align*}
        & \sum_{i=\ceil{\Psi(\bar{B}_0)}}^t D_i \leq \sum_{i=\ceil{\Psi(\bar{B}_0)}}^t (\frac{M}{2}\sqrt{2C_i} + \frac{M^2}{4}C_i ) = \frac{\sqrt{2}M}{2}\sum_{i=\ceil{\Psi(\bar{B}_0)}}^t \sqrt{C_i} + \frac{M^2}{4}\sum_{i=\ceil{\Psi(\bar{B}_0)}}^t C_i \\
        & = \frac{\sqrt{2}M}{2}\sqrt{C_0} \sum_{i=\ceil{\Psi(\bar{B}_0)}}^t \sqrt{\frac{f(x_i)-f(x_*)}{f(x_0)-f(x_*)}} + \frac{M^2}{4} C_0 \sum_{i=\ceil{\Psi(\bar{B}_0)}}^t \frac{f(x_i)-f(x_*)}{f(x_0)-f(x_*)}\\
        & \leq \frac{\sqrt{2}M}{2}\sqrt{C_0} \sum_{i=\ceil{\Psi(\bar{B}_0)}}^t \left(1 - \frac{\alpha(1 - \beta)}{3(1 + D_0)^2}\right)^{\frac{i}{2}} + \frac{M^2}{4}C_0 \sum_{i=\ceil{\Psi(\bar{B}_0)}}^t \left(1 - \frac{\alpha(1 - \beta)}{3(1 + D_0)^2}\right)^{i} \\
        & \leq \frac{\sqrt{2}M}{2}\sqrt{C_0} \sum_{i=1}^{\infty} \left(1 - \frac{\alpha(1 - \beta)}{3(1 + D_0)^2}\right)^{\frac{i}{2}} + \frac{M^2}{4}C_0 \sum_{i=1}^{\infty} \left(1 - \frac{\alpha(1 - \beta)}{3(1 + D_0)^2}\right)^{i} \\
        & \leq \frac{\sqrt{2}M}{2}\sqrt{C_0} \left(\frac{6(1 + D_0)^2}{\alpha(1 - \beta)} - 1\right) + \frac{M^2}{4}C_0 \left(\frac{3(1 + D_0)^2}{\alpha(1 - \beta)} - 1\right)
    \end{align*}
    where we used the fact that $\sum_{i=1}^{\infty}(1-\rho)^{\frac{i}{2}} = \frac{\sqrt{1-\rho}}{1-\sqrt{1-\rho}} = \frac{\sqrt{1-\rho}+1-\rho}{\rho}\leq \frac{2}{\rho}-1$ and $\sum_{i=1}^{\infty}(1-\rho)^{i} = \frac{1 - \rho}{1 - (1 - \rho)} = \frac{1}{\rho} - 1$ for any $\rho \in (0,1)$. Hence, by combining both inequalities, we have 
    \begin{equation}\label{sum_of_C_t}
    \begin{split}
        & \sum_{i = 0}^{t - 1} D_i = \sum_{i=0}^{\ceil{\Psi(\bar{B}_0)}-1} D_i + \sum_{i=\ceil{\Psi(\bar{B}_0)}}^t D_i \\
        & \leq D_0 \Psi(\bar{B}_{0}) + \frac{\sqrt{2}M}{2}\sqrt{C_0}\frac{6(1 + D_0)^2}{\alpha(1 - \beta)} + \frac{M^2}{4}C_0\frac{3(1 + D_0)^2}{\alpha(1 - \beta)} \\
        & = D_0\Psi(\bar{B}_{0}) + (M\sqrt{2C_0} + \frac{M^2}{4}C_0)\frac{3(1 + D_0)^2}{\alpha(1 - \beta)} \leq D_0\left(\Psi(\bar{B}_{0}) + \frac{3(1 + D_0)^2}{\alpha(1 - \beta)}\right),
    \end{split}
    \end{equation}
    where the last inequality is due to (d) from Lemma~\ref{lemma_omega} and the definition of $D_t$ in \eqref{suboptimality}. Hence, when 
    $$
    t \geq \Psi(\Tilde{B}_{0}) + 3D_0\left(\Psi(\bar{B}_{0}) + \frac{3(1 + D_0)^2}{\alpha(1 - \beta)}\right) \geq \Psi(\Tilde{B}_{0}) + 3\sum_{i = 0}^{t - 1} D_i,
    $$
    using result from \eqref{proposition_3_1}, we have the linear convergence rate in \eqref{eq:second_linear_rate}.

\subsection{Proof of Lemma~\ref{lem:unit_step_size}}\label{proof_lem:unit_step_size}

    Denote $\bar{x}_{t + 1} = x_t + d_t$ and $\bar{s}_t = \bar{x}_{t + 1} - x_t = d_t$. Since $\delta_1 \leq \min\{\frac{1}{16}, \frac{4}{M^2}\omega(\frac{1}{32})\}$ and $\delta_2 \geq \frac{15}{16}$, we have $f(\bar{x}_{t + 1}) \leq f(x_t)$ from Lemma~\ref{lemma_decrease}. Using Taylor's expansion, we have that $f(\bar{x}_{t + 1}) = f(x_t) + g_t^\top d_t + \frac{1}{2}d_t^\top \nabla^2{f(x_t + \hat{\tau} (\bar{x}_{t + 1} - x_{t}))}d_t$, where $\hat{\tau} \in [0, 1]$. Hence, we have
    \begin{equation*}%\label{lemma_unit_stepsize_1}
    \begin{split}
        & \frac{f(x_t) - f(\bar{x}_{k+1})}{-g_t^\top d_t} = \frac{-g_t^\top d_t - \frac{1}{2}d_t^\top \nabla^2{f(x_t + \hat{\tau} (\bar{x}_{t + 1} - x_{t}))}d_t}{-g_t^\top d_t} \\
        & = 1 - \frac{1}{2}\frac{d_t^\top \nabla^2{f(x_t + \hat{\tau} (\bar{x}_{t + 1} - x_{t}))}d_t}{-g_t^\top d_t} \geq 1 - \frac{1 + D_t}{2}\frac{d_t^\top \nabla^2{f(x_*)}d_t}{-g_t^\top d_t} = 1 - \frac{1 + D_t}{2\rho_t},
    \end{split}
    \end{equation*}
    where we apply the \eqref{eq:J_t_vs_H*} from Lemma~\ref{lemma_Hessian} since $f(\bar{x}_{t + 1}) \leq f(x_t)$. Therefore, when $C_t \leq \delta_1 \leq \frac{4}{M^2}\omega(\frac{\sqrt{2(1 - \alpha)}- 1}{2})$ and $\rho_t \geq \delta_2 \geq \frac{1}{\sqrt{2(1 - \alpha)}}$, we obtain that $\frac{f(x_t) - f(\bar{x}_{k+1})}{-g_t^\top d_t} \geq 1 - \frac{1 + D_t}{2\rho_t} = 1 - \frac{1 + 2\omega^{-1}(\frac{M}{4}C_t)}{2\rho_t} \geq \alpha$ and unit step size $\eta_t = 1$ satisfies the sufficient condition \eqref{sufficient_decrease}. 
    
    Similarly, using \eqref{eq:J_t_vs_H*} from Lemma~\ref{lemma_Hessian} since $f(\bar{x}_{t + 1}) \leq f(x_t)$ and denote $\bar{g}_{k + 1} = \nabla{f}(\bar{x}_{t + 1})$, $\bar{y}_t = \bar{g}_{k + 1} - g_t$, we have that
    \begin{equation*}%\label{lemma_unit_stepsize_2}
    \frac{\bar{y}_t^\top \bar{s}_t}{-g_t^\top \bar{s}_t} = \frac{\bar{s}_t^\top J_t \bar{s}_t}{-g_t^\top \bar{s}_t} = \frac{d_t^\top J_t d_t}{-g_t^\top d_t} \geq \frac{1}{1 + D_t}\frac{d_t^\top \nabla^2{f(x_*)}d_t}{-g_t^\top d_t} = \frac{1}{(1 + D_t)\rho_t}.
    \end{equation*}
    Therefore, when $C_t \leq \delta_1 \leq \frac{4}{M^2}\omega(\frac{1}{2}(\frac{1}{\sqrt{1 - \beta}} - 1))$ and $\rho_t \leq \delta_3 = \frac{1}{\sqrt{1 - \beta}}$, we obtain that $\frac{\bar{y}_t^\top \bar{s}_t}{-g_t^\top \bar{s}_t} \geq \frac{1}{(1 + D_t)\rho_t} = \frac{1}{(1 + 2\omega^{-1}(\frac{M^2}{4}C_t))\rho_t} \geq 1 - \beta$, which indicates that $\bar{g}_{t + 1}^\top d_t = \bar{g}_{t + 1}^\top \bar{s}_t = \bar{y}_t^\top \bar{s}_t + g_{t}^\top \bar{s}_t \geq -g_{t}^\top \bar{s}_t(1 -\beta) + g_{t}^\top \bar{s}_t = \beta g_{t}^\top \bar{s}_t = \beta g_{t}^\top d_t$. Hence, unit step size $\eta_t = 1$ satisfies the curvature condition \eqref{curvature_condition}. Therefore, we prove that when $C_t \leq \delta_1$ and $\delta_2 \leq \rho_t \leq \delta_3$, step size $\eta_t = 1$ satisfies the Armijo-Wolfe conditions \eqref{sufficient_decrease} and \eqref{curvature_condition}.

\subsection{Proof of Lemma~\ref{lem:bad_set}}\label{proof_lem:bad_set}

Since in Theorem~\ref{thm:first_linear_rate}, we already prove that
\begin{equation*}
    \frac{f(x_t) - f(x_*)}{f(x_0) - f(x_*)} \leq \left(1 - \frac{\alpha(1 - \beta)e^{-\frac{\Psi(\bar{B}_0)}{t}}}{(1 + D_0)^2}\right)^{t}.
\end{equation*}
This implies that
\begin{equation*}
    C_t \leq \left(1 - \frac{\alpha(1 - \beta)e^{-\frac{\Psi(\bar{B}_0)}{t}}}{(1 + D_0)^2}\right)^{t}C_0.
\end{equation*}
When $t \geq \Psi(\bar{B}_0)$, we obtain that
\begin{equation*}
    C_t \leq \left(1 - \frac{\alpha(1 - \beta)}{3(1 + D_0)^2}\right)^{t}C_0.
\end{equation*}
When $t \geq \frac{3(1 + D_0)^2}{\alpha(1 - \beta)}\log{\frac{C_0}{\delta_1}}$, we obtain that
\begin{equation*}
    C_t \leq \left(1 - \frac{\alpha(1 - \beta)}{3(1 + D_0)^2}\right)^{t}C_0 \leq \delta_1.
\end{equation*}
Therefore, the first claim in \eqref{definition_t_0} follows.
    
Define $I_1  = \{t_0 \leq i \leq t - 1:\; \rho_t < \delta_2\}$ and $I_2  = \{t_0 \leq i \leq t - 1:\; \rho_t > \delta_3\}$, we know that $|I| = |I_1| + |I_2|$. Notice that for $t \in I_1$, we have that $\rho_t - 1 < \delta_2 - 1 < 0$ since $\delta_2 < 1$ and the function $\omega(x)$ defined in \eqref{definition_omega} is decreasing for $-1 < x < 0$ from (a) in Lemma~\ref{lemma_omega}. Hence, we have that $ \sum_{i \in I_1}\omega(\rho_i - 1) \geq  \sum_{i \in I_1}\omega(\delta_2 - 1) = \omega(\delta_2 - 1)|I_1|$. Similarly, we have that for $t \in I_2$, we have that $\rho_i - 1 > \delta_3 - 1 > 0$ since $\delta_3 > 1$ and the function $\omega(x)$ is increasing for $x > 0$ from (a) in Lemma~\ref{lemma_omega}. Hence, we have that $ \sum_{i \in I_2}\omega(\rho_i - 1) \geq \sum_{i \in I_2}\omega(\delta_3 - 1) = \omega(\delta_3 - 1)|I_2|$. Using \eqref{eq:omega} from Proposition~\ref{lemma_BFGS_superlinear}, we have that $\sum_{i = 0}^{t - 1} \omega(\rho_i - 1) \leq \Psi(\tilde{B}_{0}) + 2\sum_{i = 0}^{t  - 1}D_i$ for any $t \geq 1$. Therefore, we obtain that
    \begin{align*}
        & \Psi(\tilde{B}_{0}) + 2\sum_{i = 0}^{t - 1}D_i \geq \sum_{i = 0}^{t - 1} \omega(\rho_i - 1) \geq \sum_{i \in I_1}\omega(\beta_i - 1) + \sum_{i \in I_2}\omega(\beta_i - 1) \\
        & \geq \omega(\delta_2 - 1)|I_1| + \omega(\delta_3 - 1)|I_2|  \geq \min\{\omega(\delta_2 - 1), \omega(\delta_3 - 1)\}(|I_1| + |I_2|),
    \end{align*}
    which leads to the result
    \begin{equation}\label{lem:bad_set_proof_1}
    \begin{split}
        |I| & = |I_1| + |I_2| \leq \frac{\Psi(\tilde{B}_{0}) + 2\sum_{i = 0}^{t - 1}D_i}{\min\{\omega(\delta_2 - 1), \omega(\delta_3 - 1)\}} \\
        & = \delta_4 \left(\Psi(\tilde{B}_{0}) + 2\sum_{i = 0}^{t - 1}D_i\right) \leq \delta_4 \left(\Psi(\tilde{B}_{0}) + 2D_0(\Psi(\bar{B}_{0}) + \frac{3(1 + D_0)^2}{\alpha(1 - \beta)})\right),
    \end{split}
    \end{equation}
    where $\delta_4 = \frac{1}{\min\{\omega(\delta_2 - 1), \omega(\delta_3 - 1)\}}$ and the last inequality is due to \eqref{sum_of_C_t}.

\subsection{Proof of Theorem~\ref{thm:superlinear_rate_2}}\label{proof_thm:superlinear_rate_2}

    First, we prove the following result:
    \begin{equation}
        \frac{f(x_t) - f(x_*)}{f(x_0) - f(x_*)} \leq \left(\frac{\delta_6 t_0 + \delta_7 \Psi(\tilde{B}_0) + \delta_8 \sum_{i = 0}^{t - 1}D_i}{t}\right)^{t}.
    \end{equation}
    
    We choose the weight matrix as $P = \nabla^2 f(x_*)$ throughout the proof. Taking the sum from $0$ to $t - 1$ in inequality \eqref{eq:potential_decrease} of the Proposition~\ref{lemma_BFGS}, we obtain that
    \begin{equation}\
        \sum_{i = 0}^{t - 1} \log{\frac{\cos^2(\hat{\theta}_i)}{\hat{m}_i}} \geq  - \Psi(\tilde{B}_{0}) + \sum_{i = 0}^{t - 1}\left(1 - \frac{\|\hat{y}_i\|^2}{\hat{y}_i^\top \hat{s}_i} \right), \qquad  \forall t \geq 1.
    \end{equation}
    Notice that $\frac{\|\hat{y}_i\|^2}{\hat{y}_i^\top \hat{s}_i} = \frac{\|\hat{J}_i \hat{s}_i\|^2}{\hat{s}_i^\top \hat{J}_i \hat{s}_i} \leq \|\hat{J}_i\| \leq 1 + D_i$ where $\hat{J}_i$ is defined in \eqref{weighted_vector} with $P = \nabla^2 f(x_*)$ and we use \eqref{eq:J_t_vs_H*} from Lemma~\ref{lemma_Hessian}. Therefore, we have that
    \begin{equation}\label{proposition_4_1}
        \prod_{i = 0}^{t - 1} \frac{\cos^2(\hat{\theta}_i)}{\hat{m}_i} \geq e^{- \Psi(\Tilde{B}_{0}) + \sum_{i = 0}^{t - 1}\left(1 - \frac{\|\hat{y}_i\|^2}{\hat{y}_i^\top \hat{s}_i} \right)} \geq e^{- \Psi(\Tilde{B}_{0}) - \sum_{i = 0}^{t - 1}D_i}.
    \end{equation}
    where $\cos(\hat{\theta}_i)$ is defined in \eqref{eq:theta}. Recall the definitions in \eqref{eq:def_alpha_q_m} and the results in Lemma~\ref{lemma_eigenvalue}, we have
    \begin{equation}\label{proposition_4_2}
        \prod_{i = 0}^{t - 1} \hat{q}_i \geq \prod_{i = 0}^{t - 1}\frac{2}{(1 + D_i)^2} \geq 2^t e^{-2\sum_{i = 0}^{t - 1}D_i}.
    \end{equation}
    Recall the definition of the set $I = \{t_0 \leq i \leq t - 1:\; \rho_i \notin [\delta_2,\delta_3]\}$ and define the set $\bar{I} = \{t_0 \leq i \leq t - 1:\; \rho_i \in [\delta_2, \delta_3]\}$ for any $t > t_0$. Then, we have that
    \begin{equation}\label{proposition_4_3}
        \prod_{i = 0}^{t - 1} \hat{p}_i\hat{n}_i = \prod_{i = 0}^{t_0 - 1} \hat{p}_i\hat{n}_i \prod_{i \in I} \hat{p}_i\hat{n}_i \prod_{i \in \bar{I}} \hat{p}_i\hat{n}_i.
    \end{equation}
    From Lemma~\ref{lemma_line_search}, we know $\hat{p}_t \geq \alpha$ and $\hat{n}_t\geq 1 - \beta$ for any $t \geq 0$, which lead to
    \begin{equation}\label{proposition_4_4}
        \prod_{i = 0}^{t_0 - 1} \hat{p}_i\hat{n}_i \geq \alpha^{t_0} (1 - \beta)^{t_0} = \frac{1}{2^{t_0}}e^{-t_0\log{\frac{1}{2\alpha(1 - \beta)}}}.
    \end{equation}
    \begin{equation}\label{proposition_4_5}
    \begin{split}
        & \prod_{i \in I} \hat{p}_i\hat{n}_i \geq \prod_{i \in I} \alpha(1 - \beta) = \frac{1}{2^{|I|}}e^{-|I|\log{\frac{1}{2\alpha(1 - \beta)}}} \\
        & \geq \frac{1}{2^{|I|}}e^{-|I|\log{\frac{1}{2\alpha(1 - \beta)}}} \geq \frac{1}{2^{|I|}}e^{-\delta_4\Big(\Psi(\tilde{B}_{0}) + 2\sum_{i = 0}^{t - 1}D_i\Big)\log{\frac{1}{2\alpha(1 - \beta)}}},
    \end{split}
    \end{equation}
    where the second inequality holds since $\log{\frac{1}{2\alpha(1 - \beta)}} > 0$ and the last inequality holds since \eqref{lem:bad_set_proof_1} from the proof of Lemma~\ref{lem:bad_set}. Notice that when index $i \in \bar{I}$, we have $C_i \leq \delta_1$ from Lemma~\ref{lem:bad_set} and $\rho_i \in [\delta_2, \delta_3]$. Applying Lemma~\ref{lem:lower_bound} and Lemma~\ref{lem:unit_step_size}, we know that for $i \in \bar{I}$, $\eta_i = 1$ satisfies the Armijo-Wolfe conditions \eqref{sufficient_decrease}, \eqref{curvature_condition} and we have $\hat{p}_i \geq 1 - \frac{1 + D_i}{2\rho_i} > 0$ and $\hat{n}_i \geq \frac{1}{(1 + D_i)\rho_i}$ from \eqref{eq:lower_bound}. Hence, we obtain that
    \begin{equation}\label{proposition_4_6}
        \prod_{i \in \bar{I}} \hat{p}_i\hat{n}_i \geq \frac{1}{2^{\bar{I}}}\prod_{i \in \bar{I}}(2 - \frac{1 + D_i}{\rho_i})\frac{1}{(1 + D_i)\rho_i} \geq \frac{1}{2^{|\bar{I}|}}e^{-\sum_{i \in \bar{I}}D_i}\prod_{i \in \bar{I}}(2 - \frac{1 + D_i}{\rho_i})\frac{1}{\rho_i},
    \end{equation}
    where the last inequality holds since $\frac{1}{1 + D_i} \geq e^{-D_i}$. Using the fact that $\log{x} \geq 1 - \frac{1}{x}$, we obtain
    \begin{equation}\label{proposition_4_7}
    \begin{split}
        & \prod_{i \in \bar{I}}(2 - \frac{1 + D_i}{\rho_i})\frac{1}{\rho_i} = \prod_{i \in \bar{I}}e^{\log{(2 - \frac{1 + D_i}{\rho_i})} - \log{\rho_i}} \geq \prod_{i \in \bar{I}}e^{1 - \frac{1}{2 - \frac{1 + D_i}{\rho_i}} - \log{\rho_i}} \\
        & = \prod_{i \in \bar{I}}e^{\frac{\rho_i - 1 - D_i}{2\rho_i - 1 - D_i} - \log{\rho_i}} = \prod_{i \in \bar{I}}e^{\frac{\rho_i - 1 - \log{\rho_i} + 2(1 - \rho_i)\log{\rho_i} - (1 - \log{\rho_i})D_i}{2\rho_i - 1 - D_i}} \\
        & = \prod_{i \in \bar{I}}e^{\frac{\omega(\rho_i - 1) + 2(1 - \rho_i)\log{\rho_i} - (1 - \log{\rho_i})D_i }{2\rho_i - 1 - D_i}} \geq \prod_{i \in \bar{I}}e^{\frac{- 2(\rho_i - 1)\log{\rho_i} - (1 - \log{\rho_i})D_i }{2\rho_i - 1 - D_i}} \\
        & = \prod_{i \in \bar{I}}e^{-\frac{2(\rho_i - 1)\log{\rho_i} + (1 - \log{\rho_i})D_i}{2\rho_i - 1 - D_i}} \geq \prod_{i \in \bar{I}}e^{-\frac{2(\rho_i - 1)\log{\rho_i} + (1 - \log{\delta_2})D_i}{2\delta_2 - 1 - {1}/{16}}} = \prod_{i \in \bar{I}}e^{-\frac{2(\rho_i - 1)\log{\rho_i} + (1 - \log{\delta_2})D_i}{2\delta_2 - {17}/{16}}} ,
    \end{split}
    \end{equation}
    where the second inequality holds since $\omega(\rho_i - 1) \geq 0$ and the third inequality holds since $\rho_i \geq \delta_2$ due to $i \in \bar{I}$ and $C_i \leq \delta_1 \leq \frac{4}{M^2}\omega(\frac{1}{32})$, $D_i = 2\omega^{-1}(\frac{M^2}{4}C_i) \leq \frac{1}{16}$ due to $i \geq t_0$ and Lemma~\ref{lem:bad_set}. Notice that $2\rho_i - 1 - D_i \geq 2\delta_2 - 1 - \frac{1}{16} > 0$ for all $i \in \bar{I}$ since $\rho_i \geq \delta_2 \geq \frac{15}{16}$.

    When $\rho_i \geq 1$, using $\log{\rho_i} \leq \rho_i - 1$, (b) in Lemma~\ref{lemma_omega} and $\rho_i \leq \delta_3$ due to $i \in \bar{I}$, we have that
    \begin{equation}\label{proposition_4_8}
        (\rho_i - 1)\log{\rho_i} \leq (\rho_i - 1)^2 \leq 2\rho_i\omega(\rho_i - 1) \leq 2\delta_3\omega(\rho_i - 1).
    \end{equation}
    Similarly, when $\rho_i < 1$, using $\log{\rho_i} \geq 1 - \frac{1}{\rho_i}$, (c) in Lemma~\ref{lemma_omega} and $\rho_i \geq \delta_2$ due to $i \in \bar{I}$, we have
    \begin{equation}\label{proposition_4_9}
        (\rho_i - 1)\log{\rho_i} \leq \frac{(\rho_i - 1)^2}{\rho_i} \leq \frac{\rho_i + 1}{\rho_i}\omega(\rho_i - 1) \leq (1 + \frac{1}{\delta_2})\omega(\rho_i - 1).
    \end{equation}
    Combining \eqref{proposition_4_7}, \eqref{proposition_4_8} and \eqref{proposition_4_9}, we obtain that
    \begin{equation}\label{proposition_4_10}
    \begin{split}
        & \prod_{i \in \bar{I}}(2 - \frac{1 + D_i}{\rho_i})\frac{1}{\rho_i} \geq \prod_{i \in \bar{I}}e^{-\frac{2(\rho_i - 1)\log{\rho_i} + (1 - \log{\delta_2})D_i}{2\delta_2 - 17/16}} = \prod_{i \in \bar{I}}e^{-\frac{2(\rho_i - 1)\log{\rho_i}}{2\delta_2 - 17/16}} \prod_{i \in \bar{I}}e^{-\frac{(1 - \log{\delta_2})D_i}{2\delta_2 - 17/16}} \\
        & = \prod_{i \in \bar{I}, \rho_i < 1}e^{-\frac{2(\rho_i - 1)\log{\rho_i}}{2\delta_2 - 17/16}} \prod_{i \in \bar{I}, \rho_i \geq 1}e^{-\frac{2(\rho_i - 1)\log{\rho_i}}{2\delta_2 - 17/16}} \prod_{i \in \bar{I}}e^{-\frac{(1 - \log{\delta_2})D_i}{2\delta_2 - 17/16}} \\
        & \geq \prod_{i \in \bar{I}, \rho_i < 1}e^{-\frac{2(1 + \frac{1}{\delta_2})\omega(\rho_i - 1)}{2\delta_2 - 17/16}} \prod_{i \in \bar{I}, \rho_i \geq 1}e^{-\frac{4\delta_3\omega(\rho_i - 1)}{2\delta_2 - 17/16}} \prod_{i \in \bar{I}}e^{-\frac{(1 - \log{\delta_2})D_i}{2\delta_2 - 17/16}} \\
        & = e^{- \frac{2 + \frac{2}{\delta_2}}{2\delta_2 - 17/16}\sum_{i \in \bar{I}, \rho_i < 1}\omega(\rho_i - 1) - \frac{4\delta_3}{2\delta_2 - 17/16}\sum_{i \in \bar{I}, \rho_i \geq 1}\omega(\rho_i - 1) - \frac{(1 - \log{\delta_2})}{2\delta_2 - 17/16}\sum_{i \in \bar{I}}D_i} \\
        & \geq e^{- \delta_5\left(\sum_{i \in \bar{I}, \rho_i < 1}\omega(\rho_i - 1) + \sum_{i \in \bar{I}, \rho_i \geq 1}\omega(\rho_i - 1)\right) - \frac{(1 - \log{\delta_2})}{2\delta_2 - 17/16}\sum_{i \in \bar{I}}D_i} \\
        & = e^{- \delta_5\sum_{i \in \bar{I}}\omega(\rho_i - 1) - \frac{(1 - \log{\delta_2})}{2\delta_2 - 17/16}\sum_{i \in \bar{I}}D_i}
    \end{split}
    \end{equation}
    where $\delta_5 = \max\{\frac{2 + \frac{2}{\delta_2}}{2\delta_2 - 17/16}, \frac{4\delta_3}{2\delta_2 - 17/16}\}$. Combining \eqref{proposition_4_6} and \eqref{proposition_4_10}, we obtain that
    \begin{equation}\label{proposition_4_11}
    \begin{split}
        & \prod_{i \in \bar{I}} \hat{p}_i\hat{n}_i \geq \frac{1}{2^{|\bar{I}|}}e^{-\sum_{i \in \bar{I}}D_i}\prod_{i \in \bar{I}}(2 - \frac{1 + D_i}{\rho_i})\frac{1}{\rho_i} \\
        & \geq \frac{1}{2^{|\bar{I}|}}e^{- \delta_5\sum_{i \in \bar{I}}\omega(\rho_i - 1) - (1 + \frac{1 - \log{\delta_2}}{2\delta_2 - 17/16})\sum_{i \in \bar{I}}D_i} \geq \frac{1}{2^{|\bar{I}|}}e^{- \delta_5\sum_{i = 0}^{t - 1}\omega(\rho_i - 1) - \frac{2\delta_2 - \delta_1 - \log{\delta_2}}{2\delta_2 - 17/16}\sum_{i = 0}^{t - 1}D_i} \\
        & \geq \frac{1}{2^{|\bar{I}|}}e^{- \delta_5\Big(\Psi(\tilde{B}_{0}) + 2\sum_{i = 0}^{t - 1}D_i\Big) - \frac{2\delta_2 - 1/16 - \log{\delta_2}}{2\delta_2 - 17/16}\sum_{i = 0}^{t - 1}D_i},
    \end{split}
    \end{equation}
    where the last inequality is due to \eqref{eq:omega} from Lemma~\ref{lemma_omega}. Combining \eqref{proposition_4_3}, \eqref{proposition_4_4}, \eqref{proposition_4_5} and \eqref{proposition_4_11}, we obtain that
    \begin{equation}\label{proposition_4_12}
    \begin{split}
        & \prod_{i = 0}^{t - 1} \hat{p}_i\hat{n}_i = \prod_{i = 0}^{t_0 - 1} \hat{p}_i\hat{n}_i\prod_{i \in I} \hat{p}_i\hat{n}_i \prod_{i \in \bar{I}} \hat{p}_i\hat{n}_i \\
        & \geq \frac{1}{2^{t_0}}e^{-t_0\log{\frac{1}{2\alpha(1 - \beta)}}}\frac{1}{2^{|I|}}e^{-\delta_4\Big(\Psi(\tilde{B}_{t_0}) + 2\sum_{i = 0}^{t - 1}D_i\Big)\log{\frac{1}{2\alpha(1 - \beta)}}} \\
        & \quad \frac{1}{2^{|\bar{I}|}}e^{- \delta_5\Big(\Psi(\tilde{B}_{0}) + 2\sum_{i = 0}^{t - 1}D_i\Big) - \frac{2\delta_2 - 1/16 - \log{\delta_2}}{2\delta_2 - 17/16}\sum_{i = 0}^{t - 1}D_i} \\
        & = \frac{1}{2^t}e^{-\Big(t_0\log{\frac{1}{2\alpha(1 - \beta)}} + (\delta_4\log{\frac{1}{2\alpha(1 - \beta)}} + \delta_5) \Psi(\tilde{B}_{0}) + (2\delta_4\log{\frac{1}{2\alpha(1 - \beta)}} + 2\delta_5 + \frac{2\delta_2 - 1/16 - \log{\delta_2}}{2\delta_2 - 17/16})\sum_{i = 0}^{t - 1}D_i \Big)}.
    \end{split}
    \end{equation}
    Leveraging \eqref{proposition_4_1}, \eqref{proposition_4_2}, \eqref{proposition_4_12} with \eqref{eq:product} from Proposition~\ref{lemma_bound}, we prove that
    \begin{align*}
        & \frac{f(x_{t}) - f(x_*)}{f(x_{0}) - f(x_*)} \leq \left[1 - \left(\prod_{i = 0}^{t - 1} \hat{p}_i \hat{q}_{i} \hat{n}_i\frac{\cos^2(\hat{\theta}_{i})}{\hat{m}_{i}}\right)^{\frac{1}{t}}\right]^{t} = \left[1 - \left(\prod_{i = 0}^{t - 1} \hat{p}_i \hat{n}_{i} \prod_{i = 0}^{t - 1} \hat{q}_i \prod_{i = 0}^{t - 1}\frac{\cos^2(\hat{\theta}_{i})}{\hat{m}_{i}}\right)^{\frac{1}{t}}\right]^{t} \\
        & \leq \left(1 - e^{- \frac{t_0\log{\frac{1}{2\alpha(1 - \beta)}} + (1 + \delta_4\log{\frac{1}{2\alpha(1 - \beta)}} + \delta_5) \Psi(\tilde{B}_{0}) + (2 + 2\delta_4\log{\frac{1}{2\alpha(1 - \beta)}} + 2\delta_5 + \frac{2\delta_2 - 1/16 - \log{\delta_2}}{2\delta_2 - 17/16})\sum_{i = 0}^{t - 1}D_i}{t}}\right)^{t} \\
        & = \left(1 - e^{-\frac{\delta_6 t_0 + \delta_7 \Psi(\tilde{B}_0) + \delta_8 \sum_{i = 0}^{t - 1}D_i}{t}}\right)^{t} \leq \left(\frac{\delta_6 t_0 + \delta_7 \Psi(\tilde{B}_0) + \delta_8 \sum_{i = 0}^{t - 1}D_i}{t}\right)^{t},
    \end{align*}
    where the inequality is due to the fact that $1 - e^{-x} \leq x$ for any $x \in \mathbb{R}$ and $\delta_6, \delta_7, \delta_8$ are defined in \eqref{def_delta_5678}. Hence, we prove that for any $t > t_0$,
    \begin{equation}\label{proof_theorem_superlinear_1}
    \frac{f(x_{t}) - f(x_*)}{f(x_{0}) - f(x_*)} \leq \left(1 - e^{-\frac{\delta_6 t_0 + \delta_7 \Psi(\tilde{B}_0) + \delta_8 \sum_{i = 0}^{t - 1}D_i}{t}}\right)^{t} \leq \left(\frac{\delta_6 t_0 + \delta_7 \Psi(\tilde{B}_0) + \delta_8 \sum_{i = 0}^{t - 1}D_i}{t}\right)^{t}.
    \end{equation}
    From \eqref{sum_of_C_t} in Theorem~\ref{thm:second_linear_rate}, we have that
    \begin{equation}
        \sum_{i = 0}^{t - 1}D_i \leq D_0\left(\Psi(\bar{B}_{0}) + \frac{3(1 + D_0)^2}{\alpha(1 - \beta)}\right).
    \end{equation}
    Therefore, combing the above inequality with \eqref{proof_theorem_superlinear_1}, we prove that
    \begin{equation*}
    \begin{split}
        & \frac{f(x_{t}) - f(x_*)}{f(x_{0}) - f(x_*)} \leq \left(\frac{\delta_6 t_0 + \delta_7 \Psi(\tilde{B}_0) + \delta_8 \sum_{i = 0}^{t - 1}D_i}{t}\right)^{t} \\
        & \leq \left(\frac{\delta_6 t_0 + \delta_7 \Psi(\tilde{B}_0) + \delta_8 D_0\left(\Psi(\bar{B}_{0}) + \frac{3(1 + D_0)^2}{\alpha(1 - \beta)}\right)}{t}\right)^{t}.
    \end{split}
    \end{equation*}

\section{Proof of Iteration Complexity}\label{proof_of_complexity}

We treat the line search parameters $\alpha$ and $\beta$ as absolute constants. The first linear rate from Theorem~\ref{thm:first_linear_rate} leads to the global complexity of 
\begin{equation}\label{complexity_1}
    \mathcal{O}( \Psi(\bar{B_0}) + (1 + D_0)^2 \log{\frac{1}{\epsilon}} )
\end{equation}
The second linear rate from Theorem~\ref{thm:second_linear_rate} leads to the global complexity of 
\begin{equation}\label{complexity_2}
    \mathcal{O}( \Psi(\tilde{B_0}) + (\Psi(\bar{B_0}) + (1 + D_0)^2)D_0 + \log{\frac{1}{\epsilon}} )
\end{equation}
where the first term is the number of iterations required to reach the linear rate in \eqref{eq:second_linear_rate}. For the analysis of the superlinear convergence rate, we denote that 
\begin{equation*}
    \Omega = \Psi(\tilde{B_0}) + (\Psi(\bar{B_0}) + (1 + D_0)^2)D_0
\end{equation*}
From Theorem~\ref{thm:superlinear_rate_2}, we have that
\begin{align*}
    \frac{f(x_t) - f(x_*)}{f(x_0) - f(x_*)} \leq (\frac{\Omega}{t})^{t}
\end{align*}
Let $T_*$ be the number such that the inequality $(\frac{\Omega}{t})^{t} \leq \epsilon$ above becomes equality. we have
\begin{align*}
    \log{\frac{1}{\epsilon}} = T_* \log{\frac{T_*}{\Omega}} \leq T_*(\frac{T_*}{\Omega} - 1),
\end{align*}
which leads to
\begin{align*}
    T_* \geq \frac{\Omega + \sqrt{\Omega^2 + 4\Omega\log{\frac{1}{\epsilon}}}}{2}.
\end{align*}
Hence, we have that
\begin{align*}
    \log{\frac{1}{\epsilon}} = T_* \log{\frac{T_*}{\Omega}} \geq T_*\log{\frac{\Omega + \sqrt{\Omega^2 + 4\Omega\log{\frac{1}{\epsilon}}}}{2\Omega}} \geq T_*\log{\left(\frac{1}{2} + \sqrt{\frac{1}{4} + \frac{\log{\frac{1}{\epsilon}}}{\Omega}}\right)},
\end{align*}
which implies that
\begin{align*}
    T_* \leq \frac{\log{\frac{1}{\epsilon}}}{\log{\left(\frac{1}{2} + \sqrt{\frac{1}{4} + \frac{\log{\frac{1}{\epsilon}}}{\Omega}}\right)}}.
\end{align*}
Hence, to reach the accuracy of $\epsilon$, we need the number of iterations $t$ to be at least
\begin{equation}\label{complexity_3}
    \mathcal{O}(\frac{\log{\frac{1}{\epsilon}}}{\log{\left(\frac{1}{2} + \sqrt{\frac{1}{4} + \frac{1}{\Omega}\log{\frac{1}{\epsilon}}}\right)}}).
\end{equation}
Therefore, we prove the iteration complexity by choosing the minimal from \eqref{complexity_1}, \eqref{complexity_2}, and \eqref{complexity_3}. For the special case of $B_0 = a I$ for $a > 0$, just replace $\Psi(\bar{B_0})$ and $\Psi(\tilde{B_0})$ by $\Delta_1$ and $\Delta_2$ defined in \eqref{Delta_1}, \eqref{Delta_2}, respectively.

\section{Proof of Line Search Complexity}\label{proof_of_linesearch}

\begin{algorithm}%[t]
\caption{Log Bisection Algorithm for Weak Wolfe Conditions}\label{algo_wolfe} 
\begin{algorithmic}[1] 
{\REQUIRE Initial step size $\eta^{(0)} = 1$, $\eta^{(0)}_{min} = 0$, $\eta^{(0)}_{max} = +\infty$
\FOR {$i = 0, 1, 2, \ldots$}
    \IF{$f(x_t + \eta^{(i)} d_t) > f(x_t) + \alpha \eta^{(i)} \nabla{f}(x_t)^\top d_t$}\label{line:sufficient_decrease}
        \STATE 
        Set $\eta^{(i + 1)}_{max} = \eta^{(i)}$ and  $\eta^{(i + 1)}_{min} = \eta^{(i)}_{min}$
        \IF{$\eta^{(i)}_{min} = 0$}
            \STATE $\eta^{(i + 1)} = (\frac{1}{2})^{2^{i + 1} - 1}$
        \ELSE
            \STATE $\eta^{(i + 1)} = \sqrt{\eta^{(i + 1)}_{max}\eta^{(i + 1)}_{min}}$
        \ENDIF
    \ELSIF{$\nabla{f}(x_t + \eta^{(i)} d_t)^\top d_t < \beta \nabla{f}(x_t)^\top d_t$}\label{line:curvature}
        \STATE Set $\eta^{(i + 1)}_{max} = \eta^{(i)}_{max}$ and $\eta^{(i + 1)}_{min} = \eta^{(i)}$
        \IF{$\eta^{(i)}_{max} = +\infty$}
            \STATE $\eta^{(i + 1)} = 2^{^{2^{i + 1} - 1}}$
        \ELSE
            \STATE $\eta^{(i + 1)} = \sqrt{\eta^{(i + 1)}_{max}\eta^{(i + 1)}_{min}}$
        \ENDIF
    \ELSE
        \STATE Return $\eta^{(i)}$
    \ENDIF
\ENDFOR}
\end{algorithmic}
\end{algorithm}

\begin{proposition}\label{proposition_stepsize}
    Suppose that Assumption~\ref{ass_self_concodant} holds. Consider the BFGS method with inexact line search defined in \eqref{sufficient_decrease} and \eqref{curvature_condition} and we choose the step size $\eta_t$ according to Algorithm~\ref{algo_wolfe}. At iteration $t$, denote $\lambda_t$ as the number of loops in Algorithm~\ref{algo_wolfe} to terminate and return the $\eta_t$ satisfying the Wolfe conditions \eqref{sufficient_decrease} and \eqref{curvature_condition}. Then $\lambda_t$ is finite and upper bounded by
    \begin{equation}
    \begin{split}
        & \lambda_t \leq 2 + \log_{2}\Big(1 + \frac{(1 - \beta)(1 + 2D_t)}{\beta - \alpha}\Big) \\
        & \qquad + 2\log_{2}\Big(1 + \log_{2}\big(2(1 - \alpha)(1 + D_t)\big) + \max\{\log_{2}\rho_t, \log_{2}\frac{1}{\rho_t}\}\Big).
    \end{split}
    \end{equation}
\end{proposition}

\begin{proof}
   Please check Proposition K.2 in \cite{qiujiang2024quasinewton}. The only difference is $C_t$ is replaced by $D_t$ defined in \eqref{suboptimality}.
\end{proof}

    We can prove the line search complexity in Proposition~\ref{prop:line_search} using result from Proposition~\ref{proposition_stepsize}. We have that
    \begin{equation}\label{thm:line_search_proof_1}
    \begin{split}
        \Lambda_t & = \frac{1}{t}\sum_{i = 0}^{t - 1} \lambda_i  \leq 2 + \frac{1}{t}\sum_{i = 0}^{t - 1}\log_{2}\Big(1 + \frac{(1 - \beta)(1 + 2D_i)}{\beta - \alpha}\Big) \\
        & + \frac{2}{t}\sum_{i = 0}^{t - 1}\log_{2}\Big(1 + \log_{2}\big(2(1 - \alpha)(1 + D_i)\big) + \max\{\log_{2}\rho_i, \log_{2}\frac{1}{\rho_i}\}\Big).
    \end{split}
    \end{equation}
    Using Jensen's inequality, we have that
    \begin{equation}\label{thm:line_search_proof_2}
    \begin{split}
        \frac{1}{t}\sum_{i = 0}^{t - 1}\log_{2}\Big(1 + \frac{(1 - \beta)(1 + 2D_i)}{\beta - \alpha}\Big) \leq \log_{2}\Big(1 + \frac{1 - \beta}{\beta - \alpha} + \frac{2(1 - \beta)}{\beta - \alpha}\frac{\sum_{i = 0}^{t - 1}D_i}{t}\Big).
    \end{split}
    \end{equation}
    \begin{equation}\label{thm:line_search_proof_3}
    \begin{split}
        & \frac{1}{t}\sum_{i = 0}^{t - 1}\log_{2}\Big(1 + \log_{2}\big(2(1 - \alpha)(1 + D_i)\big) + \max\{\log_{2}\rho_i, \log_{2}\frac{1}{\rho_i}\}\Big) \\
        & \leq \log_{2}\Big(1 +\log_{2}2(1 - \alpha) + \frac{1}{t}\sum_{i = 0}^{t - 1}\log_{2}(1 + D_i) + \frac{1}{t}\sum_{i = 0}^{t - 1}\max\{\log_{2}\rho_i, \log_{2}\frac{1}{\rho_i}\}\Big) \\
        & \leq \log_{2}\Big(1 +\log_{2}2(1 - \alpha) + \log_{2}\big(1 + \frac{\sum_{i = 0}^{t - 1}D_i}{t}) + \frac{1}{t}\sum_{i = 0}^{t - 1}\max\{\log_{2}\rho_i, \log_{2}\frac{1}{\rho_i}\}\Big).
    \end{split}
    \end{equation}
    We also have that
    \begin{equation}\label{thm:line_search_proof_4}
    \begin{split}
        & \quad \frac{1}{t}\sum_{i = 0}^{t - 1}\max\{\log_{2}\rho_i, \log_{2}\frac{1}{\rho_i}\} = \frac{1}{t}\sum_{i = 0, \rho_i \geq 1}^{t - 1}\log_{2}\rho_i + \frac{1}{t}\sum_{i = 0, 0 \leq \rho_i < 1}^{t - 1}\log_{2}\frac{1}{\rho_i} \\
        & = \frac{1}{t}\sum_{i = 0, \rho_i \geq 2}^{t - 1}\log_{2}\rho_i + \frac{1}{t}\sum_{i = 0, 1 \leq \rho_i < 2}^{t - 1}\log_{2}\rho_i + \frac{1}{t}\sum_{i = 0, \frac{1}{2} < \rho_i < 1}^{t - 1}\log_{2}\frac{1}{\rho_i} + \frac{1}{t}\sum_{i = 0, \rho_i \leq \frac{1}{2}}^{t - 1}\log_{2}\frac{1}{\rho_i} \\
        & \leq 2 + \frac{1}{t}\sum_{i = 0, \rho_i \geq 2}^{t - 1}\log_{2}\rho_i + \frac{1}{t}\sum_{i = 0, \rho_i \leq \frac{1}{2}}^{t - 1}\log_{2}\frac{1}{\rho_i},
    \end{split}
    \end{equation}
    where the inequality is due to $\log_{2}\rho_i \leq 1$ for $\rho_i < 2$ and $\log_{2}\frac{1}{\rho_i} \leq 1$ for $\rho_i > \frac{1}{2}$. Using the definition of $\omega$ and (b) in Lemma~\ref{lemma_omega}, we obtain that
    \begin{equation}\label{thm:line_search_proof_5}
    \begin{split}
        & \frac{1}{t}\sum_{i = 0, \rho_i \geq 2}^{t - 1}\log_{2}\rho_i = \frac{\log_{2}e}{t}\sum_{i = 0, \rho_i \geq 2}^{t - 1}\log\rho_i = \frac{\log_{2}e}{t}\sum_{i = 0, \rho_i \geq 2}^{t - 1}(\rho_i - 1 - \omega(\rho_i - 1)) \\
        & \leq \frac{\log_{2}e}{t}\sum_{i = 0, \rho_i \geq 2}^{t - 1}(\frac{2\rho_i}{\rho_i - 1}\omega(\rho_i - 1) - \omega(\rho_i - 1)) = \frac{\log_{2}e}{t}\sum_{i = 0, \rho_i \geq 2}^{t - 1}\frac{\rho_i + 1}{\rho_i - 1}\omega(\rho_i - 1) \\
        & \leq \frac{3\log_{2}e}{t}\sum_{i = 0, \rho_i \geq 2}^{t - 1}\omega(\rho_i - 1).
    \end{split}
    \end{equation}
    Similarly, using (c) in Lemma~\ref{lemma_omega}, we obtain that
    \begin{equation}\label{thm:line_search_proof_6}
    \begin{split}
        & \frac{1}{t}\sum_{i = 0, \rho_i \leq \frac{1}{2}}^{t - 1}\log_{2}\frac{1}{\rho_i} = \frac{\log_{2}e}{t}\sum_{i = 0, \rho_i \leq \frac{1}{2}}^{t - 1}\log\frac{1}{\rho_i} = \frac{\log_{2}e}{t}\sum_{i = 0, \rho_i \leq \frac{1}{2}}^{t - 1}(\omega(\rho_i - 1) + 1 - \rho_i) \\
        & \leq \frac{\log_{2}e}{t}\sum_{i = 0, \rho_i \leq \frac{1}{2}}^{t - 1}(\omega(\rho_i - 1) + \frac{1+ \rho_i}{1 - \rho_i}\omega(\rho_i - 1)) = \frac{\log_{2}e}{t}\sum_{i = 0, \rho_i \leq \frac{1}{2}}^{t - 1}\frac{2}{1 - \rho_i}\omega(\rho_i - 1) \\
        & \leq \frac{4\log_{2}e}{t}\sum_{i = 0, \rho_i \leq \frac{1}{2}}^{t - 1}\omega(\rho_i - 1).
    \end{split}
    \end{equation}
    Combining \eqref{thm:line_search_proof_4}, \eqref{thm:line_search_proof_5} and \eqref{thm:line_search_proof_6}, we prove that
    \begin{equation}\label{thm:line_search_proof_7}
    \begin{split}
        & \frac{1}{t}\sum_{i = 0}^{t - 1}\max\{\log_{2}\rho_i, \log_{2}\frac{1}{\rho_i}\} \leq 2 + \frac{1}{t}\sum_{i = 0, \rho_i \geq 2}^{t - 1}\log_{2}\rho_i + \frac{1}{t}\sum_{i = 0, \rho_i \leq \frac{1}{2}}^{t - 1}\log_{2}\frac{1}{\rho_i} \\
        & \leq 2 + \frac{4\log_{2}e}{t}\sum_{i = 0}^{t - 1}\omega(\rho_i - 1) \leq 2 + \frac{6}{t}\Big(\Psi(\tilde{B}_{0}) + 2\sum_{i = 0}^{t - 1}D_i\Big) .
    \end{split}
    \end{equation}
    where we use the fact that $\omega(\rho_i - 1) \geq 0$ for any $i \geq 0$ and the last inequality is due to \eqref{eq:omega} in Proposition~\ref{lemma_BFGS_superlinear}. Leveraging \eqref{thm:line_search_proof_1}, \eqref{thm:line_search_proof_2}, \eqref{thm:line_search_proof_3} and \eqref{thm:line_search_proof_7}, we have that
    \begin{equation*}
    \begin{split}
        \Lambda_t & \leq 2 + \log_{2}\Big(1 + \frac{1 - \beta}{\beta - \alpha} + \frac{2(1 - \beta)}{\beta - \alpha}\frac{\sum_{i = 0}^{t - 1}D_i}{t}\Big) \\
        & + 2\log_{2}\Big(3 +\log_{2}2(1 - \alpha) + \log_{2}\big(1 + \frac{\sum_{i = 0}^{t - 1}D_i}{t}) + \frac{6}{t}\big(\Psi(\tilde{B}_{0}) + 2\sum_{i = 0}^{t - 1}D_i\big)\Big) \\
        & \leq 2 + \log_{2}\Big(1 + \frac{1 - \beta}{\beta - \alpha} + \frac{2(1 - \beta)}{\beta - \alpha}\frac{\sum_{i = 0}^{t - 1}D_i}{t}\Big) \\
        & \quad + 2\log_{2}\Big(\log_{2}16(1 - \alpha) + \log_{2}\big(1 + \frac{\sum_{i = 0}^{t - 1}D_i}{t}) + \frac{6\Psi(\tilde{B}_{0}) + 12\sum_{i = 0}^{t - 1}D_i}{t}\Big) \\
        & \leq 2 + \log_{2}\Big(1 + \frac{1 - \beta}{\beta - \alpha} + \frac{2(1 - \beta)}{\beta - \alpha}\frac{\sum_{i = 0}^{t - 1}D_i}{t}\Big) \\
        & \quad + 2\log_{2}\Big(\log_{2}16(1 - \alpha) + \log_{2}\big(1 + \frac{6\Psi(\tilde{B}_{0}) + 14\sum_{i = 0}^{t - 1}D_i}{t}\Big).
    \end{split}
    \end{equation*}
    Using \eqref{sum_of_C_t} from the proof of Theorem~\ref{thm:second_linear_rate}, i.e.,
    \begin{equation*}
        \sum_{i = 0}^{t - 1}D_i \leq D_0\left(\Psi(\bar{B}_{0}) + \frac{3(1 + D_0)^2}{\alpha(1 - \beta)}\right).
    \end{equation*}
    We prove the line search complexity as
    \begin{equation*}
        \Lambda_t = \mathcal{O}\left(\log(1 + \frac{\Gamma}{t}) + \log\log(1 + \frac{\Psi(\tilde{B}_{0}) + \Gamma}{t})\right)
    \end{equation*}
    where
    \begin{equation*}
        \Gamma = \mathcal{O}\left(D_0(\Psi(\bar{B}_{0}) + (1 + D_0)^2)\right)
    \end{equation*}
    For the special case of $B_0 = a I$ for $a > 0$, just replace $\Psi(\bar{B_0})$ and $\Psi(\tilde{B_0})$ by $\Delta_1$ and $\Delta_2$ defined in \eqref{Delta_1}, \eqref{Delta_2}, respectively.

\section{Proof of Strong Self-Concordance}\label{proof_of_concordance}

Consider the log-sum-exp function defined as $f(x) = \log{(\sum_{i = 1}^{n}\exp({c_i^\top x - b_i}))} + \frac{1}{2}\sum_{i = 1}^{n}(c_i^\top x)^2$, we have that $\nabla{f}(x) = \sum_{i = 1}^{n}\pi_i c_i + \sum_{i = 1}^{n}(c_i^\top x)c_i$ where $\pi_i = \frac{\exp({c_i^\top x - b_i}))}{\sum_{j = 1}^{n}\exp({c_j^\top x - b_j}))}$ and $\nabla^2{f}(x) = \sum_{i = 1}^{n}(\pi_i + 1) c_i c_i^\top - (\sum_{i = 1}^{n}\pi_i c_i)(\sum_{i = 1}^{n}\pi_i c_i)^\top$. From proof in \cite{logsumconvex}, this log-sum-exp function is strictly convex. Moreover, we also need to prove that this function is strongly self-concordant. Notice that, with respect to the operator $B = \sum_{i = 1}^{n}c_i c_i^\top$, this function f is strongly convex with parameter $1$ and its Hessian is Lipschitz continuous with parameter $2$ (check example 1 of \cite{uniformlyconvex}\footnote{N. Doikov and Y. Nesterov. Minimizing uniformly convex functions by cubic regularization of newton method. arXiv, 1905.02671, 2019}). Hence, using results from Example 4.1 in \cite{rodomanov2020greedy}, the log-sum-exp function is strongly self-concordant.

The proof of that the logistic regression function $f(x) = \frac{1}{N}\sum_{i = 1}^{N}\ln{(1 + e^{-y_i z_i^\top x})}$ without $l_2$ regularization is strongly self-concordant is almost the same. It has the similar structure with the log-sum-exp function $f(x) = log{(\sum_{i = 1}^{n}e^{c_i^\top x - b_i})} + \frac{1}{2}\sum_{i = 1}^{n}(c_i^\top x)^2$. Notice that in our BFGS method, we use the line search scheme such that we always have $f(x_t) \leq f(x_0)$ for any $t \geq 0$. Hence, the iterations generated by BFGS method with weak-Wolfe line search conditions always stay in the bounded set $\{x|f(x) \leq f(x_0)\}$ where $x_0$ is the initial point. In this bounded set, the logistic regression function is strongly convex and its Hessian is smooth with respect to the operator matrix $B = \sum_{i = 1}^{n}z_iz_i^\top$. According to the Example 4.1 from that greedy quasi-Newton paper, if a function is strongly convex and its Hessian is smooth with respect to some matrix $B$, then the function is strongly self-concordant. Hence, the logistic regression function is strongly self-concordant. Similarly, for the hard cubic function, we can show that it is strongly convex and its Hessian is Lipschitz smooth with respect to some operator matrix $B$ and therefore the hard cubic function is also strongly self-concordant.

\vspace{-3mm}

\section{Additional Numerical Experiments}\label{add_experiments}

\vspace{-3mm}

Additional numerical experiments on the hard cubic function and the logistic regression for different dimensions are presented in figures~\ref{fig:5} and \ref{fig:6}. The empirical results of the performance of different optimization methods for the hard cubic function with respect to the number of gradient evaluations and the time in seconds are in Figures~\ref{fig:7} and \ref{fig:8}. Additional numerical results of the values of the step sizes of BFGS method are in Figure~\ref{fig:9}. Additional results of the performance of different optimization methods with transformation matrix are in Figure~\ref{fig:10}. The convergence performance of BFGS method is similar to the empirical results from Figures~\ref{fig:1}, \ref{fig:2}, \ref{fig:3}, and \ref{fig:4} in section~\ref{sec:experiments}. 

\begin{figure}[t!]
    \centering
    \subfigure[$d = 50$.]{\includegraphics[width=0.49\linewidth]{Figures/Hard_Cubic/d5e1.pdf}}
    \subfigure[$d = 300$.]{\includegraphics[width=0.49\linewidth]{Figures/Hard_Cubic/d3e2.pdf}}
    \subfigure[$d = 600$.]{\includegraphics[width=0.49\linewidth]{Figures/Hard_Cubic/d6e2.pdf}}
    \subfigure[$d = 2000$.]{\includegraphics[width=0.49\linewidth]{Figures/Hard_Cubic/d2e3.pdf}}
    \caption{Convergence rates of BFGS with different $B_0$, gradient descent and accelerated gradient descent for solving the hard cubic function with different dimensions.}
    \label{fig:5}
\end{figure}

\begin{figure}[t!]
    \centering
    \subfigure[$d = 50$.]{\includegraphics[width=0.49\linewidth]{Figures/Logistic/d5e1.pdf}}
    \subfigure[$d = 300$.]{\includegraphics[width=0.49\linewidth]{Figures/Logistic/d3e2.pdf}}
    \subfigure[$d = 600$.]{\includegraphics[width=0.49\linewidth]{Figures/Logistic/d6e2.pdf}}
    \subfigure[$d = 2000$.]{\includegraphics[width=0.49\linewidth]{Figures/Logistic/d2e3.pdf}}
    \caption{Convergence rates of BFGS with different $B_0$, gradient descent and accelerated gradient descent for solving the logistic regression function with different dimensions.}
    \label{fig:6}
\end{figure}

\begin{figure}[t!]
    \centering
    \subfigure[$d = 100$.]{\includegraphics[width=0.49\linewidth]{Figures/GradEval/d100.pdf}}
    \subfigure[$d = 200$.]{\includegraphics[width=0.49\linewidth]{Figures/GradEval/d200.pdf}}
    \subfigure[$d = 500$.]{\includegraphics[width=0.49\linewidth]{Figures/GradEval/d500.pdf}}
    \subfigure[$d = 1000$.]{\includegraphics[width=0.49\linewidth]{Figures/GradEval/d1000.pdf}}
    \caption{Convergence rates of BFGS with different $B_0$, gradient descent and accelerated gradient descent for solving the hard cubic function with respect to the number of gradient evaluations.}
    \label{fig:7}
\end{figure}

\begin{figure}[t!]
    \centering
    \subfigure[$d = 100$.]{\includegraphics[width=0.49\linewidth]{Figures/Time/d100.pdf}}
    \subfigure[$d = 200$.]{\includegraphics[width=0.49\linewidth]{Figures/Time/d200.pdf}}
    \subfigure[$d = 500$.]{\includegraphics[width=0.49\linewidth]{Figures/Time/d500.pdf}}
    \subfigure[$d = 1000$.]{\includegraphics[width=0.49\linewidth]{Figures/Time/d1000.pdf}}
    \caption{Convergence rates of BFGS with different $B_0$, gradient descent and accelerated gradient descent for solving the hard cubic function with respect to the time in seconds.}
    \label{fig:8}
\end{figure}

\begin{figure}[t!]
    \centering
    \subfigure[$d = 100$.]{\includegraphics[width=0.45\linewidth]{Figures/StepSize/d200.pdf}}
    \subfigure[$d = 1000$.]{\includegraphics[width=0.45\linewidth]{Figures/StepSize/d500.pdf}}
    \caption{Step size of BFGS with different $B_0$ using inexact line search for solving the hard cubic function with different dimensions.}
    \label{fig:9}
\end{figure}

\begin{figure}[t!]
    \centering
    \subfigure[$d = 100$.]{\includegraphics[width=0.45\linewidth]{Figures/Transform/d200.pdf}}
    \subfigure[$d = 1000$.]{\includegraphics[width=0.45\linewidth]{Figures/Transform/d500.pdf}}
    \caption{Convergence rates of BFGS with different $B_0$, gradient descent and accelerated gradient descent for solving the hard cubic function with transformation matrix $A$.}
    \label{fig:10}
\end{figure}

\clearpage

\printbibliography

\end{document}
